% ===========================================================================
%  TEMPLATE — RESUMO EXPANDIDO
%  Seminário de Iniciação Científica (SIC 2026)
%  IFNMG — Campus Januária
% ===========================================================================

\documentclass[
  12pt,
  a4paper,
  oneside,
  article,
  brazil
]{abntex2}

% ---------------------------------------------------------------------------
% PACOTES
% ---------------------------------------------------------------------------
\usepackage[T1]{fontenc}
\usepackage[utf8]{inputenc}
\usepackage{lmodern}
\usepackage{microtype}
\usepackage{graphicx}
\usepackage{float}
\usepackage{booktabs}
\usepackage{caption}
\usepackage{subcaption}
\usepackage{amsmath,amssymb}
\usepackage{indentfirst}
\usepackage{geometry}
\usepackage{fancyhdr}
\usepackage{lastpage}
\usepackage{xcolor}
\usepackage{enumitem}
\usepackage{url}
\usepackage{multirow}
\usepackage{setspace}

% ---------------------------------------------------------------------------
% CITAÇÕES E REFERÊNCIAS ABNT
% ---------------------------------------------------------------------------
\usepackage[
  backend=biber,
  style=abnt,
  giveninits=true
]{biblatex}

\addbibresource{referencias.bib}

% ---------------------------------------------------------------------------
% MARGENS ABNT NBR 6022
% ---------------------------------------------------------------------------
\geometry{top=2cm, bottom=2cm, left=3cm, right=2cm}

\setlength{\parindent}{1.25cm}
\setlength{\parskip}{0pt}

% ---------------------------------------------------------------------------
% CORES IFNMG
% ---------------------------------------------------------------------------
\definecolor{ifnmgverde}{RGB}{0,104,55}
\definecolor{ifnmgcinza}{RGB}{88,88,90}

% ---------------------------------------------------------------------------
% HYPERLINKS
% ---------------------------------------------------------------------------
\usepackage{hyperref}
\hypersetup{
  colorlinks=true,
  linkcolor=black,
  citecolor=black,
  urlcolor=blue,
  pdftitle={Geração de Layouts de Fases 3D Assistida por Modelos de Linguagem de Grande Porte},
  pdfauthor={Gustavo Rafael Nunes Durães}
}

% ---------------------------------------------------------------------------
% CABECALHO E RODAPE
% ---------------------------------------------------------------------------
\pagestyle{fancy}
\fancyhf{}
\renewcommand{\headrulewidth}{0.4pt}
\fancyhead[L]{\small\textit{Projeto de TCC}}
\fancyhead[R]{\small\textit{\nouppercase{\leftmark}}}
\fancyfoot[C]{\small Página \thepage\ de \pageref{LastPage}}

\fancypagestyle{primeirapagina}{
  \fancyhf{}
  \renewcommand{\headrulewidth}{0pt}
  \fancyfoot[C]{\small Página \thepage\ de \pageref{LastPage}}
}

\aliaspagestyle{chapter}{fancy}
\aliaspagestyle{abntchapfirst}{fancy}
\aliaspagestyle{cleared}{fancy}
\aliaspagestyle{part}{fancy}

% ---------------------------------------------------------------------------
% SECOES — estilo ABNT
% ---------------------------------------------------------------------------
\renewcommand{\ABNTEXchapterfont}{\bfseries}
\renewcommand{\ABNTEXchapterfontsize}{\normalsize}
\renewcommand{\ABNTEXsectionfont}{\bfseries}
\renewcommand{\ABNTEXsectionfontsize}{\normalsize}
\renewcommand{\ABNTEXsubsectionfont}{\bfseries\itshape}
\renewcommand{\ABNTEXsubsectionfontsize}{\normalsize}
\renewcommand{\ABNTEXsubsubsectionfont}{\itshape}
\renewcommand{\ABNTEXsubsubsectionfontsize}{\normalsize}

\captionsetup{
  font=small,
  labelfont=bf,
  labelsep=endash,
  justification=raggedright,
  singlelinecheck=false
}

% ===========================================================================
%  INÍCIO DO DOCUMENTO
% ===========================================================================
\begin{document}
\textual
\aliaspagestyle{plain}{fancy}
\aliaspagestyle{chapter}{fancy}
\aliaspagestyle{abntchapfirst}{fancy}
\pagestyle{fancy}
\thispagestyle{primeirapagina}
\OnehalfSpacing

\vspace{0.15cm}

\noindent\rule{\linewidth}{1pt}
\vspace{0.05cm}

% ---------------------------------------------------------------------------
% TITULO
% ---------------------------------------------------------------------------
\begin{center}
  {\large\bfseries
    GERAÇÃO DE LAYOUTS DE FASES 3D ASSISTIDA POR MODELOS DE LINGUAGEM DE GRANDE PORTE
  }
\end{center}

\vspace{0.1cm}

% ---------------------------------------------------------------------------
% AUTORES E AFILIAÇÕES
% ---------------------------------------------------------------------------
\begin{center}
  \begin{SingleSpace}
    {\normalsize
      Gustavo Rafael Nunes DURÃES\textsuperscript{1,*},\;
      Wagner Ferreira de BARROS\textsuperscript{2}
    }\\[3pt]
    {\small
      \textsuperscript{1}Discente do Curso de Bach. em Ciência da Computação, IFNMG --- Campus Montes Claros.
      \texttt{grnd@aluno.ifnmg.edu.br}%
      \;$|$\;\textit{*autor correspondente}\\
      \textsuperscript{2}Docente do IFNMG --- Campus Montes Claros.
      \texttt{wagner.barros@ifnmg.edu.br}
    }
  \end{SingleSpace}
\end{center}

\vspace{0.2cm}
\noindent\rule{\linewidth}{0.4pt}
\vspace{0.05cm}

% ---------------------------------------------------------------------------
% RESUMO
% ---------------------------------------------------------------------------
{\centering\textbf{RESUMO}\par}
\vspace{0.1cm}
\begin{SingleSpace}
\noindent
A criação manual de fases para jogos digitais é um processo exaustivo que exige a definição minuciosa de caminhos, salas, regras de progressão e a validação geométrica de espaços tridimensionais (3D). Embora os Modelos de Linguagem de Grande Porte (LLMs) apresentem excelente capacidade de interpretação semântica e planejamento lógico a partir de linguagem natural, eles demonstram baixa confiabilidade no raciocínio espacial e geométrico exigido para a modelagem 3D. O presente trabalho propõe uma arquitetura que separa a geração procedimental de fases em duas responsabilidades distintas: um modelo atuando como "arquiteto", encarregado de traduzir intenções textuais do usuário em um grafo lógico estruturado (JSON), e um algoritmo determinístico atuando como "engenheiro", desenvolvido na \textit{engine} Godot, encarregado de validar a acessibilidade das salas e instanciar os componentes 3D no espaço sem colisões indevidas. O objetivo da pesquisa é avaliar o impacto da etapa de validação determinística na coerência das fases geradas, comparando a abordagem assistida estruturalmente com abordagens de geração unificada. Como resultados esperados, almeja-se a criação de um \textit{addon} funcional para Godot que demonstre que a divisão de tarefas (lógica vs. geometria) supera gargalos de inconsistência arquitetural presentes nas abordagens baseadas puramente em LLMs, garantindo a geração de fases solváveis, logicamente conectadas e geometricamente precisas.

\vspace{0.3cm}
\noindent\textbf{Palavras-chave:} Geração Procedural de Conteúdo. Modelos de Linguagem de Grande Porte. Godot Engine. Level Design. JSON.
\end{SingleSpace}

\vspace{0.1cm}
\noindent\rule{\linewidth}{0.4pt}
\vspace{0.1cm}

% ---------------------------------------------------------------------------
% ABSTRACT
% ---------------------------------------------------------------------------
{\centering\textbf{ABSTRACT}\par}
\vspace{0.1cm}
\begin{SingleSpace}
\noindent
The manual creation of levels for digital games is an exhausting process that requires the meticulous definition of paths, rooms, progression rules, and the geometric validation of three-dimensional (3D) spaces. Although Large Language Models (LLMs) exhibit excellent semantic interpretation and logical planning capabilities from natural language, they show low reliability in the spatial and geometric reasoning required for 3D modeling. This ongoing work proposes an architecture that separates procedural level generation into two distinct responsibilities: a model acting as an "architect", tasked with translating the user's textual intentions into a structured logical graph (JSON), and a deterministic algorithm acting as an "engineer", developed in the Godot engine, tasked with validating room accessibility and instantiating 3D components in space without improper collisions. The aim of this research is to evaluate the impact of the deterministic validation step on the coherence of the generated levels, comparing the structurally assisted approach with unified generation approaches. As expected results, we aim to create a functional addon for Godot that demonstrates that dividing tasks (logic vs. geometry) overcomes architectural inconsistency bottlenecks present in purely LLM-based approaches, ensuring the generation of solvable, logically connected, and geometrically accurate levels.

\vspace{0.3cm}
\noindent\textbf{Keywords:} Procedural Content Generation. Large Language Models. Godot Engine. Level Design. JSON.
\end{SingleSpace}

\vspace{0.3cm}
\noindent\rule{\linewidth}{1pt}
\vspace{0.3cm}

% ===========================================================================
%  CORPO DO TEXTO
% ===========================================================================

% ---------------------------------------------------------------------------
\section{Introdução}
% ---------------------------------------------------------------------------

A Geração Procedural de Conteúdo (PCG - \textit{Procedural Content Generation}) é amplamente utilizada na indústria de jogos para criar mundos expansivos e variados de forma algorítmica, mitigando o esforço manual do \textit{Level Design} \cite{summerville2018procedural}. Tradicionalmente, métodos determinísticos garantem que o conteúdo gerado (como masmorras ou mapas) atenda a critérios estritos de jogabilidade, como conectividade e solvabilidade. 

Com o advento dos Modelos de Linguagem de Grande Porte (LLMs), surgiu um novo paradigma denominado PCG via Machine Learning (PCGML). Trabalhos recentes demonstraram a viabilidade de usar LLMs para gerar níveis estruturados em representações textuais bidimensionais, como na matriz de posições de blocos para \textit{Sokoban} \cite{todd2023level}. Ademais, abordagens integradas com motores de jogos demonstram o alto potencial dessas redes na orquestração semântica e na assistência criativa contínua \cite{tang2025unrealllm}.

Apesar da alta proficiência em raciocínio lógico e interpretação de intenções em linguagem natural, os LLMs lutam significativamente na coordenação espacial e raciocínio puramente geométrico em espaços tridimensionais (3D). Ao tentar gerar o posicionamento 3D absoluto de estruturas arquitetônicas complexas através de uma rede neural estocástica, frequentemente surgem erros críticos como intersecção de malhas, corredores sem saída e chaves geradas em salas inacessíveis \cite{xu2025constraint, sun2024roomplanner}.


O objetivo principal deste estudo é desenvolver e avaliar um sistema no qual o LLM processa o \textit{prompt} do usuário e gera exclusivamente a macroestrutura lógica da fase, representada através de um grafo em formato JSON, propondo uma abordagem arquitetural de responsabilidades segregadas. A premissa fundamental, que atua como problema motivador da pesquisa, é a de que o LLM pensa e planeja a fase conceitualmente, enquanto um algoritmo determinístico a verifica e a constrói fisicamente. A instânciação geométrica e a validação topológica, por sua vez, são delegadas a um módulo procedimental desenvolvido na \textit{engine} Godot. 

Espera-se que essa segregação reduza o escopo de falha do modelo probabilístico, mantendo a expressividade oriunda da linguagem natural e assegurando a criação de mapas totalmente jogáveis.
% ---------------------------------------------------------------------------
\section{Material e Métodos}
% ---------------------------------------------------------------------------

Para operacionalizar a geração assistida, a metodologia adota um fluxo de trabalho em \textit{pipeline}, dividido em três etapas sequenciais: (1) Processamento da Linguagem, (2) Validação Determinística e (3) Construção e Testes, todas integradas através de um \textit{addon} em desenvolvimento para o Godot Engine.

\textbf{1. Processamento da Linguagem:} Inicialmente, um LLM recebe a requisição estrutural do desenvolvedor/jogador em linguagem natural (e.g., `crie uma masmorra com uma sala inicial, uma bifurcação onde um lado contém um tesouro e o outro lado demanda uma chave localizada no tesouro para acessar a sala do chefe''). O modelo é instruído via \textit{Prompt Engineering} a omitir cálculos espaciais (X, Y, Z) e exportar as conexões semânticas na forma de um grafo serializado em JSON. Essa matriz lógica reflete relações hierárquicas, ramificações e pré-requisitos mecânicos (portas trancadas vs. chaves).

\textbf{2. Validação Determinística:} O Godot intercepta o arquivo JSON e processa o grafo gerado através de um sistema baseado em regras clássicas. O algoritmo averígua condicionais estritas necessárias para uma fase funcional. As regras avaliam premissas topológicas, como: a existência imperativa de uma única sala inicial; a alcançabilidade da sala do chefe; a não alocação de chaves atrás da porta que elas abrem, tornando a sala inacessível; e a garantia de que as salas bônus (como tesouros) não estrangulem a rota principal de progresso. Caso o grafo apresente anomalias lógicas impeditivas, o sistema está sendo modelado para recusar o mapa e retroalimentar o LLM com as diretrizes do erro lógico.

\textbf{3. Construção do Mapa e Testes Espaciais:} Validado o grafo, o gerador espacial entra em vigor. Ele realiza a alocação procedural \textit{tile-based} ou via \textit{prefabs} de salas de combate, corredores e transições \cite{xu2025constraint}. O algoritmo determinístico calcula dinamicamente as posições espaciais com base no tamanho individual de cada malha tridimensional, impedindo sobreposições visuais e inconsistência física (\textit{clipping}). Subsequentemente, a \textit{engine} aplica ferramentas nativas, como nós de \textit{NavigationRegion3D} e projeção de \textit{RayCasts}, para atestar que o caminho principal possui uma malha de navegação (NavMesh) intacta, conectando fisicamente a entrada à saída.
% ---------------------------------------------------------------------------
\section{Metodologia Experimental e Resultados Esperados}
% ---------------------------------------------------------------------------

Tendo em vista que o trabalho encontra-se na fase de desenvolvimento e implementação arquitetural, esta seção detalha a formatação dos experimentos programados para validar a premissa da pesquisa, bem como as inferências qualitativas e quantitativas previstas.

\subsection{Delineamento dos Experimentos}

Para atestar o impacto da validação determinística, os experimentos serão divididos em dois tratamentos comparativos gerando amostras de fases de dimensões variadas (pequena, média e grande escala, traduzidas em número de cômodos):

\begin{itemize}[leftmargin=*]
  \item \textbf{Grupo Controle (Geração Unificada):} O LLM será solicitado a gerar não somente as conexões lógicas, mas também as coordenadas 3D de instanciação para cada \textit{asset} e cômodo.
  \item \textbf{Grupo Experimental (Geração Segregada):} A metodologia proposta na qual o LLM produz estritamente a árvore/grafo JSON da progressão, com o algoritmo determinístico encarregado de resolver e instanciar os obstáculos e coordenadas espaciais.
\end{itemize}

As amostras de fases resultantes de ambos os tratamentos serão expostas à verificações por meio de rotinas automatizadas no Godot que auditarão: (a) o número de colisões estruturais inválidas (sobreposição geométrica); (b) o tempo (milissegundos) para computar a malha navegável do início ao fim do nível; (c) a presença de bloqueios mecânicos (ex: chaves retidas em salas incomunicáveis).

\subsection{Resultados Esperados}

De acordo com o estado da arte e observações empíricas com modelos de linguagem aplicados a problemas de \textit{Facility Layout Optimization} \cite{xu2025constraint}, preveem-se os seguintes resultados primários:

\begin{itemize}[leftmargin=*]
    \item \textbf{Eliminação de Anomalias Geométricas:} Espera-se que a abordagem segregada demonstre uma taxa de falhas espaciais tendendo a zero, visto que a instaciação passará a obedecer a matrizes de colisão rigorosas que não são compreendidas eficientemente pela rede estocástica do LLM.
    \item \textbf{Resiliência e Coerência de \textit{Design}:} Pela validação por grafo anterior à instanciação espacial, a metodologia deve impedir inteiramente a geração de bloqueios intransponíveis provocados por falhas semânticas em rotas ramificadas, permitindo níveis longos e coesos que combinem salas de recompensas com salas de desafio.
    \item \textbf{Automatização em \textit{Addon}:} Como substrato prático da pesquisa, espera-se finalizar um \textit{plugin} integrado à interface nativa do Godot. O pacote proverá um campo de \textit{prompt} diretamente no inspetor do editor, convertendo linguagem natural em cenários físicos editáveis, otimizando drasticamente o \textit{workflow} do artista técnico ou \textit{Level Designer}.
\end{itemize}

% ---------------------------------------------------------------------------
\section{Considerações Finais}
% ---------------------------------------------------------------------------

O presente trabalho propõe atacar a deficiência inerente dos LLMs no raciocínio espacial contínuo transferindo os cálculos matemáticos para as \textit{game engines}. A arquitetura em desenvolvimento pauta-se no conceito de cooperação entre um componente estocástico-criativo (o modelo, como arquiteto) e um componente determinístico-rígido (o algoritmo, como engenheiro). 

À medida em que a prova de conceito (\textit{Addon} para Godot) avançar em testes empíricos de malhas navegáveis e consistência estrutural, será possível aferir a eventual redução da taxa de mapas inacessíveis, validando o arranjo teórico. Por fim, caso o método demonstre melhoras na solução problema investigado, projeta-se que a simbiose entre as descrições de \textit{layout} via LLM e algoritmos de construtivismo direcional fomentará metodologias de prototipagem hiper-ágeis para o mercado de jogos digitais, permitindo aos desenvolvedores concentrarem-se primordialmente na parte artística, enquanto delegam o trabalho estrutural bruto a modelos inteligentes aliados à lógica formal.

%% ---------------------------------------------------------------------------
\section{Fundamentação Teórica}

% ---------------------------------------------------------------------------

Este capítulo apresenta a base conceitual necessária para o desenvolvimento da arquitetura metodológica do projeto. O escopo abrange desde os fundamentos da Geração Procedural de Conteúdo (PCG), passando pelas capacidades e limitações espaciais dos Modelos de Linguagem de Grande Porte (LLMs), até culminar no estado da arte das arquiteturas híbridas neuro-simbólicas. Tais conceitos são indispensáveis para justificar a delegação de responsabilidades entre o modelo generativo -- que atua como arquiteto -- e o motor de jogo determinístico, que atua como engenheiro validador das restrições espaciais tridimensionais.

\subsection{Geração Procedural de Conteúdo (PCG) e Machine Learning (PCGML)}

A Geração Procedural de Conteúdo (PCG - \textit{Procedural Content Generation}) é uma subárea fundamental na pesquisa de Inteligência Artificial aplicada a jogos. Historicamente, PCG refere-se à criação algorítmica de dados do jogo (como níveis, narrativas, texturas e regras) de maneira autônoma ou semi-autônoma, muitas vezes combinada à aleatoriedade para garantir variabilidade e rejogabilidade. Tradicionalmente, as abordagens de PCG podiam ser categorizadas como métodos construtivos (como autômatos celulares e gramáticas de forma) ou métodos baseados em busca (algoritmos genéticos e otimização evolucionária) \cite{summerville2018procedural}. 
No que tange aos métodos clássicos, como destacado na literatura de geração procedural \cite{vieira2024metodologia}, a principal vantagem reside na garantia de resolução: uma vez que as regras determinísticas de particionamento de espaço (como as encontradas na estruturação modular baseada em \textit{tiles}) são definidas, a saída estrutural é previsível. Contudo, essa abordagem de PCG pura esbarra na falta de semântica criativa; os geradores construtivos produzem artefatos funcionalmente corretos, mas frequentemente carecem de uma narrativa ou intenção arquitetônica clara.
Para contornar a rigidez dos métodos construtivos tradicionais, emergiu a vertente \textit{Procedural Content Generation via Machine Learning} (PCGML). Segundo o trabalho fundamental de \cite{summerville2018procedural}, o PCGML propõe treinar modelos a partir de uma base de dados estruturada de conteúdo de jogos existentes. Ao contrário dos algoritmos baseados em busca, os quais geram níveis testando-os através de uma função de aptidão (\textit{fitness function}), o modelo treinado gera o conteúdo diretamente emulando os padrões latentes no corpus de dados originais. Inicialmente limitados a modelos estatísticos mais modestos (como cadeias de Markov bidimensionais ou LSTMs), a ascensão dos grandes modelos fundacionais abriu portas para uma formulação textual mais declarativa para a criação dos mundos. É sob essa ótica evolutiva do PCGML que o uso atual de Modelos de Linguagem de Grande Porte (LLMs) se apoia \cite{Gallotta_2024}.

\subsection{O Problema da Ancoragem Espacial (\textit{Spatial Grounding}) em LLMs}
Os Modelos de Linguagem de Grande Porte (LLMs) revolucionaram diversas esferas da automação e da IA generativa por meio de redes neurais baseadas na arquitetura \textit{Transformer}. Eles demonstram proficiência notável no raciocínio estruturado, planejamento narrativo e na manutenção de coerência temática. Sua capacidade inata de entender dependências lógicas os torna ideais para desenhar fluxos de progressão, gerando rapidamente esquemas semânticos formatados, por exemplo, como estruturas JSON. No entanto, o paradigma do treinamento em PLN (Processamento de Linguagem Natural), que é estritamente baseado no sequenciamento unidimensional (1D) de \textit{tokens}, introduz sérias desvantagens quando aplicado em contextos espaciais tridimensionais (3D).
Pesquisas recentes de metodologias aplicadas apontam a carência de \textit{spatial grounding} (ancoragem física e espacial) como a principal barreira técnica enfrentada pelos modelos generativos atuais em domínios arquitetônicos de jogos. Estudo exploratório sobre o \textit{workflow} do level designer perante a IA indica que LLMs são excepcionalmente propensos a cometer "alucinações geométricas" em arranjos 3D \cite{re2026integration}. Sem estar amparado a uma verificação geométrica contínua, o LLM falha ao dimensionar distâncias, prever áreas de secção transacional e posicionar elementos num \textit{grid} tridimensional complexo.
Tais achados são sustentados de maneira recorrente pela literatura do estado da arte. Por exemplo, trabalhos como o \textit{Code-as-Room} \cite{yang2026codeasroomgenerating3drooms} e o estudo de \cite{feng2026repurposing3dgenerativemodel} acerca da geração espacial baseada em layout autorregressivo demonstram evidências empíricas de que a geração puramente pautada na inferência probabilística leva inevitavelmente ao \textit{overlap} de \textit{bounding boxes} (intersecção e sobreposição de malhas) e ao deslocamento de portas ou vias cruciais para áreas inacessíveis. \cite{sun2024roomplanner} chegam a conclusões semelhantes em seu estudo sobre planejadores explícitos: modelos de texto puro não compreendem instintivamente as propriedades da geometria euclidiana, frequentemente posicionando mobiliário flutuando e gerando cômodos cujas portas desembocam no vácuo ou se conflitam. Portanto, essas limitações inviabilizam o uso "direto" de um LLM como motor que edita livremente um espaço virtual.

\subsection{Modelagem de Fases via Grafos de Progressão}

Cientes da deficiência dos LLMs no espaço e sua eficácia na lógica semântica pura, autores vêm recorrendo à teoria dos grafos. A estruturação de fases em \textit{Level Design} pode ser separada em dois paradigmas: o da "topologia" e o da "topografia" (espaço contínuo e dimensão física).
Topologicamente, uma fase inteira de um jogo digital (como uma masmorra ou um complexo labiríntico) pode ser abstratamente resumida a um grafo relacional ou grafo de progressão lógico \cite{chen2025narrative}. Nesse grafo, os vértices (nós) representam as salas, pátios ou instâncias temáticas e os nós direcionais ou arestas delineiam a conectividade (portas, bloqueios lógicos ou teletransportes). 
Esta representação abstrai completamente as coordenadas flutuantes ou as restrições tridimensionais, aproximando-se perfeitamente da expertise linguística de um LLM \cite{Gallotta_2024}. Delegar a criação de grafos relacionais (como uma cadeia estruturada de chaves JSON determinando o nó A conecta-se ao nó B condicionalmente) ao invés do mapeamento métrico do espaço resolve a incapacidade da Inteligência Artificial em operar sobre a geometria. Este modelo de divisão em grafos viabiliza que a ferramenta generativa se limite à elaboração do ritmo do percurso e ao arranjo sistêmico das recompensas e ameaças, blindando a arquitetura do projeto de alucinações matemáticas espaciais.

\subsection{Arquiteturas Híbridas e Lógica Determinística}
A síntese entre o poder generativo de um LLM com a confiabilidade rígida do processamento clássico resultou na concepção de sistemas neuro-simbólicos ou \textit{pipelines} de arquiteturas híbridas. Na literatura atual envolvendo PCG avançado, esta dicotomia costuma ser referida pelo paradigma Arquiteto-Engenheiro ou "Dual-Agent", sendo amplamente considerada o estado da arte para sínteses em 3D.
Frameworks baseados neste paradigma delegam ao LLM estritamente a tarefa de ideação ou macro-planejamento semântico (\textit{Arquiteto}), enquanto injetam essas descrições simbólicas para um validador algorítmico tradicional (\textit{Engenheiro}). A execução e instanciamento destas diretrizes se valem de programação estrita baseada em restrições lógicas (\textit{Constraints}) \cite{xu2025constraint}. Trabalhos como \textit{SceneCraft} \cite{hu2024scenecraft}, voltados à síntese programática no Blender, bem como implementações direcionadas à interpretação de cenas em jogos virtuais \cite{xu2025databasedrivenframework3dlevel}, utilizam explicitamente LLMs apenas para a estrutura narrativa base, deixando a manipulação do mundo e a aplicação de coordenadas matemáticas sob estrito controle algorítmico.
Em contextos equivalentes, \cite{her2025zeroshot3dmapgeneration} exploram a eficácia da abordagem de agência dupla (\textit{Dual-Agent Architecture}) para mapas 3D \textit{Zero-shot}, onde um agente projeta os requisitos criativos da cena e um outro atua como validador implacável de regras operacionais em uma \textit{game engine}. Da mesma forma, propostas de metodologias narrativas para cenários procedurais (como o modelo \textit{Word2World} \cite{nasir2024word2worldgeneratingstoriesworlds}) confiam em \textit{strings} de texto abstratas processadas e transcodificadas deterministicamente por uma *engine* em tempo de execução para gerar instâncias 2D ou 3D robustas e livres de erro de colisão. A premissa central e o embasamento dessa arquitetura são de que, por trás do ruído natural do aprendizado não determinístico de máquina, haja sempre um filtro algorítmico rigoroso em código tradicional para sustentar a interface com o jogador.

\subsection{Testes de Restrições geométricas e integridade de navegação}
Na camada inferior do paradigma híbrido -- que no projeto é encapsulado pelas lógicas em C\# ou GDScript rodando nativamente sobre as ferramentas do Godot Engine -- a abstração oriunda do JSON transforma-se em vértices instanciáveis. Motores de jogo modernos detêm subsistemas potentes, particularmente motores de física dotados da habilidade de computar \textit{Axis-Aligned Bounding Boxes} (AABBs), hierarquias em volumes limitantes e lançamento de raios espaciais (\textit{Raycasting}). Tais operações compõem a base algorítmica capaz de varrer o plano matemático tridimensional, impedindo sobreposições visuais indesejáveis de malhas (\textit{Meshes}). Em abordagens mistas aplicadas no mercado, a participação humana na interface PCG necessita atuar numa camada de testes estritos como forma de garantir uma disposição geométrica plausível \cite{walton2020evaluating}.
O principal subproduto gerado pelas restrições geométricas, além de assegurar que módulos modulares ou corredores pré-fabricados não intercedam no mesmo polígono de forma impenetrável, é o cálculo da jogabilidade efetiva, muitas vezes verificada na computação da Malha de Navegação (\textit{Navigation Mesh} ou \textit{NavMesh}). Uma NavMesh atua mapeando de forma conexa todos os espaços fisicamente percorríveis de um ambiente virtual para atores virtuais ou controladores de jogador.
O projeto de um jogo não detém o aval do sucesso construtivo se instâncias procedurais falharem no isolamento topológico; caminhos traçados que induzem \textit{dead-ends} imprevisíveis ou rotas não conectadas condenam toda a geração. É exigido que o algoritmo do motor atue aplicando verificações iterativas na fase gerada: se uma sobreposição violar o percurso até o nó de saída estipulado no JSON inicial, algoritmos complementares de resolução de caminhos, como o famoso \textit{A*}, atuarão identificando um bloqueio insolúvel. Tais abordagens teóricas de integridade de malha referendam e justificam a adoção do motor de jogo não apenas como visualizador, mas como a última barreira formal da validação da coerência entre o plano abstrato (IA) e o plano funcional (Jogo) que a pesquisa objetiva integrar.
%\section{Fundamentação Teórica}
\label{sec:fundamentacao}

Este capítulo apresenta a base conceitual necessária para a compreensão da arquitetura do projeto. Inicia-se com a definição formal de fases e \textit{Level Design}, avançando pelas metodologias procedurais clássicas e de aprendizado de máquina. Em seguida, discutem-se as limitações dos modelos de linguagem na ancoragem espacial, culminando na teoria dos grafos de progressão, representação 3D modular e validação de restrições em sistemas híbridos.

\subsection{Fases, \textit{Level Design} e Layouts 3D}
\label{sec:level_design}

Em jogos digitais, uma "fase" (ou nível) compreende o espaço interativo no qual as mecânicas de jogo ocorrem e onde o jogador experimenta os desafios propostos. A disciplina de projetar essas estruturas é conhecida como \textit{Level Design}. O \textit{Level Design} é responsável por ditar o ritmo da experiência \cite{dormans2010}, englobando tanto a elaboração narrativa das missões quanto o arranjo físico dos obstáculos.

Para conceber adequadamente o \textit{design} de um nível, a literatura sugere separar a organização de uma fase em duas camadas independentes. Há uma diferença crucial na literatura de \textit{Level Design} \cite{dormans2010} entre a \textit{missão} (a progressão lógica de tarefas, chaves e bloqueios) e o \textit{espaço} (a topografia física que sedia a missão e que exige um arranjo dimensional). 

No contexto tridimensional, o planejamento desses ambientes materializa-se através de \textit{layouts} 3D. Um \textit{layout} define a disposição espacial arquitetônica dos elementos -- paredes, pisos, plataformas e interações. Planejar um layout 3D exige um rigoroso controle sobre a métrica do espaço virtual para que a progressão planejada na missão não seja corrompida por falhas estruturais.

\subsection{Geração Procedural de Conteúdo (PCG)}
\label{sec:pcg}

Para reduzir o custo de produção e garantir variabilidade, a indústria e a academia têm investido fortemente na Geração Procedural de Conteúdo (PCG -- \textit{Procedural Content Generation}). A PCG refere-se à criação algorítmica autônoma ou semiautônoma de dados e instâncias para jogos. Conforme detalhado na literatura basilar \cite{pcgbook}, o campo pode ser classificado em diferentes abordagens.

Dentre as vertentes mais consolidadas estão os métodos construtivos e os métodos baseados em busca. Os métodos construtivos, como autômatos celulares e gramáticas de formas, constroem o conteúdo de forma determinística passo a passo, garantindo que regras de particionamento e distribuição de espaço nunca sejam violadas. Já as abordagens baseadas em busca geram o conteúdo e o avaliam posteriormente por meio de uma função de aptidão (\textit{fitness function}), selecionando a iteração que melhor resolve o espaço \cite{pcgbook}. Tais algoritmos garantem robustez, contudo, costumam carecer de intencionalidade criativa e de semântica complexa.

\subsection{Procedural Content Generation via Machine Learning (PCGML) e LLMs}
\label{sec:pcgml}

Visando superar a rigidez dos algoritmos clássicos, a comunidade introduziu o \textit{Procedural Content Generation via Machine Learning} (PCGML). No escopo inicial dessas pesquisas, o PCGML foi definido \cite{summerville2018procedural} como o treinamento de modelos probabilísticos a partir de \textit{corpora} de conteúdos de jogos já existentes. A máquina aprende os padrões implícitos no \textit{design} humano e os reproduz.

Com a consolidação da arquitetura \textit{Transformer}, os Modelos de Linguagem de Grande Porte (LLMs) começaram a ser aplicados nesse domínio. Diferente das antigas cadeias de Markov \cite{summerville2018procedural}, os LLMs processam contexto em longa escala e permitem direcionamento via linguagem natural, habilitando a geração de narrativas e conexões que transcendem a pura geometria estocástica. Para que a saída desses modelos seja funcionalmente processada por outras ferramentas computacionais, é indispensável a adoção de uma representação estruturada intermediária. Modelos conceituais arquitetados pela rede são frequentemente serializados em formatos consolidados de troca de dados, como o JSON, atuando como a ponte semântica entre o processamento de linguagem e os sistemas lógicos estritos.

\subsection{O Problema da Ancoragem Espacial (\textit{Spatial Grounding}) em LLMs}
\label{sec:spatial_grounding}

Apesar de sua notável capacidade de interpretação semântica, os LLMs carecem de um raciocínio espacial inato -- a habilidade autônoma de compreender e computar noções absolutas de posição, distância, direção e adjacência no espaço euclidiano. Esta lacuna é formalmente conhecida como o problema da ancoragem espacial (\textit{spatial grounding}). 

Sem o devido raciocínio espacial, um modelo pode descrever a lógica de "um corredor conectando a sala A à sala B", mas falha ao deduzir as coordenadas cartesianas absolutas necessárias para desenhar esse corredor sem invadir o espaço de um cômodo vizinho. Pesquisas recentes voltadas a planejadores explícitos e síntese 3D indicam que, embora o LLM possa sugerir relações qualitativas formidáveis, ele encontra graves dificuldades em satisfazer regras matemáticas contínuas \cite{sun2024roomplanner, hu2024scenecraft}. Portanto, transferir a responsabilidade geométrica irrestrita para uma IA generativa de texto resulta tipicamente em anomalias arquitetônicas.

\subsection{Modelagem de Fases via Grafos de Progressão}
\label{sec:grafos}

Para viabilizar as qualidades do LLM sem esbarrar em suas limitações físicas, a modelagem foca primeiramente na camada de missão, comumente abstraída utilizando a base matemática da Teoria dos Grafos. 

Matematicamente, um grafo é definido como $G = (V, E)$, onde $V$ representa um conjunto de vértices (ou nós) e $E$ representa um conjunto de arestas que conectam esses vértices. Para descrever arranjos e restrições complexas de uma fase, a teoria dos grafos fornece propriedades essenciais: \textit{caminhos} (sequência finita de arestas conectando nós iterativamente), \textit{conectividade} e \textit{alcançabilidade} (garantia de que um nó final pode ser acessado a partir do nó inicial), \textit{ciclos} (caminhos que retornam à origem), \textit{pesos} (custo para transpor uma aresta) e \textit{rótulos} (condições qualitativas atreladas aos nós ou arestas).

Aplicando esse rigor ao design de níveis, cria-se o grafo de progressão lógico ou grafo de ações \cite{inproceedings}. Nesse modelo, os vértices não detêm coordenadas X, Y, Z, mas representam "instâncias de estado" ou salas temáticas. As arestas representam o fluxo de conectividade direcional e os rótulos demarcam as dependências e pré-condições de jogabilidade (por exemplo, "para transpor esta aresta rumo à Sala do Chefe, o rótulo de pré-condição exige a Chave de Ouro"). Ao estruturar a lógica em um grafo abstrato, os requisitos do jogo são matematicamente resolvidos muito antes da instanciação das paredes.

\subsection{Representação de Layouts Espaciais 3D}
\label{sec:layout_3d}

Uma vez estabelecida a estrutura lógica (grafo), o espaço físico deve ser representado computacionalmente para permitir a materialização procedural. A literatura de PCG aborda comumente o espaço físico por meio da representação modular, onde cenários inteiros são compostos por \textit{tiles} (peças de grade bidimensional), módulos maiores, salas pré-construídas ou \textit{prefabs} tridimensionais \cite{linden2014dungeons}.

Nessa representação, cada módulo detém propriedades absolutas essenciais: exata ocupação de volume, orientação vetorial e pontos âncora de encaixe (ex.: os batentes das portas). Para que o \textit{layout} seja estabelecido de forma válida, a instância espacial precisa calcular a translação e rotação global de cada componente modular, monitorando constantemente as intersecções de volumes para certificar-se de que nenhum modelo tridimensional ocupe o mesmo espaço que outro no plano contínuo.

\subsection{Arquiteturas Híbridas e Validação}
\label{sec:arquiteturas_hibridas}

A separação metodológica descrita anteriormente propicia a utilização de sistemas ou arquiteturas híbridas de geração, que combinam geradores estocásticos com restrições operacionais exatas. Trabalhos notórios no campo de iniciativa mista (\textit{Mixed-Initiative}), como o \textit{Tanagra} \cite{tanagra}, evidenciam que a geração baseada em restrições (\textit{Constraint-Based Generation}) é indispensável para conciliar criatividade fluida com coerência matemática.

Em um \textit{pipeline} híbrido neuro-simbólico, o modelo probabilístico propõe o esquema abstrato, e um algoritmo clássico de regras entra em ação na etapa de satisfação de restrições (\textit{Constraint Satisfaction}). Se a diretriz do gerador acarretar uma incompatibilidade física no espaço 3D, a etapa de validação tem a autonomia matemática de recusar a instrução, aplicar rotinas de reparo imediato ou recomputar posições, servindo de filtro definitivo entre a lógica intencional e a execução espacial.

\subsection{Restrições de Nível: Lógica, Geometria, Navegação e Solvabilidade}
\label{sec:validade}

Para assegurar que a geração de um nível foi tecnicamente bem-sucedida, é imperativo avaliar as restrições e qualidades em espectros individuais, separando rigorosamente propriedades estruturais de dinâmicas de jogo \cite{inproceedings, dormans2010}:

\begin{itemize}
    \item \textbf{Validade Geométrica:} Uma propriedade primária. Refere-se estritamente à integridade volumétrica dos elementos gerados. Um nível é geometricamente válido quando não ocorrem sobreposições ilógicas (\textit{clipping}) ou intersecções indesejadas de malhas no espaço euclidiano.
    
    \item \textbf{Navegabilidade:} Propriedade que atesta que o cenário possui superfícies contínuas adequadas para a locomoção de agentes e jogadores (ex: escadas com degraus viáveis, chãos não-fragmentados). Pode ser mapeada em um projeto pela existência de \textit{Navigation Meshes} (\textit{NavMesh}) bem formadas.
    
    \item \textbf{Pathfinding (Busca de Caminho):} O mecanismo algorítmico utilizado para comprovar, na prática, que a malha de navegabilidade realmente permite traçar uma rota ininterrupta entre o ponto A e o ponto B estipulado no grafo.
    
    \item \textbf{Validade Lógica:} A garantia interna de que a estrutura semântica ou o grafo de dependências do nível não contém contradições (ex.: portas que bloqueiam a si mesmas ou ausência de rotas direcionais).
    
    \item \textbf{Solvabilidade:} É a comprovação funcional final de que o jogador consegue vencer a fase cumprindo todas as regras de \textit{gameplay}. 
\end{itemize}

Essa taxonomia é fundamental pois os fenômenos operam de maneira paralela. Uma fase pode deter total \textbf{Validade Geométrica} (nenhuma intersecção de modelo) e permitir perfeito \textbf{Pathfinding} (caminho navegável do início ao fim), mas falhar completamente na \textbf{Solvabilidade} (ex: o jogador necessitar de uma chave vermelha para abrir a porta que abriga essa mesma chave vermelha). Portanto, assegurar um \textit{layout} robusto exige que algoritmos dedicados atuem validando não apenas a matemática do espaço, mas o fluxo temporal da missão gerada pelo modelo.

%% ---------------------------------------------------------------------------
\section{Fundamentação Teórica}
\label{sec:fundamentacao}
% ---------------------------------------------------------------------------

Este capítulo apresenta os fundamentos teóricos que sustentam o estudo da geração automática de fases tridimensionais para jogos digitais. A exposição organiza-se de modo progressivo: parte-se da definição de fases, de \textit{Level Design} e da distinção conceitual entre organização lógica e organização espacial; em seguida, abordam-se os princípios da Geração Procedural de Conteúdo (PCG), sua vertente baseada em aprendizado de máquina (PCGML) e a caracterização dos Modelos de Linguagem de Grande Porte (LLMs), incluindo suas limitações de raciocínio e ancoragem espacial. Por fim, discutem-se a modelagem lógica de fases por meio da Teoria dos Grafos, a representação modular e geométrica de \textit{layouts} tridimensionais, as arquiteturas híbridas guiadas por restrições e os critérios formais de validade, navegabilidade e solvabilidade.

% ---------------------------------------------------------------------------
\subsection{Fases de Jogos, \textit{Level Design} e \textit{Layouts} 3D}
\label{sec:level_design}
% ---------------------------------------------------------------------------

\subsubsection{Fases, Níveis e \textit{Level Design}}
\label{sec:fases_niveis}

No contexto dos jogos digitais, uma fase --- também denominada nível (\textit{level}) ou mapa --- constitui a unidade espacial, temporal e interativa na qual as regras e mecânicas do sistema ludológico se manifestam para o jogador \cite{tanagra, pcgbook}. Mais do que um cenário estático, uma fase opera como um contêiner estruturado para a experiência de jogo (\textit{gameplay}), delimitando o território de exploração, os recursos disponíveis, os obstáculos a serem transpostos e as condições de vitória ou transição \cite{tanagra, linden2014dungeons}. Em gêneros como ação e aventura, jogos de plataforma e \textit{Role-Playing Games} (RPGs), a qualidade da experiência interativa depende diretamente de como esses espaços são construídos e articulados \cite{dormans2010, linden2014dungeons}.

A disciplina responsável pelo planejamento, concepção e estruturação de fases é o \textit{Level Design} \cite{dormans2010, tanagra}. Seus objetivos fundamentais abrangem a condução da atenção e do deslocamento do jogador, a introdução gradual de mecânicas, o balanceamento entre desafio e recompensa e a regulação do ritmo da experiência (\textit{pacing}) \cite{dormans2010, inproceedings, linden2014dungeons}. Em fases do tipo masmorra (\textit{dungeon}), por exemplo, o \textit{Level Design} busca conciliar a liberdade de exploração espacial com um controle rigoroso sobre a progressão de desafios e quebra-cabeças (\textit{puzzles}), assegurando que o jogador não encontre obstáculos intransponíveis para o conjunto de habilidades ou recursos que detém em determinado instante da partida \cite{linden2014dungeons, shaker2016constructive}.

\subsubsection{\textit{Layout} de uma Fase}
\label{sec:layout_fase}

O \textit{layout} de uma fase refere-se ao arranjo estrutural e espacial dos elementos que compõem o ambiente jogável \cite{dormans2010, linden2014dungeons}. Em ambientes internos e masmorras, o \textit{layout} é classicamente decomposto em unidades arquitetônicas discretas, destacando-se as salas (ou câmaras), os corredores (ou vias de transição), os conectores (portas, portais, escadas e elevadores) e as regiões funcionais \cite{linden2014dungeons, shaker2016constructive, xu2025databasedrivenframework3dlevel}. As salas atuam tipicamente como arenas de interação concentrada --- abrigando combates, enigmas, recursos ou eventos narrativos ---, ao passo que os corredores estabelecem as ligações físicas entre essas áreas, modulando o tempo de travessia e a visibilidade entre setores adjacentes \cite{inproceedings, shaker2016constructive}.

Além da geometria estrutural que delimita o espaço navegável (pisos, paredes, tetos e desníveis), o \textit{layout} compreende a disposição interna de elementos funcionais e estéticos, frequentemente denominados instalações ou objetos de cena (\textit{facilities} ou \textit{props}), bem como entidades dinâmicas, tais como personagens não jogáveis (NPCs), inimigos, armadilhas e itens coletáveis \cite{linden2014dungeons, xu2025constraint}. Em ambientes tridimensionais (3D), a especificação de um \textit{layout} envolve não apenas a distribuição planar em dois eixos, mas também a verticalidade, a sobreposição de múltiplos pavimentos, as relações de linha de visão (\textit{line of sight}) e as dimensões volumétricas exatas de cada componente no espaço contínuo \cite{xu2025constraint, xu2025databasedrivenframework3dlevel}.

\subsubsection{Organização Lógica e Organização Espacial}
\label{sec:organizacao_logica_espacial}

Um princípio teórico central na literatura de \textit{Level Design} consiste na distinção analítica entre duas estruturas complementares que coexistem em uma fase: a \textit{missão} (organização lógica) e o \textit{espaço} (organização espacial) \cite{dormans2010, linden2014dungeons}. A missão corresponde à sequência lógica ou parcialmente ordenada de tarefas, ações e objetivos que o jogador deve cumprir para concluir o nível, tais como derrotar um guardião, obter uma chave específica, desativar um bloqueio e alcançar a saída \cite{dormans2010, inproceedings}. O espaço, por sua vez, refere-se à realização física e geométrica do ambiente virtual onde essas tarefas estão situadas, englobando as dimensões das salas, as distâncias métricas e a configuração arquitetônica das conexões \cite{dormans2010, linden2014dungeons}.

Embora o espaço físico acomode a missão e esta seja mapeada sobre aquele, as duas estruturas são conceitualmente independentes e regidas por princípios distintos \cite{dormans2010}. Uma missão estritamente linear --- na qual os objetivos precisam ser cumpridos em uma ordem única e fixa --- pode ser distribuída sobre um espaço físico não linear, como um arranjo em estrela (\textit{hub-and-spoke}) ou em anel, exigindo que o jogador retorne a áreas previamente visitadas à medida que desbloqueia novos acessos \cite{dormans2010}. Inversamente, uma missão não linear, dotada de múltiplos desafios paralelos ou caminhos alternativos, pode ser acomodada em uma geometria espacial simples \cite{dormans2010, inproceedings}. Essa dissociação teórica demonstra que a coerência da progressão de um nível reside primariamente em sua estrutura lógica de dependências, enquanto sua viabilidade física e navegacional depende da correta realização geométrica no espaço tridimensional \cite{dormans2010, inproceedings, linden2014dungeons}.

% ---------------------------------------------------------------------------
\subsection{Geração Procedural de Conteúdo}
\label{sec:pcg}
% ---------------------------------------------------------------------------

\subsubsection{Conceito e Objetivos da PCG}
\label{sec:conceito_pcg}

A Geração Procedural de Conteúdo (PCG --- \textit{Procedural Content Generation}) é definida na literatura como a criação algorítmica de conteúdos para jogos digitais de maneira autônoma ou com intervenção humana limitada e indireta \cite{pcgbook, summerville2018procedural, maleki2024procedural}. O escopo do termo ``conteúdo'' abrange praticamente todos os artefatos que integram um jogo, incluindo níveis, mapas, terrenos, malhas vegetais, texturas, efeitos sonoros, narrativas, missões (\textit{quests}) e conjuntos de regras, distinguindo-se do comportamento em tempo de execução dos agentes de inteligência artificial que controlam NPCs \cite{pcgbook, summerville2018procedural, maleki2024procedural}.

Historicamente introduzida em títulos pioneiros da década de 1980, como \textit{Rogue} (1980) e \textit{Elite} (1984), para contornar severas restrições de memória e armazenamento por meio da síntese algorítmica em tempo de execução \cite{pcgbook, summerville2018procedural, araujo2018geracao}, a PCG consolidou-se tanto na indústria quanto na pesquisa acadêmica em função de múltiplos objetivos \cite{linden2014dungeons, pcgbook}. Dentre as principais motivações para o seu emprego destacam-se: (i) a automação e a redução do tempo e do custo financeiro associados à autoria manual de ativos em larga escala; (ii) o aumento da variabilidade e da diversidade dos ambientes gerados, ampliando o fator de rejogabilidade (\textit{replayability}); (iii) o suporte à adaptabilidade centrada no jogador, permitindo ajustar dinamicamente os desafios às habilidades observadas; e (iv) a exploração de novas estéticas e mecânicas emergentes \cite{linden2014dungeons, pcgbook, summerville2018procedural}.

A avaliação e o projeto de geradores procedurais apoiam-se em propriedades técnicas fundamentais, notadamente a reprodutibilidade, a expressividade e a controlabilidade \cite{tanagra, linden2014dungeons, maleki2024procedural}. A reprodutibilidade decorre frequentemente do uso de Geradores de Números Pseudoaleatórios (PRNGs), nos quais a especificação de um valor inicial --- denominado semente (\textit{seed}) --- permite reconstruir de maneira determinística a mesma sequência de decisões algorítmicas \cite{maleki2024procedural}. A expressividade, frequentemente analisada por meio do espaço expressivo (\textit{expressive range}), caracteriza a variedade e as propriedades estilísticas ou estruturais dos artefatos que o gerador é capaz de produzir \cite{tanagra, inproceedings}. A controlabilidade, por fim, refere-se ao grau de domínio oferecido ao projetista ou ao sistema para direcionar intencionalmente as características do conteúdo gerado por meio de parâmetros, regras ou restrições, garantindo resultados consistentes sem sacrificar a diversidade \cite{linden2014dungeons, shaker2016constructive}.

\subsubsection{Principais Abordagens de PCG}
\label{sec:abordagens_pcg}

A literatura clássica de PCG classifica os métodos generativos em três grandes paradigmas operacionais: métodos construtivos, métodos baseados em busca (\textit{search-based}) e métodos baseados em restrições ou solucionadores (\textit{solver-based} / \textit{constraint-based}) \cite{pcgbook, summerville2018procedural, maleki2024procedural}.

Os métodos construtivos caracterizam-se por gerar o conteúdo em uma única passagem sequencial de passos algorítmicos, com tempo de execução previsível e sem realizar ciclos de avaliação e regeneração sobre o artefato produzido \cite{shaker2016constructive, summerville2018procedural}. Exemplos clássicos incluem funções de ruído coerente (como o ruído de Perlin para síntese de terrenos e texturas), autômatos celulares, algoritmos de particionamento espacial e gramáticas generativas \cite{linden2014dungeons, shaker2016constructive, maleki2024procedural}. Importa ressaltar que, embora procedimentos construtivos sejam computacionalmente rápidos e possam operar de forma determinística a partir de uma semente fixa, o determinismo algorítmico assegura apenas a repetibilidade da saída para uma mesma entrada; ele não garante, isoladamente, que o conteúdo resultante seja estruturalmente válido, conectado ou solucionável pelo jogador, a menos que as regras construtivas tenham sido estritamente projetadas para preservar tais invariantes em todas as etapas \cite{linden2014dungeons, shaker2016constructive}.

Os métodos baseados em busca (\textit{search-based PCG}) enquadram-se como um caso especial do paradigma gerar-e-testar (\textit{generate-and-test}) \cite{pcgbook, maleki2024procedural}. Em vez de aceitar imediatamente o artefato sintetizado ou simplesmente descartá-lo de forma binária, algoritmos baseados em busca atribuem a cada instância candidata um valor numérico contínuo por meio de uma função de avaliação ou aptidão (\textit{fitness function}), utilizando esse escore para guiar iterativamente a modificação e a seleção de novas soluções no espaço de busca \cite{linden2014dungeons, araujo2018geracao, maleki2024procedural}. Essa categoria emprega amplamente técnicas de computação evolutiva (como Algoritmos Genéticos) \cite{linden2014dungeons, maleki2024procedural}, além de meta-heurísticas como têmpera simulada (\textit{Simulated Annealing}), otimização por enxame de partículas (PSO) e Busca em Árvore de Monte Carlo (MCTS) \cite{maleki2024procedural, xu2025constraint}.

Os métodos baseados em restrições (\textit{constraint-based} ou \textit{solver-based PCG}), por sua vez, especificam declarativamente as propriedades e relações obrigatórias que o artefato final deve satisfazer, delegando a construção a solucionadores matemáticos ou lógicos (\textit{constraint solvers}), tais como programação por conjuntos de respostas (\textit{Answer Set Programming} --- ASP), solucionadores de satisfatibilidade módulo teorias (SMT) e bibliotecas de programação por restrições \cite{tanagra, linden2014dungeons, pcgbook, summerville2018procedural, chen2025narrative}, além de algoritmos de síntese baseados em propagação local de padrões e restrições, como o \textit{Wave Function Collapse} (WFC) \cite{maleki2024procedural}. Nessa abordagem, o foco desloca-se da definição de passos procedurais imperativos para a formulação rigorosa do espaço de design admissível \cite{tanagra, linden2014dungeons}.

\subsubsection{Geração Procedural de Fases}
\label{sec:pcg_fases}

A aplicação de PCG à criação de fases e níveis constitui um dos tópicos mais investigados na área, apresentando exigências técnicas particulares que a diferenciam da geração de ativos puramente cosméticos \cite{summerville2018procedural, maleki2024procedural}. A literatura distingue conteúdos cosméticos (como texturas e efeitos visuais, cuja alteração não modifica o resultado das ações do jogador) de conteúdos funcionais, que afetam diretamente a mecânica e a dinâmica da partida \cite{summerville2018procedural}. Fases de jogos são artefatos intrinsecamente funcionais: uma fase visualmente atraente, mas que apresente falhas topológicas ou bloqueios estruturais que impeçam sua conclusão, é considerada inválida para o propósito do jogo \cite{summerville2018procedural, linden2014dungeons}.

Na geração procedural de níveis e masmorras, um método típico envolve três componentes conceituais: (1) um modelo representacional abstrato que descreve a estrutura global da fase; (2) um método para construir esse modelo representacional; e (3) um método para sintetizar a geometria efetiva da fase a partir do modelo abstrato \cite{linden2014dungeons, shaker2016constructive}. Diversas famílias algorítmicas foram desenvolvidas para essas etapas. Algoritmos de particionamento de espaço, como o Particionamento Binário do Espaço (\textit{Binary Space Partitioning} --- BSP), \textit{quadtrees} (em 2D) e \textit{octrees} (em 3D), subdividem recursivamente o domínio espacial em células disjuntas, garantindo estruturalmente que salas alocadas em partições distintas não se sobreponham \cite{shaker2016constructive}. Em contrapartida, abordagens baseadas em agentes escavadores (\textit{agent-based dungeon growing}) operam em escala microscópica, deslocando um agente virtual que escava túneis e posiciona salas estocasticamente; embora produzam configurações orgânicas, agentes sem mecanismos de inspeção antecipada (\textit{look-ahead}) estão sujeitos a gerar salas sobrepostas e corredores sem saída \cite{shaker2016constructive}.

De modo complementar, autômatos celulares operam sobre grades discretas aplicando regras de transição baseadas na vizinhança local (como as vizinhanças de Moore ou de von Neumann) para formar cavernas naturais, mas não asseguram, por si sós, a conectividade global entre câmaras isoladas, exigindo etapas adicionais de verificação e abertura de passagens \cite{linden2014dungeons, shaker2016constructive}. Gramáticas generativas --- incluindo gramáticas de grafos e gramáticas de formas (\textit{shape grammars}) --- expandem um axioma inicial por meio de regras de produção e reescrita, permitindo capturar princípios de progressão e hierarquia em níveis de ação e aventura \cite{dormans2010, inproceedings, shaker2016constructive}. Apesar da maturidade dessas técnicas em domínios bidimensionais, a transição para a geração de fases tridimensionais complexas permanece um desafio aberto na literatura devido ao aumento expressivo das restrições geométricas, de navegação vertical e de acoplamento espacial \cite{linden2014dungeons, maleki2024procedural, xu2025constraint}.

% ---------------------------------------------------------------------------
\subsection{PCG Baseada em Aprendizado de Máquina e Uso de LLMs}
\label{sec:pcgml_llms}
% ---------------------------------------------------------------------------

\subsubsection{\textit{Procedural Content Generation via Machine Learning} --- PCGML}
\label{sec:pcgml}

A Geração Procedural de Conteúdo via Aprendizado de Máquina (\textit{Procedural Content Generation via Machine Learning} --- PCGML) é definida formalmente como a geração de conteúdos de jogos diretamente por modelos estatísticos ou de aprendizado de máquina que foram treinados sobre conjuntos de conteúdos de jogos preexistentes \cite{summerville2018procedural}. A distinção fundamental entre o PCGML e as abordagens clássicas baseadas em busca ou em solucionadores reside no mecanismo gerador: enquanto a busca evolucionária pode eventualmente empregar uma rede neural como função de avaliação para pontuar candidatos em um espaço de busca codificado manualmente, no PCGML o modelo treinado aprende a distribuição latente dos dados de treinamento e produz novas instâncias de conteúdo diretamente a partir de sua inferência \cite{summerville2018procedural}.

A taxonomia clássica de PCGML organiza os métodos segundo duas dimensões: a representação subjacente dos dados (sequências unidimensionais, grades bidimensionais/tridimensionais e grafos) e o paradigma de treinamento ou modelagem empregado \cite{summerville2018procedural}. Entre as principais técnicas investigadas na literatura figuram: (i) modelos probabilísticos baseados em contagem de frequência, como modelos de $n$-gramas, Cadeias de Markov Multidimensionais (MdMCs) e Campos Aleatórios de Markov (MRFs), utilizados para amostrar níveis fatia por fatia ou bloco por bloco; (ii) métodos de fatoração de matrizes e redução de dimensionalidade, como Análise de Componentes Principais (PCA) e fatoração de matrizes não negativas; (iii) Redes Neurais Recorrentes, em especial redes \textit{Long Short-Term Memory} (LSTMs), aplicadas à geração sequencial de fases de plataforma; (iv) \textit{Autoencoders} e \textit{Variational Autoencoders} (VAEs), capazes de comprimir níveis em espaços latentes contínuos para geração, interpolação e reparo de segmentos inválidos; e (v) Redes Adversárias Generativas (GANs) e Aprendizado por Reforço (PCGRL) \cite{summerville2018procedural, maleki2024procedural}.

Embora o PCGML reduza a necessidade de codificar manualmente regras de estilo ao extrair padrões diretamente de exemplos de design humano, a literatura aponta limitações estruturais importantes nesse paradigma \cite{summerville2018procedural, maleki2024procedural, her2025zeroshot3dmapgeneration}. A primeira refere-se à escassez crônica de dados (\textit{data scarcity}): diferentemente de domínios como visão computacional ou processamento de texto, jogos digitais dispõem de \textit{corpora} públicos reduzidos e carecem de formatos de representação padronizados entre diferentes títulos \cite{summerville2018procedural}. A segunda limitação diz respeito à dificuldade de assegurar restrições funcionais globais: modelos puramente estatísticos treinados sobre matrizes visuais ou sequências locais de blocos frequentemente reproduzem padrões estéticos plausíveis, mas falham em garantir propriedades de longo alcance, como a conectividade global e a solvabilidade do nível sem o auxílio de restrições externas \cite{summerville2018procedural, maleki2024procedural}.

\subsubsection{Modelos de Linguagem de Grande Porte Aplicados à Geração de Conteúdo}
\label{sec:llms_pcg}

Os Modelos de Linguagem de Grande Porte (\textit{Large Language Models} --- LLMs) constituem uma classe de modelos fundacionais de aprendizado profundo baseados na arquitetura \textit{Transformer}, introduzida para superar os gargalos computacionais e de memória dos modelos sequenciais precursores \cite{Gallotta_2024, maleki2024procedural}. Na evolução do Processamento de Linguagem Natural (PLN), abordagens estatísticas iniciais baseadas em cadeias de Markov e $n$-gramas deram lugar a representações vetoriais de palavras (\textit{word embeddings}) e a Redes Neurais Recorrentes (RNNs e LSTMs) \cite{Gallotta_2024, maleki2024procedural}. Contudo, como as arquiteturas recorrentes processam a informação sequencialmente passo a passo, elas sofrem de limitações na captura de dependências de longo alcance, associadas ao problema do desaparecimento de gradientes (\textit{vanishing gradient problem}) em sequências extensas \cite{maleki2024procedural}.

A arquitetura \textit{Transformer} contorna essas limitações ao substituir a recorrência sequencial pelo mecanismo de autoatenção (\textit{self-attention}) \cite{Gallotta_2024, maleki2024procedural}. Em vez de processar a entrada estritamente elemento por elemento, o mecanismo de autoatenção pondera dinamicamente a relevância mútua entre todos os elementos (\textit{tokens}) de uma sequência, permitindo ao modelo relacionar diferentes partes do contexto de entrada e capturar simultaneamente relações sintáticas, semânticas e estruturais de longo alcance \cite{Gallotta_2024, maleki2024procedural}.

Na literatura de modelos de linguagem baseados em \textit{Transformers}, destacam-se duas famílias principais: modelos bidirecionais de codificação (como o BERT, otimizados por modelagem de linguagem mascarada para extração de características e preenchimento contextual) \cite{Gallotta_2024, maleki2024procedural} e modelos autorregressivos (como a família GPT \cite{Gallotta_2024, maleki2024procedural} e modelos abertos como LLaMA e Mistral \cite{Gallotta_2024}, voltados à geração de texto e código). Na formulação autorregressiva causal, o modelo estima a distribuição de probabilidade condicional do próximo \textit{token} $x_t$ dado o contexto precedente $(x_1, x_2, \dots, x_{t-1})$ \cite{Gallotta_2024}. O desenvolvimento de um LLM moderno compreende tipicamente o pré-treinamento autossupervisionado em larga escala sobre vastos \textit{corpora} textuais, seguido pelo ajuste fino por instruções (\textit{instruction fine-tuning}) e pelo alinhamento com preferências humanas, notadamente via Aprendizado por Reforço com \textit{Feedback} Humano (RLHF) \cite{Gallotta_2024, her2025zeroshot3dmapgeneration}.

Em nível de inferência, o escalonamento para centenas de milhões ou bilhões de parâmetros faz emergir capacidades como o aprendizado em contexto (\textit{In-Context Learning} --- ICL) \cite{ocal2024sceneteller, her2025zeroshot3dmapgeneration}, pelo qual o modelo resolve novas tarefas condicionadas apenas pelas instruções do \textit{prompt} (\textit{zero-shot}) ou por poucos exemplos demonstrativos de entrada e saída (\textit{few-shot}), sem necessidade de atualização de seus pesos \cite{Gallotta_2024, ocal2024sceneteller, her2025zeroshot3dmapgeneration}. Estratégias complementares de engenharia de \textit{prompt} e acoplamento externo incluem a Cadeia de Pensamento (\textit{Chain-of-Thought} --- CoT), que decompõe problemas complexos em etapas explícitas de raciocínio intermediário, e a Geração Aumentada por Recuperação (\textit{Retrieval-Augmented Generation} --- RAG), que integra bases de conhecimento ou bancos de dados vetoriais externos ao contexto de geração \cite{Gallotta_2024, tang2025unrealllm, her2025zeroshot3dmapgeneration, wang2026designerlyloop}.

No domínio de PCG, a introdução dos LLMs mitigou parcialmente o problema da escassez de conjuntos de dados específicos de jogos ao transferir o conhecimento semântico e estrutural adquirido durante o pré-treinamento geral \cite{Gallotta_2024, todd2023level, maleki2024procedural}. A literatura documenta a aplicação de LLMs na geração de narrativas interativas, diálogos e árvores de missões; na síntese direta de níveis 2D codificados como grades de caracteres (a exemplo de \textit{Sokoban}, \textit{Super Mario Bros.} e mundos baseados em \textit{tiles}) \cite{todd2023level, nasir2024word2worldgeneratingstoriesworlds, maleki2024procedural}; e na atuação como planejadores semânticos de alto nível, capazes de traduzir intenções em linguagem natural em especificações de cena, parâmetros de ferramentas procedurais e regras de progressão \cite{Gallotta_2024, hu2024scenecraft, tang2025unrealllm, xu2025constraint, her2025zeroshot3dmapgeneration}.

\subsubsection{Geração de Representações Estruturadas}
\label{sec:representacoes_estruturadas}

Quando modelos de linguagem são integrados a sistemas computacionais mais amplos --- tais como motores gráficos, otimizadores matemáticos ou verificadores formais ---, a saída em linguagem natural livre torna-se inadequada devido à sua ambiguidade inerente e à ausência de garantias sintáticas para interpretação automatizada (\textit{parsing}) \cite{xu2025constraint, regmi2025system}. Para viabilizar a interoperabilidade entre o raciocínio probabilístico de um LLM e módulos algorítmicos determinísticos, a literatura recorre à geração de representações estruturadas intermediárias \cite{hu2024scenecraft, tang2025unrealllm, xu2025constraint, chen2025narrative, wang2026designerlyloop}.

Sob o ponto de vista teórico e de engenharia de software, é fundamental distinguir dois níveis conceituais nessa comunicação: o \textit{modelo conceitual} (ou estrutura abstrata de informação) e o \textit{formato de serialização} \cite{xu2025constraint, xu2025databasedrivenframework3dlevel, chen2025narrative}. O modelo conceitual corresponde à abstração matemática ou simbólica que organiza o domínio do problema, podendo assumir a forma de um grafo de cena, de raciocínio ou de progressão \cite{xu2025databasedrivenframework3dlevel, wang2026designerlyloop}, de um conjunto de triplas relacionais do tipo \texttt{[Objeto-Relação-Objeto]} \cite{chen2025narrative}, de uma árvore hierárquica de componentes, ou de sentenças formais em uma Linguagem de Descrição de Domínio Específico (DSL), como a \textit{Game Level Description Language} (GLDL) \cite{xu2025constraint} ou códigos de síntese programática \cite{hu2024scenecraft, yang2026codeasroomgenerating3drooms}.

O formato de serialização, por outro lado, consiste na notação textual padronizada utilizada para codificar, armazenar e transmitir esse modelo conceitual entre processos distintos \cite{sun2024roomplanner, xu2025databasedrivenframework3dlevel}. O formato JSON (\textit{JavaScript Object Notation}), assim como XML, YAML ou folhas de estilo estruturadas (CSS) \cite{ocal2024sceneteller}, constitui um exemplo amplamente adotado de representação serializada por estruturar dados em pares chave-valor e listas hierárquicas facilmente validáveis por esquemas formais (\textit{schemas}) \cite{tang2025unrealllm, sun2024roomplanner, xu2025databasedrivenframework3dlevel}. Portanto, o formato JSON não se confunde com o modelo lógico da fase em si: ele atua como o veículo sintático que codifica a abstração conceitual (por exemplo, os vértices, arestas e atributos de um grafo) para que sistemas externos possam processá-la sem ambiguidade \cite{tang2025unrealllm, xu2025databasedrivenframework3dlevel}.

% ---------------------------------------------------------------------------
\subsection{Raciocínio e Ancoragem Espacial em LLMs}
\label{sec:spatial_grounding}
% ---------------------------------------------------------------------------

\subsubsection{Raciocínio Espacial}
\label{sec:raciocinio_espacial}

O raciocínio espacial compreende a capacidade de representar mental ou computacionalmente objetos em um espaço dimensional, inferir suas propriedades geométricas e deduzir as consequências de suas relações e transformações mútuas \cite{sun2024roomplanner, xu2025constraint}. No contexto da síntese de ambientes virtuais e \textit{layouts} tridimensionais, o raciocínio espacial abrange tanto relações qualitativas e topológicas quanto relações métricas e contínuas \cite{xu2025constraint, chen2025narrative}.

As relações espaciais podem ser categorizadas em propriedades fundamentais que governam a organização de uma cena \cite{xu2025constraint, sun2024roomplanner, chen2025narrative}: (i) \textit{posição}, que define a localização exata de uma entidade em relação a um referencial de coordenadas global ou local; (ii) \textit{distância e proximidade}, que quantificam o afastamento métrico euclidiano entre pares de elementos (como estar próximo ou distante de um ponto de referência); (iii) \textit{direção e orientação}, que determinam para qual vetor ou ângulo a face frontal de um objeto está voltada (\textit{facing/focus}) e como seus eixos locais se alinham aos eixos do ambiente; (iv) \textit{adjacência e conectividade}, que indicam contato imediato ou comunicação entre regiões vizinhas (como duas salas que compartilham uma parede ou porta); (v) \textit{contenção}, que exige que objetos internos permaneçam estritamente dentro dos limites volumétricos de um compartimento maior; e (vi) \textit{interseção e ocupação espacial}, que descrevem se dois volumes tridimensionais compartilham pontos comuns no espaço \cite{xu2025constraint, sun2024roomplanner, xu2025databasedrivenframework3dlevel}. Manter a consistência simultânea de todas essas relações quando dezenas de elementos interagem em três dimensões constitui um problema combinatório e geométrico de elevada complexidade \cite{xu2025constraint, feng2026repurposing3dgenerativemodel}.

\subsubsection{\textit{Spatial Grounding}}
\label{sec:conceito_spatial_grounding}

Na interseção entre inteligência artificial simbólica, visão computacional e processamento de linguagem natural, o conceito de ancoragem (\textit{grounding}) refere-se ao estabelecimento de um vínculo semântico e operacional entre símbolos abstratos (como palavras, predicados ou \textit{tokens}) e seus referentes no mundo físico ou em um ambiente simulado \cite{Gallotta_2024, chen2025narrative}. Especificamente, a ancoragem espacial (\textit{spatial grounding}) consiste no processo de mapear descrições linguísticas e conceitos relacionais qualitativos --- tais como ``posicionar uma mesa de centro em frente ao sofá'', ``alocar uma estante encostada à parede norte'' ou ``conectar a sala de guarda ao corredor principal'' --- em parâmetros espaciais concretos, coordenadas métricas, escalas proporcionais e restrições físicas verificáveis dentro de uma representação geométrica 2D ou 3D \cite{hu2024scenecraft, sun2024roomplanner, chen2025narrative, ocal2024sceneteller}.

Uma ancoragem espacial efetiva exige que o sistema generativo não apenas compreenda a associação semântica e funcional (\textit{affordance}) entre os objetos mencionados no texto (por exemplo, saber que cadeiras costumam acompanhar mesas ou que portas conectam compartimentos) \cite{Gallotta_2024, chen2025narrative}, mas também traduza essas intenções em configurações que respeitem as dimensões reais dos modelos tridimensionais (\textit{bounding boxes}), as leis de suporte físico, a orientação funcional dos componentes e a preservação de áreas livres para circulação e interação \cite{sun2024roomplanner, xu2025constraint, ocal2024sceneteller}.

\subsubsection{Limitações dos LLMs em Tarefas Espaciais Tridimensionais}
\label{sec:limitacoes_espaciais_llms}

Em modelos generativos de linguagem, denomina-se ``alucinação'' a síntese de conteúdos que aparentam fluência e plausibilidade sintática, mas que são factualmente incorretos, contradictórios ou desprovidos de aderência às restrições do domínio \cite{Gallotta_2024, xu2025constraint, bandeali2026generative}. Quando LLMs são aplicados diretamente à geração de cenas e níveis tridimensionais --- sendo solicitados a produzir coordenadas numéricas $(x, y, z)$, dimensões e rotações para múltiplos ativos ---, esse fenômeno manifesta-se como \textit{alucinação espacial} (ou inconsistência geométrica) \cite{Gallotta_2024, sun2024roomplanner, yang2026codeasroomgenerating3drooms, feng2026repurposing3dgenerativemodel}.

A literatura recente analisa múltiplos fatores técnicos que explicam as limitações dos LLMs em tarefas espaciais 3D \cite{Gallotta_2024, hu2024scenecraft, sun2024roomplanner, feng2026repurposing3dgenerativemodel, yang2026codeasroomgenerating3drooms}. Embora os LLMs apresentem forte capacidade de planejamento semântico e extração de relações qualitativas, seu pré-treinamento ocorre sobre distribuições estatísticas textuais, sem acoplamento nativo a um verificador de geometria euclidiana ou motor de simulação física durante a inferência \cite{Gallotta_2024, sun2024roomplanner, feng2026repurposing3dgenerativemodel}. Além disso, ao tratar a síntese de \textit{layout} puramente como uma tarefa de geração textual, ocorre uma severa \textit{compressão de informação}: representar objetos tridimensionais complexos apenas por descritores textuais e caixas delimitadoras omite interações geométricas finas de superfície \cite{feng2026repurposing3dgenerativemodel}. Adicionalmente, em formulações autorregressivas ou em planejamentos de passo único (\textit{single-step planning}), verifica-se o fenômeno de \textit{viés de exposição} (\textit{exposure bias}) e acúmulo de erros: pequenas imprecisões métricas cometidas nas primeiras predições de coordenadas propagam-se ao longo da sequência gerada, degradando a coerência espacial global à medida que o número de objetos ou salas aumenta \cite{ocal2024sceneteller, sun2024roomplanner, feng2026repurposing3dgenerativemodel, yang2026codeasroomgenerating3drooms}.

Estudos empíricos na área de geração de cenas e níveis 3D documentam falhas recorrentes associadas a essa limitação de ancoragem espacial \cite{ocal2024sceneteller, hu2024scenecraft, sun2024roomplanner, xu2025constraint, feng2026repurposing3dgenerativemodel, yang2026codeasroomgenerating3drooms, bandeali2026generative}:
\begin{itemize}
    \item \textbf{Colisões e sobreposições volumétricas (\textit{overlapping} / \textit{clipping}):} atribuição de coordenadas cujos volumes delimitadores se interceptam ou se interpenetram, fazendo com que móveis, paredes ou salas inteiras ocupem simultaneamente a mesma região do espaço tridimensional \cite{ocal2024sceneteller, sun2024roomplanner, xu2025constraint, feng2026repurposing3dgenerativemodel};
    \item \textbf{Violações de suporte e de limites (\textit{out-of-bounds} e objetos flutuantes):} posicionamento de elementos fora do perímetro das salas, atravessando paredes estruturais ou suspensos no ar sem suporte físico adequado \cite{ocal2024sceneteller, sun2024roomplanner, xu2025constraint, feng2026repurposing3dgenerativemodel, bandeali2026generative};
    \item \textbf{Erros de orientação e alinhamento:} rotação incorreta de elementos funcionais (como portas, escadas ou mobiliário voltados contra paredes cegas) e desalinhamento entre aberturas de cômodos adjacentes \cite{sun2024roomplanner, xu2025constraint, yang2026codeasroomgenerating3drooms};
    \item \textbf{Obstrução topológica e perda de acessibilidade:} alocação de objetos ou estruturas que bloqueiam passagens críticas, portas ou corredores, fragmentando o espaço navegável e isolando itens essenciais à progressão \cite{sun2024roomplanner, xu2025constraint, xu2025databasedrivenframework3dlevel}.
\end{itemize}

Tais evidências na literatura demonstram que, embora os LLMs sejam ferramentas expressivas para a concepção semântica e estrutural de ambientes, a geração direta e irrestrita de coordenadas tridimensionais apenas por inferência textual mostra-se frequentemente insuficiente para assegurar a validade física e funcional de fases 3D, motivando o emprego de abstrações intermediárias, planejadores explícitos e algoritmos de otimização ou verificação espacial \cite{hu2024scenecraft, sun2024roomplanner, xu2025constraint, xu2025databasedrivenframework3dlevel, yang2026codeasroomgenerating3drooms}.

% ---------------------------------------------------------------------------
\subsection{Representação Lógica de Fases por Grafos}
\label{sec:grafos}
% ---------------------------------------------------------------------------

\subsubsection{Conceitos Fundamentais de Teoria dos Grafos}
\label{sec:conceitos_grafos}

A Teoria dos Grafos fornece o arcabouço matemático padrão para modelar relações discretas entre entidades \cite{summerville2018procedural}. Formalmente, um grafo é definido como um par ordenado $G = (V, E)$, no qual $V$ constitui um conjunto finito de vértices (ou nós) e $E$ representa um conjunto de arestas (ou arcos) que estabelecem ligações entre pares de vértices de $V$ \cite{maleki2024procedural, yang2026codeasroomgenerating3drooms}.

De acordo com a natureza de suas conexões, um grafo classifica-se como \textit{não direcionado}, quando suas arestas representam ligações simétricas e bidirecionais entre pares de vértices, ou como \textit{direcionado} (dígrafo), quando cada aresta possui uma orientação explícita que parte de um vértice de origem em direção a um vértice de destino \cite{dormans2010, summerville2018procedural}. O número de arestas conectadas a um vértice define o seu \textit{grau} \cite{inproceedings, xu2025databasedrivenframework3dlevel}. Adicionalmente, quando valores numéricos são associados às arestas, obtém-se um \textit{grafo ponderado}, no qual os pesos podem representar custos de travessia, distâncias ou níveis de dificuldade; quando símbolos, tipos ou conjuntos de propriedades são associados aos vértices ou arestas, tem-se um \textit{grafo rotulado} ou \textit{grafo de atributos} \cite{dormans2010, inproceedings, linden2014dungeons, chen2025narrative}.

Sobre essa estrutura definem-se conceitos topológicos essenciais para a análise de redes e níveis de jogos \cite{inproceedings, linden2014dungeons, summerville2018procedural}:
\begin{itemize}
    \item \textbf{Caminho (\textit{path}):} uma sequência finita de vértices conectados sucessivamente por arestas de $E$, permitindo transitar de um nó de origem a um nó de destino \cite{inproceedings, linden2014dungeons};
    \item \textbf{Ciclo (\textit{cycle}):} um caminho fechado que parte de um vértice e retorna a ele mesmo após percorrer uma sequência de arestas \cite{dormans2010, linden2014dungeons}. Um grafo desprovido de ciclos é denominado acíclico, destacando-se as \textit{árvores} (grafos conexos e acíclicos) e os Grafos Acíclicos Direcionados (DAGs) \cite{linden2014dungeons, wang2026designerlyloop};
    \item \textbf{Conectividade:} propriedade estrutural que indica a existência de caminhos ligando os vértices do grafo, assegurando que o conjunto de nós forme uma rede integrada e sem subcomponentes isolados \cite{linden2014dungeons, summerville2018procedural, xu2025databasedrivenframework3dlevel};
    \item \textbf{Alcançabilidade (\textit{reachability}):} propriedade que atesta se um vértice de destino pode ser efetivamente atingido a partir de um vértice inicial por meio de um caminho válido em $G$ \cite{dormans2010, linden2014dungeons, xu2025databasedrivenframework3dlevel};
    \item \textbf{Planaridade:} propriedade de um grafo que pode ser representado em um plano bidimensional de modo que suas arestas não se cruzem fora dos vértices onde incidem \cite{inproceedings}.
\end{itemize}

\subsubsection{Grafos Aplicados ao \textit{Level Design}}
\label{sec:grafos_level_design}

Na literatura de \textit{Level Design} e de geração procedural de masmorras, os grafos são amplamente adotados como modelos representacionais devido à sua capacidade de codificar relações estruturais de forma compacta \cite{dormans2010, inproceedings, linden2014dungeons, summerville2018procedural}. Essa modelagem manifesta-se principalmente sob duas perspectivas complementares: grafos topológicos de espaço e grafos de missão ou de ações \cite{dormans2010, inproceedings, linden2014dungeons}.

Nos \textit{grafos topológicos de espaço} (frequentemente chamados de grafos de salas ou redes de conectividade espacial), cada vértice $v \in V$ representa uma região física discreta do nível --- como uma sala, uma câmara, um pátio ou uma plataforma ---, enquanto cada aresta $e \in E$ representa uma ligação física de adjacência ou passagem entre duas regiões, tais como portas, corredores, pontes, escadas ou portais de teletransporte \cite{dormans2010, linden2014dungeons, summerville2018procedural}. Atributos associados aos vértices descrevem o tipo funcional ou temático da sala (ex.: sala inicial, sala de combate, sala de tesouro, sala de chefe) e os objetos nela contidos, ao passo que atributos nas arestas descrevem a natureza da transição (ex.: passagem aberta, porta trancada, passagem unidirecional) \cite{linden2014dungeons, xu2025databasedrivenframework3dlevel}.

Nos \textit{grafos de missão} ou \textit{grafos de ações} (\textit{mission graphs} / \textit{action graphs}), por sua vez, os vértices representam tarefas, eventos narrativos ou ações que o jogador deve realizar (por exemplo, ``derrotar inimigo'', ``coletar chave'', ``destrancar porta'', ``interagir com altar''), e as arestas direcionadas especificam as relações de precedência causal ou temporal entre essas ações \cite{dormans2010, inproceedings}. Trabalhos clássicos demonstram que um grafo de ações pode ser progressivamente transformado em um grafo de salas por meio do agrupamento de ações que devem ocorrer no mesmo local (\textit{co-located actions}), servindo como ponte formal entre o planejamento de \textit{gameplay} e a topologia do espaço \cite{inproceedings, linden2014dungeons}.

\subsubsection{Grafos de Progressão e Dependências}
\label{sec:grafos_progressao}

A combinação de topologia espacial com regras de precedência origina os \textit{grafos de progressão}, os quais modelam explicitamente os estados, pré-condições e dependências lógicas que governam o avanço do jogador através da fase \cite{dormans2010, inproceedings, linden2014dungeons, xu2025databasedrivenframework3dlevel}. O mecanismo canônico para expressar dependências em fases de aventura e exploração é a estrutura de chave e fechadura (\textit{lock-and-key puzzle}) \cite{dormans2010, linden2014dungeons, xu2025databasedrivenframework3dlevel}. Conceitualmente, uma ``fechadura'' representa qualquer barreira condicional associada a uma aresta ou nó (uma porta trancada, um abismo, um obstáculo destrutível ou um guardião), enquanto a ``chave'' corresponde ao recurso, item, habilidade ou estado (uma chave física, uma alavanca, um item como bombas ou um artefato especial) cuja obtenção prévia é pré-condição obrigatória para transpor a barreira \cite{dormans2010, inproceedings, shaker2016constructive}.

Em um grafo de progressão, essas dependências impõem restrições de ordem topológica ($\tau$): se uma aresta ou sala $R_L$ está condicionada a um bloqueio que exige um recurso situado na sala $R_K$, a precedência lógica requer que $R_K$ seja alcançável a partir do vértice inicial sem que seja necessário atravessar $R_L$ ($\tau(R_K) < \tau(R_L)$) \cite{xu2025databasedrivenframework3dlevel}. A organização dessas arestas e dependências define padrões estruturais importantes para o ritmo da fase: a rota crítica (caminho obrigatório entre o início e o objetivo final), as ramificações opcionais (subgrafos laterais que abrigam recompensas secundárias ou recursos auxiliares sem obstruir o fluxo principal) e os ciclos controlados, que permitem criar atalhos ou exigir retorno estratégico (\textit{backtracking}) após a obtenção de uma chave \cite{dormans2010, inproceedings, linden2014dungeons}.

\subsubsection{Estrutura Lógica versus Representação Espacial}
\label{sec:logica_vs_espacial}

Um aspecto teórico fundamental na modelagem por grafos reside no fato de que um grafo é uma estrutura puramente topológica e combinatória, e não uma representação métrica ou geométrica \cite{dormans2010, summerville2018procedural}. Os vértices e arestas de um grafo de progressão definem quais elementos existem, quais funções desempenham e como se conectam logicamente, mas não determinam, por si sós, coordenadas cartesianas $(x, y, z)$, dimensões em metros, formatos poligonais de salas ou comprimentos e curvaturas de corredores \cite{dormans2010, inproceedings, linden2014dungeons, summerville2018procedural}.

Essa abstração implica que a correspondência entre um grafo lógico e um \textit{layout} tridimensional é de um-para-muitos: um único grafo de progressão admite múltiplas realizações espaciais geometricamente distintas no espaço 3D, desde que as relações de adjacência e conectividade especificadas pelas arestas sejam preservadas \cite{dormans2010, inproceedings}. Ao mesmo tempo, nem todo grafo abstrato pode ser mapeado diretamente em uma grade ortogonal ou em um pavimento plano sem adaptações geométricas --- por exemplo, um grafo não planar ou com vértices cujo grau exceda o número físico de portas disponíveis em um módulo de sala requer a introdução de nós intermediários (como corredores ramificados ou múltiplos andares) durante a etapa de resolução de \textit{layout} \cite{inproceedings, linden2014dungeons, xu2025databasedrivenframework3dlevel}.

% ---------------------------------------------------------------------------
\subsection{Representação de \textit{Layouts} Espaciais 3D}
\label{sec:layout_3d}
% ---------------------------------------------------------------------------

\subsubsection{Representação Modular de Ambientes}
\label{sec:representacao_modular}

Para materializar um ambiente virtual a partir de uma especificação lógica, a literatura de PCG e a indústria de jogos empregam extensivamente a representação espacial modular \cite{linden2014dungeons, shaker2016constructive, xu2025databasedrivenframework3dlevel}. Nesse paradigma, o cenário tridimensional não é esculpido como uma malha única monolítica, mas construído pela composição de unidades geométricas discretas e reutilizáveis \cite{linden2014dungeons, shaker2016constructive, xu2025databasedrivenframework3dlevel}.

As unidades modulares variam em granularidade conforme a arquitetura do gerador \cite{linden2014dungeons, shaker2016constructive, vieira2024metodologia}:
\begin{itemize}
    \item \textbf{Grades e ladrilhos (\textit{grids}, \textit{tiles} e \textit{voxels}):} subdivisões regulares do espaço em células quadradas (2D) ou cúbicas (3D) de tamanho uniforme, nas quais cada célula recebe uma peça elementar (piso, parede, canto ou espaço vazio) \cite{pcgbook, shaker2016constructive, summerville2018procedural};
    \item \textbf{Blocos e segmentos (\textit{chunks}):} agrupamentos pré-autorados de elementos geométricos e interativos que formam seções parciais de um nível \cite{linden2014dungeons, shaker2016constructive};
    \item \textbf{Salas pré-construídas (\textit{room templates}) e \textit{prefabs} 3D:} módulos arquitetônicos completos (como salas temáticas de dimensões conhecidas, corredores e escadarias) ou ativos pré-fabricados (\textit{prefabs}) de objetos internos \cite{linden2014dungeons, xu2025databasedrivenframework3dlevel, regmi2025system}. Um \textit{prefab} encapsula em um único pacote reutilizável a malha poligonal tridimensional (\textit{mesh}), texturas, volumes de colisão física e metadados semânticos \cite{xu2025databasedrivenframework3dlevel, regmi2025system}.
\end{itemize}

Para viabilizar o acoplamento preciso entre módulos adjacentes (como unir uma sala a um corredor ou duas salas entre si), cada componente modular define pontos de conexão explícitos --- denominados conectores, âncoras, portas ou portais (\textit{connectors} / \textit{doorways} / \textit{exits}) \cite{inproceedings, linden2014dungeons, shaker2016constructive, sun2024roomplanner}. Esses conectores especificam a posição local da abertura na fronteira do módulo, suas dimensões de passagem e a direção ou orientação espacial para a qual a conexão se abre \cite{inproceedings, shaker2016constructive, sun2024roomplanner, xu2025databasedrivenframework3dlevel}.

\subsubsection{Relações e Transformações Espaciais}
\label{sec:transformacoes_espaciais}

A instanciação de componentes modulares em um espaço tridimensional contínuo fundamenta-se na geometria analítica e na álgebra linear, operando sobre dois referenciais distintos: o sistema de coordenadas local (intrínseco a cada módulo ou objeto, geralmente centrado em seu ponto de pivô ou origem geométrica) e o sistema de coordenadas global do mundo (\textit{world space}) \cite{sun2024roomplanner, xu2025constraint, ocal2024sceneteller}.

A configuração espacial (ou \textit{pose}) de cada componente $i$ no espaço global $\mathbb{R}^3$ é determinada por transformações afins que combinam três parâmetros fundamentais \cite{sun2024roomplanner, xu2025constraint}: (i) \textit{translação} (posição $\mathbf{p}_i = (x_i, y_i, z_i) \in \mathbb{R}^3$), que desloca a origem local do módulo para sua coordenada no mundo; (ii) \textit{rotação} (orientação $\mathbf{R}_i$, expressa por ângulos de Euler, matrizes de rotação ortogonais ou quatérnios), frequentemente parametrizada em \textit{layouts} arquitetônicos pelo ângulo de guinada $\theta_i$ em torno do eixo vertical para alinhar paredes e portas às direções cardeais \cite{sun2024roomplanner, xu2025databasedrivenframework3dlevel}; e (iii) \textit{escala} ou \textit{dimensão} ($\mathbf{s}_i = (w_i, l_i, h_i) \in \mathbb{R}^3$), que define a largura, o comprimento e a altura física ocupados pelo elemento \cite{sun2024roomplanner, xu2025constraint}.

Para simplificar o cálculo de ocupação espacial e de distâncias entre malhas poligonais arbitrárias, a computação gráfica e os sistemas de geração de \textit{layout} utilizam volumes delimitadores (\textit{bounding volumes}) tridimensionais \cite{sun2024roomplanner, xu2025constraint, xu2025databasedrivenframework3dlevel}. A representação padrão consiste na Caixa Delimitadora Alinhada aos Eixos (\textit{Axis-Aligned Bounding Box} --- AABB), definida pelos intervalos mínimos e máximos $[x_{\min}, x_{\max}] \times [y_{\min}, y_{\max}] \times [z_{\min}, z_{\max}]$ de cada objeto ao longo dos eixos coordenados, ajustada conforme a orientação espacial do elemento na cena \cite{sun2024roomplanner, xu2025constraint}.

\subsubsection{Restrições Geométricas}
\label{sec:restricoes_geometricas}

Para que um conjunto de módulos e objetos transformados no espaço $\mathbb{R}^3$ constitua um \textit{layout} arquitetonicamente íntegro, diversas restrições geométricas precisam ser satisfeitas simultaneamente \cite{linden2014dungeons, sun2024roomplanner, xu2025constraint, xu2025databasedrivenframework3dlevel}:
\begin{itemize}
    \item \textbf{Restrição de não sobreposição (\textit{non-overlap}):} dois módulos de salas, corredores ou objetos sólidos distintos não devem ocupar simultaneamente o mesmo volume no espaço contínuo; formalmente, a interseção volumétrica entre seus interiores (frequentemente quantificada pela métrica de Interseção sobre União tridimensional, $\text{IoU}_{3D}$, ou pelo volume de interseção $V_{\text{overlap}}$ de suas caixas delimitadoras) deve ser nula \cite{sun2024roomplanner, xu2025constraint};
    \item \textbf{Restrição de contenção e limites (\textit{boundary constraints}):} elementos internos alocados a uma sala devem ter seus volumes delimitadores integralmente contidos dentro do perímetro útil definido pelas paredes, piso e teto daquela sala, sem projetar-se para o exterior \cite{sun2024roomplanner, xu2025constraint};
    \item \textbf{Compatibilidade e alinhamento de conexões:} para que duas salas adjacentes ou uma sala e um corredor se conectem fisicamente, seus respectivos portais/portas devem coincidir em posição e dimensão, apresentando orientações compatíveis ao longo da fronteira compartilhada, sem que paredes sólidas obstruam a abertura \cite{inproceedings, shaker2016constructive, sun2024roomplanner, xu2025databasedrivenframework3dlevel};
    \item \textbf{Espaço livre e desobstrução (\textit{clearance}):} a disposição de objetos internos deve preservar corredores de circulação livres ao redor das entradas, saídas e pontos de interação, evitando o congestionamento excessivo (\textit{clustering}) que impeça a passagem física de agentes \cite{sun2024roomplanner, xu2025constraint, xu2025databasedrivenframework3dlevel}.
\end{itemize}

% ---------------------------------------------------------------------------
\subsection{Arquiteturas Híbridas e Geração Baseada em Restrições}
\label{sec:arquiteturas_hibridas}
% ---------------------------------------------------------------------------

\subsubsection{Sistemas Híbridos e Neuro-Simbólicos}
\label{sec:sistemas_hibridos}

A constatação de que nenhum paradigma isolado de PCG é isento de limitações --- métodos construtivos e baseados em restrições oferecem garantias estruturais, mas exigem codificação manual complexa de regras criativas, ao passo que modelos de aprendizado profundo e LLMs oferecem alta expressividade semântica, mas carecem de garantias espaciais e lógicas exatas --- impulsionou o desenvolvimento de arquiteturas híbridas e composicionais \cite{linden2014dungeons, Gallotta_2024, maleki2024procedural}. Na geração procedural composicional, métodos distintos são encadeados ou aninhados em um mesmo fluxo, de modo que as potencialidades de uma técnica compensem as fragilidades da outra \cite{linden2014dungeons, maleki2024procedural}.

Nas pesquisas contemporâneas que envolvem modelos fundacionais, essa integração assume frequentemente a forma de \textit{arquiteturas híbridas neuro-simbólicas} \cite{Gallotta_2024, hu2024scenecraft, chen2025narrative, xu2025constraint}. Nesses sistemas, o componente neural ou probabilístico (como um LLM) opera nas camadas de abstração semântica e simbólica --- interpretando intenções em linguagem natural, estruturando narrativas, definindo grafos de cena ou especificando restrições qualitativas ---, enquanto componentes simbólicos, formais ou algorítmicos (como solucionadores de restrições, algoritmos de busca em grafos, otimizadores matemáticos ou rotinas geométricas em motores gráficos) encarregam-se da resolução métrica, da instanciação espacial e da verificação rigorosa de regras \cite{hu2024scenecraft, tang2025unrealllm, chen2025narrative, xu2025constraint, xu2025databasedrivenframework3dlevel, her2025zeroshot3dmapgeneration}.

De maneira correlata, a literatura destaca os sistemas de \textit{iniciativa mista} (\textit{Mixed-Initiative PCG} ou cocriação humano-IA), exemplificados por ferramentas como \textit{Tanagra} \cite{tanagra} e \textit{Sentient Sketchbook} \cite{linden2014dungeons}, nos quais projetistas humanos e algoritmos generativos colaboram iterativamente sobre uma representação compartilhada \cite{tanagra, lai2020towards, walton2020evaluating, wang2026designerlyloop}. Nesses ambientes, o usuário pode especificar diretrizes de alto nível, curar cadeias de raciocínio intermediárias ou fixar partes da geometria, cabendo ao sistema procedural preencher, adaptar e validar o restante do nível mantendo o respeito ao controle do projetista e a fluidez do processo criativo \cite{tanagra, lai2020towards, bandeali2026generative, wang2026designerlyloop}.

\subsubsection{Geração Baseada em Regras e Restrições}
\label{sec:geracao_restricoes}

A geração controlada em sistemas formais e híbridos apoia-se teoricamente na modelagem por regras e na Satisfação de Restrições \cite{tanagra, linden2014dungeons, pcgbook, xu2025constraint}. Na formulação de um Problema de Satisfação de Restrições (\textit{Constraint Satisfaction Problem} --- CSP), o artefato a ser gerado é modelado por meio de um conjunto de variáveis de decisão (por exemplo, as coordenadas de posição, dimensões ou orientações atribuídas a cada componente do nível), de domínios numéricos ou discretos que delimitam os valores possíveis para cada variável (como as dimensões máximas da área de construção) e de um conjunto de restrições algébricas ou lógicas que especificam quais combinações simultâneas de valores são admissíveis \cite{tanagra, linden2014dungeons}.

Na formulação de problemas de \textit{layout} espacial e de design de fases, a literatura diferencia duas classes de restrições \cite{xu2025constraint, xu2025databasedrivenframework3dlevel}:
\begin{itemize}
    \item \textbf{Restrições rígidas (\textit{hard constraints}):} condições invioláveis cujo cumprimento integral é obrigatório para que a solução seja considerada admissível, tais como a ausência de sobreposição entre salas, a contenção de objetos dentro dos limites do mapa e a precedência topológica de uma chave em relação à porta que ela destranca \cite{tanagra, sun2024roomplanner, xu2025constraint};
    \item \textbf{Restrições flexíveis (\textit{soft constraints}):} preferências estilísticas, funcionais ou de qualidade espacial que não invalidam estruturalmente a fase caso sofram pequenos desvios, mas cujo atendimento deve ser maximizado, tais como distâncias ideais de proximidade entre objetos correlatos, grau de agrupamento (\textit{clustering}) ou distribuição equilibrada de espaços vazios (\textit{sparsity}) \cite{xu2025constraint, xu2025databasedrivenframework3dlevel}.
\end{itemize}

Quando o problema combina restrições rígidas e flexíveis no espaço contínuo, ele é frequentemente modelado como um Problema de Otimização de \textit{Layout} de Instalações (\textit{Facility Layout Problem} --- FLP), no qual uma função objetivo agrega termos de penalidade ponderados para cada violação de restrição, buscando uma configuração espacial que minimize o custo total \cite{xu2025constraint, xu2025databasedrivenframework3dlevel}.

\subsubsection{Validação e Reparo de Conteúdo Gerado}
\label{sec:validacao_reparo}

Como geradores estocásticos, heurísticos ou baseados em aprendizado de máquina podem produzir instâncias candidatas que violem uma ou mais restrições do domínio, a literatura de PCG estabelece diferentes estratégias teóricas para a garantia de qualidade e correção do conteúdo gerado \cite{pcgbook, summerville2018procedural, maleki2024procedural, xu2025databasedrivenframework3dlevel}:
\begin{enumerate}
    \item \textbf{Geração e teste com filtragem (\textit{generate-and-test} / rejeição):} o gerador produz um artefato completo ou parcial e o submete a rotinas determinísticas de verificação; caso alguma restrição rígida seja violada, a instância é integralmente rejeitada e uma nova geração é solicitada \cite{pcgbook, maleki2024procedural};
    \item \textbf{Busca construtiva com retrocesso (\textit{backtracking}):} durante a montagem passo a passo do nível (como na alocação sequencial de salas ou padrões geométricos), se o algoritmo detecta que a decisão atual inviabiliza a satisfação das restrições restantes, ele desfaz a última atribuição, poda o ramo inválido do espaço de busca e tenta uma configuração alternativa ou relaxa restrições locais \cite{tanagra, linden2014dungeons};
    \item \textbf{Reparo pós-geração (\textit{post-generation repair}):} em vez de descartar um nível quase válido, algoritmos especializados aplicam transformações locais corretivas sobre o artefato defeituoso, tais como o deslocamento incremental de objetos em colisão, a abertura algorítmica de passagens entre regiões desconectadas ou a remoção seletiva de elementos que obstruam portas e rotas de navegação \cite{summerville2018procedural, sun2024roomplanner, xu2025databasedrivenframework3dlevel};
    \item \textbf{Refinamento iterativo guiado por diagnóstico (\textit{feedback loop}):} quando o sistema conta com um planejador generativo ou módulo de otimização, os erros detectados pelos validadores formais (sintáticos, lógicos ou espaciais) são estruturados como sinais de diagnóstico para orientar uma nova iteração de ajuste da representação \cite{hu2024scenecraft, nasir2024word2worldgeneratingstoriesworlds, xu2025constraint, her2025zeroshot3dmapgeneration}.
\end{enumerate}

% ---------------------------------------------------------------------------
\subsection{Validade, Navegabilidade e Solvabilidade de Fases}
\label{sec:validade_navegabilidade_solvabilidade}
% ---------------------------------------------------------------------------

A avaliação formal de uma fase gerada proceduralmente exige decompor o conceito geral de ``qualidade funcional'' em propriedades distintas e verificáveis. Embora frequentemente relacionados na prática, os conceitos de validade lógica, validade geométrica, navegabilidade, busca de caminhos (\textit{pathfinding}) e solvabilidade operam em camadas diferentes de abstração e não são sinônimos \cite{dormans2010, inproceedings, linden2014dungeons, sun2024roomplanner, xu2025constraint}.

\subsubsection{Validade Lógica e Estrutural}
\label{sec:validade_logica}

A validade lógica e estrutural refere-se à consistência interna da fase em seu nível de representação abstrata ou topológica, independentemente de sua realização geométrica no espaço 3D \cite{dormans2010, inproceedings, linden2014dungeons}. Uma estrutura de nível é logicamente válida quando satisfaz as invariantes sintáticas e semânticas do domínio do jogo: a existência bem definida de um estado ou região inicial e de pelo menos uma condição ou região de objetivo final; a conectividade topológica do grafo subjacente (ausência de componentes obrigatórias isoladas); e a consistência causal das relações de dependência e pré-condições \cite{dormans2010, inproceedings, xu2025databasedrivenframework3dlevel}.

Violações de validade lógica ocorrem, por exemplo, quando o grafo de progressão contém dependências circulares insolúveis (\textit{deadlocks} lógicos) --- tais como posicionar uma chave exclusivamente no interior de uma sala cuja única porta de acesso requer aquela mesma chave para ser aberta ---, ou quando uma ramificação opcional absorve recursos indispensáveis sem permitir retorno ao caminho principal \cite{dormans2010, linden2014dungeons, xu2025databasedrivenframework3dlevel}.

\subsubsection{Validade Geométrica}
\label{sec:validade_geometrica}

A validade geométrica diz respeito à conformidade física e métrica da fase quando instanciada no espaço euclidiano bidimensional ou tridimensional \cite{sun2024roomplanner, xu2025constraint, feng2026repurposing3dgenerativemodel}. Um \textit{layout} 3D detém validade geométrica quando todas as salas, corredores e objetos internos estão posicionados dentro dos limites espaciais permitidos, sem interpenetrações ou sobreposições volumétricas inválidas (\textit{clipping}) entre malhas sólidas, e com alinhamento geométrico correto nas interfaces de conexão (como portais e batentes de portas entre módulos adjacentes) \cite{sun2024roomplanner, xu2025constraint, xu2025databasedrivenframework3dlevel}.

Na computação gráfica e em motores de jogos, a verificação da validade geométrica apoia-se em estruturas de dados e consultas espaciais, tais como testes de interseção entre caixas delimitadoras (como AABBs), árvores de particionamento espacial (\textit{quadtrees}, \textit{octrees} e árvores BSP) e lançamento de raios ou volumes (\textit{raycasting} e \textit{sphere casting}) \cite{shaker2016constructive, sun2024roomplanner, xu2025constraint, xu2025databasedrivenframework3dlevel}. Cabe destacar, sob o rigor conceitual, que estruturas como caixas delimitadoras, árvores espaciais e \textit{raycasting} não previnem nem corrigem, por si sós, sobreposições de malhas: elas constituem mecanismos matemáticos de consulta e detecção de colisão ou visibilidade, os quais fornecem os dados necessários para que algoritmos de posicionamento, otimização ou reparo atuem evitando ou eliminando configurações inválidas \cite{shaker2016constructive, xu2025constraint}.

\subsubsection{Navegabilidade}
\label{sec:navegabilidade}

A navegabilidade (\textit{navigability} ou \textit{traversability}) caracteriza a propriedade de um ambiente espacial de oferecer superfícies contínuas e fisicamente transponíveis para a locomoção de um agente virtual ou jogador, levando em conta suas restrições cinemáticas e dimensionais (como raio de colisão do corpo, altura vertical livre, inclinação máxima de rampas e altura máxima de degraus superáveis) \cite{sun2024roomplanner, xu2025constraint, xu2025databasedrivenframework3dlevel}. Mesmo que um cenário seja geometricamente válido (sem colisões entre paredes ou objetos), ele pode carecer de navegabilidade caso uma passagem seja estreita demais para o volume do agente, caso haja desníveis verticais sem escadas ou rampas compatíveis, ou caso móveis dispostos perto de uma porta obstruam o raio de passagem \cite{sun2024roomplanner, xu2025databasedrivenframework3dlevel}.

Para representar computacionalmente o subespaço navegável de um ambiente 3D, utilizam-se estruturas como grades de ocupação bidimensionais (\textit{occupancy grids}) \cite{sun2024roomplanner} e, de forma predominante em ambientes tridimensionais contínuos, as Malhas de Navegação (\textit{Navigation Meshes} --- \textit{NavMesh}) \cite{xu2025constraint, xu2025databasedrivenframework3dlevel, regmi2025system}. Uma \textit{NavMesh} consiste em uma representação geométrica das superfícies caminháveis (\textit{walkable surfaces} ou \textit{walkable areas}) do cenário tridimensional, delimitando as regiões livres de obstáculos estáticos sobre as quais agentes virtuais e jogadores podem se deslocar e executar algoritmos de busca de caminho com desvio de colisões \cite{xu2025constraint, xu2025databasedrivenframework3dlevel, regmi2025system}.

\subsubsection{\textit{Pathfinding}}
\label{sec:pathfinding}

A busca de caminhos (\textit{pathfinding}) consiste no processo algorítmico de encontrar uma rota contínua e livre de obstáculos entre uma posição de origem e uma posição de destino sobre uma representação discreta do espaço navegável, tal como uma grade de ocupação ou uma \textit{NavMesh} \cite{sun2024roomplanner, nasir2024word2worldgeneratingstoriesworlds, xu2025databasedrivenframework3dlevel}.

Enquanto algoritmos de busca em largura (\textit{Breadth-First Search} --- BFS) e de preenchimento por inundação (\textit{flood-fill}) exploram uniformemente os nós adjacentes para verificar a conectividade e identificar regiões conexas \cite{xu2025databasedrivenframework3dlevel}, a determinação de rotas em espaços métricos emprega amplamente algoritmos de busca informada por heurística, com destaque para o algoritmo $A^*$ (\textit{A-star}) \cite{sun2024roomplanner, nasir2024word2worldgeneratingstoriesworlds, xu2025databasedrivenframework3dlevel}. O algoritmo $A^*$ explora o espaço navegável combinando o custo acumulado desde o ponto de partida (como uma entrada ou porta) com uma estimativa heurística de distância até o destino para determinar se existe um trajeto viável e desobstruído entre dois pontos de interesse \cite{sun2024roomplanner, nasir2024word2worldgeneratingstoriesworlds, xu2025databasedrivenframework3dlevel}. No contexto de validação de níveis, a execução bem-sucedida de um algoritmo de \textit{pathfinding} sobre a grade de ocupação ou \textit{NavMesh} comprova a existência de conectividade física navegável entre as regiões e objetos do espaço \cite{sun2024roomplanner, xu2025databasedrivenframework3dlevel}.

\subsubsection{Solvabilidade de uma Fase}
\label{sec:solvabilidade}

A solvabilidade (\textit{solvability}, \textit{winnability} ou \textit{completability}) constitui o critério funcional mais abrangente de uma fase de jogo: ela atesta que existe ao menos uma sequência finita e executável de ações e deslocamentos que permite ao jogador transitar do estado inicial da fase até a condição de conclusão ou vitória, respeitando simultaneamente a geometria física do espaço e todas as regras, estados dinâmicos, pré-condições e restrições de mecânicas do jogo \cite{tanagra, linden2014dungeons, summerville2018procedural, nasir2024word2worldgeneratingstoriesworlds, xu2025constraint}.

A distinção teórica rigorosa entre validade geométrica, navegabilidade/\textit{pathfinding} e solvabilidade é indispensável porque o sucesso em uma dessas dimensões não implica automaticamente o sucesso nas demais \cite{dormans2010, linden2014dungeons, summerville2018procedural, xu2025constraint}:
\begin{itemize}
    \item \textbf{Validade geométrica sem navegabilidade:} duas salas podem estar posicionadas sem qualquer sobreposição de malhas (\textit{clipping} zero), mas separadas por uma falha na malha de piso ou por um obstáculo que impeça a formação de uma \textit{NavMesh} contínua entre elas \cite{sun2024roomplanner, xu2025databasedrivenframework3dlevel};
    \item \textbf{Navegabilidade e \textit{pathfinding} sem solvabilidade:} a existência de um caminho geométrico encontrado pelo algoritmo $A^*$ entre a sala inicial e a sala final sobre o piso estático não garante, por si só, que a fase seja solucionável. Se o caminho atravessar uma porta trancada cuja chave foi alocada após a própria porta, ou se os recursos exigidos para superar os desafios intermediários forem insuficientes, haverá conectividade espacial na malha geométrica, mas a fase será logicamente insolúvel no espaço de estados do jogo \cite{dormans2010, inproceedings, xu2025constraint};
    \item \textbf{Validade lógica abstrata sem solvabilidade espacial:} inversamente, um grafo de progressão pode ter todas as suas dependências de chaves e portas perfeitamente ordenadas no nível simbólico, mas resultar em uma fase insolúvel caso a etapa subsequente de instanciação geométrica 3D posicione um elemento que bloqueie fisicamente o corredor de acesso à sala da chave \cite{linden2014dungeons, xu2025databasedrivenframework3dlevel}.
\end{itemize}

Por conseguinte, a verificação plena da solvabilidade em fases com progressão condicional requer que a análise de alcançabilidade considere tanto o espaço de estados lógicos (transições condicionadas à posse de itens, chaves e cumprimento de objetivos) quanto a confirmação de que cada transição admitida no grafo lógico corresponde a uma rota fisicamente navegável no ambiente tridimensional gerado \cite{dormans2010, linden2014dungeons, xu2025constraint, xu2025databasedrivenframework3dlevel}.

% ---------------------------------------------------------------------------
\section{Fundamentação Teórica}
\label{sec:fundamentacao}
% ---------------------------------------------------------------------------

Este capítulo apresenta os fundamentos teóricos que sustentam o estudo da geração automática de fases tridimensionais para jogos digitais. A exposição organiza-se de modo progressivo: parte-se da definição de fases, de \textit{Level Design} e da distinção conceitual entre organização lógica e organização espacial; em seguida, abordam-se os princípios da Geração Procedural de Conteúdo (PCG), sua vertente baseada em aprendizado de máquina (PCGML) e a caracterização dos Modelos de Linguagem de Grande Porte (LLMs), incluindo suas limitações de raciocínio e ancoragem espacial. Por fim, discutem-se a modelagem lógica de fases por meio da Teoria dos Grafos, a representação modular e geométrica de \textit{layouts} tridimensionais, as arquiteturas híbridas guiadas por restrições e os critérios formais de validade, navegabilidade e solvabilidade.

% ---------------------------------------------------------------------------
\subsection{Fases de Jogos, \textit{Level Design} e \textit{Layouts} 3D}
\label{sec:level_design}
% ---------------------------------------------------------------------------

\subsubsection{Fases, Níveis e \textit{Level Design}}
\label{sec:fases_niveis}

No contexto dos jogos digitais, uma fase --- também denominada nível (\textit{level}) ou mapa --- constitui a unidade espacial, temporal e interativa na qual as regras e mecânicas do sistema de regras do jogo se manifestam para o jogador \cite{tanagra, pcgbook}. Mais do que um cenário estático, uma fase opera como um contêiner estruturado para a experiência de jogo (\textit{gameplay}), delimitando o território de exploração, os recursos disponíveis, os obstáculos a serem transpostos e as condições de vitória ou transição \cite{tanagra, linden2014dungeons}. Em gêneros como ação e aventura, jogos de plataforma e \textit{Role-Playing Games} (RPGs), a qualidade da experiência interativa depende diretamente de como esses espaços são construídos e articulados \cite{dormans2010, linden2014dungeons}.

A disciplina responsável pelo planejamento, concepção e estruturação de fases é o \textit{Level Design} \cite{dormans2010, tanagra, totten2014architectural}. Seus objetivos fundamentais abrangem a condução da atenção e do deslocamento do jogador no espaço tridimensional, a introdução gradual de mecânicas, o balanceamento entre desafio e recompensa e a regulação do ritmo da experiência (\textit{pacing}) \cite{dormans2010, linden2013designing, linden2014dungeons, totten2014architectural}. Em fases do tipo masmorra (\textit{dungeon}), por exemplo, o \textit{Level Design} busca conciliar a liberdade de exploração espacial com um controle rigoroso sobre a progressão de desafios e quebra-cabeças (\textit{puzzles}), assegurando que o jogador não encontre obstáculos intransponíveis para o conjunto de habilidades ou recursos que detém em determinado instante da partida \cite{linden2014dungeons, shaker2016constructive}.

\subsubsection{\textit{Layout} de uma Fase}
\label{sec:layout_fase}

O \textit{layout} de uma fase refere-se ao arranjo estrutural e espacial dos elementos que compõem o ambiente jogável \cite{dormans2010, linden2014dungeons}. Em ambientes internos e masmorras, o \textit{layout} é classicamente decomposto em unidades arquitetônicas discretas, destacando-se as salas (ou câmaras), os corredores (ou vias de transição), os conectores (portas, portais, escadas e elevadores) e as regiões funcionais \cite{linden2014dungeons, shaker2016constructive, xu2025databasedrivenframework3dlevel}. As salas atuam tipicamente como arenas de interação concentrada --- abrigando combates, enigmas, recursos ou eventos narrativos ---, ao passo que os corredores estabelecem as ligações físicas entre essas áreas, modulando o tempo de travessia e a visibilidade entre setores adjacentes \cite{linden2013designing, shaker2016constructive}.

Além da geometria estrutural que delimita o espaço navegável (pisos, paredes, tetos e desníveis), o \textit{layout} compreende a disposição interna de elementos funcionais e estéticos, frequentemente denominados instalações ou objetos de cena (\textit{facilities} ou \textit{props}), bem como entidades dinâmicas, tais como personagens não jogáveis (NPCs), inimigos, armadilhas e itens coletáveis \cite{linden2014dungeons, xu2025constraint}. Em ambientes tridimensionais (3D), a especificação de um \textit{layout} envolve não apenas a distribuição planar em dois eixos, mas também a verticalidade, a sobreposição de múltiplos pavimentos, as relações de linha de visão (\textit{line of sight}) e as dimensões volumétricas exatas de cada componente no espaço contínuo \cite{xu2025constraint, xu2025databasedrivenframework3dlevel}.

\subsubsection{Organização Lógica e Organização Espacial}
\label{sec:organizacao_logica_espacial}

Um princípio teórico central na literatura de \textit{Level Design} consiste na distinção analítica entre duas estruturas complementares que coexistem em uma fase: a \textit{missão} (organização lógica) e o \textit{espaço} (organização espacial) \cite{dormans2010, linden2014dungeons}. A missão corresponde à sequência lógica ou parcialmente ordenada de tarefas, ações e objetivos que o jogador deve cumprir para concluir o nível, tais como derrotar um guardião, obter uma chave específica, desativar um bloqueio e alcançar a saída \cite{dormans2010, linden2013designing}. O espaço, por sua vez, refere-se à realização física e geométrica do ambiente virtual onde essas tarefas estão situadas, englobando as dimensões das salas, as distâncias métricas e a configuração arquitetônica das conexões \cite{dormans2010, linden2014dungeons}.

Embora o espaço físico acomode a missão e esta seja mapeada sobre aquele, as duas estruturas são conceitualmente independentes e regidas por princípios distintos \cite{dormans2010}. Uma missão estritamente linear --- na qual os objetivos precisam ser cumpridos em uma ordem única e fixa --- pode ser distribuída sobre um espaço físico não linear, como um arranjo em estrela (\textit{hub-and-spoke}) ou em anel, exigindo que o jogador retorne a áreas previamente visitadas à medida que desbloqueia novos acessos \cite{dormans2010}. Inversamente, uma missão não linear, dotada de múltiplos desafios paralelos ou caminhos alternativos, pode ser acomodada em uma geometria espacial simples \cite{dormans2010, linden2013designing}. Essa dissociação teórica evidencia que a coerência da progressão de um nível reside primariamente em sua estrutura lógica de dependências, enquanto sua viabilidade física e navegacional depende da correta realização geométrica no espaço tridimensional \cite{dormans2010, linden2013designing, linden2014dungeons}.

% ---------------------------------------------------------------------------
\subsection{Geração Procedural de Conteúdo}
\label{sec:pcg}
% ---------------------------------------------------------------------------

\subsubsection{Conceito e Objetivos da PCG}
\label{sec:conceito_pcg}

A Geração Procedural de Conteúdo (PCG --- \textit{Procedural Content Generation}) é definida na literatura como a criação algorítmica de conteúdos para jogos digitais de maneira autônoma ou com intervenção humana limitada e indireta \cite{pcgbook, summerville2018procedural, maleki2024procedural}. O escopo do termo ``conteúdo'' abrange praticamente todos os artefatos que integram um jogo, incluindo níveis, mapas, terrenos, vegetação, texturas, efeitos sonoros, narrativas, missões (\textit{quests}) e conjuntos de regras, distinguindo-se do comportamento em tempo de execução dos agentes de inteligência artificial que controlam NPCs \cite{pcgbook, summerville2018procedural, maleki2024procedural}.

Historicamente introduzida em títulos pioneiros da década de 1980, como \textit{Rogue} (1980) e \textit{Elite} (1984), para contornar severas restrições de memória e armazenamento por meio da síntese algorítmica em tempo de execução \cite{pcgbook, summerville2018procedural, araujo2018geracao}, a PCG consolidou-se tanto na indústria quanto na pesquisa acadêmica em função de múltiplos objetivos \cite{linden2014dungeons, pcgbook}. Dentre as principais motivações para o seu emprego destacam-se: (i) a automação e a redução do tempo e do custo financeiro associados à autoria manual de ativos em larga escala; (ii) o aumento da variabilidade e da diversidade dos ambientes gerados, ampliando o fator de rejogabilidade (\textit{replayability}); (iii) o suporte à adaptabilidade centrada no jogador, permitindo ajustar dinamicamente os desafios às habilidades observadas; e (iv) a exploração de novas estéticas e mecânicas emergentes \cite{linden2014dungeons, pcgbook, summerville2018procedural}.

A avaliação e o projeto de geradores procedurais apoiam-se em propriedades técnicas fundamentais, notadamente a reprodutibilidade, a expressividade e a controlabilidade \cite{tanagra, smith2010analyzing, linden2014dungeons, maleki2024procedural}. A reprodutibilidade decorre frequentemente do uso de Geradores de Números Pseudoaleatórios (PRNGs), nos quais a especificação de um valor inicial --- denominado semente (\textit{seed}) --- permite reconstruir de maneira determinística a mesma sequência de decisões algorítmicas \cite{maleki2024procedural}. A expressividade, caracterizada pela análise do espaço expressivo (\textit{expressive range}), quantifica e mapeia a variedade, a distribuição e as propriedades estilísticas ou estruturais dos artefatos que o gerador é capaz de produzir no espaço de design \cite{smith2010analyzing, linden2013designing}. A controlabilidade, por sua vez, reflete a capacidade do projetista ou do sistema de direcionar intencionalmente as características do conteúdo produzido por meio de parâmetros, regras ou restrições \cite{linden2014dungeons, shaker2016constructive}. Um desafio central no projeto de geradores consiste em equilibrar a tensão entre controle e diversidade: restrições excessivamente rígidas podem limitar a variabilidade dos resultados, ao passo que a ausência de controle eleva o risco de geração de artefatos inválidos ou incoerentes \cite{smith2010analyzing, linden2014dungeons}.

\subsubsection{Principais Abordagens de PCG}
\label{sec:abordagens_pcg}

Uma classificação recorrente na literatura de PCG organiza os métodos generativos em três grandes paradigmas operacionais: métodos construtivos, métodos baseados em busca (\textit{search-based}) e métodos baseados em restrições ou solucionadores (\textit{solver-based} / \textit{constraint-based}) \cite{togelius2011search, pcgbook, summerville2018procedural, maleki2024procedural}.

Os métodos construtivos caracterizam-se por gerar o conteúdo em uma única passagem sequencial de passos algorítmicos, com tempo de execução previsível e sem realizar ciclos iterativos de avaliação e regeneração sobre o artefato produzido \cite{shaker2016constructive, summerville2018procedural}. Exemplos clássicos incluem funções de ruído coerente (como o ruído de Perlin para síntese de terrenos e texturas), autômatos celulares, algoritmos de particionamento espacial e gramáticas generativas \cite{linden2014dungeons, shaker2016constructive, maleki2024procedural}. Importa ressaltar que procedimentos construtivos podem fornecer garantias estruturais formais sobre o conteúdo gerado (como ausência de sobreposição ou conectividade garantida) quando suas regras forem rigorosamente formuladas para preservar tais invariantes em cada etapa da construção \cite{shaker2016constructive}. Por outro lado, caso dependam de decisões puramente estocásticas desprovidas de verificação antecipada, o determinismo algorítmico assegura apenas a repetibilidade da saída a partir de uma semente fixa, não prevenindo falhas estruturais ou becos sem saída \cite{linden2014dungeons, shaker2016constructive}.

Os métodos baseados em busca (\textit{search-based PCG}) enquadram-se como um caso particular do paradigma gerar-e-testar (\textit{generate-and-test}) guiado por otimização \cite{togelius2011search, pcgbook, maleki2024procedural}. Em vez de aceitar imediatamente o artefato sintetizado ou simplesmente descartá-lo de forma binária, algoritmos baseados em busca atribuem a cada instância candidata um valor numérico contínuo por meio de uma função de avaliação ou aptidão (\textit{fitness function}), utilizando esse escore para guiar iterativamente a modificação e a seleção de novas soluções no espaço de busca \cite{togelius2011search, linden2014dungeons, araujo2018geracao, maleki2024procedural}. Essa categoria emprega amplamente técnicas de computação evolutiva (como Algoritmos Genéticos) \cite{togelius2011search, linden2014dungeons, maleki2024procedural}, além de meta-heurísticas como têmpera simulada (\textit{Simulated Annealing}), otimização por enxame de partículas (PSO) e Busca em Árvore de Monte Carlo (MCTS) \cite{maleki2024procedural, xu2025constraint}.

Os métodos baseados em restrições (\textit{constraint-based} ou \textit{solver-based PCG}), por sua vez, especificam declarativamente as propriedades e relações obrigatórias que o artefato final deve satisfazer, delegando a construção a solucionadores matemáticos ou lógicos (\textit{constraint solvers}) \cite{tanagra, smith2011answer, linden2014dungeons, pcgbook, summerville2018procedural}. Destacam-se nesse paradigma a Programação por Conjuntos de Respostas (\textit{Answer Set Programming} --- ASP), que formaliza regras e restrições de design como predicados lógicos satisfatíveis \cite{smith2011answer}, solucionadores de satisfatibilidade módulo teorias (SMT) e solucionadores de Problemas de Satisfação de Restrições (CSP) \cite{tanagra, chen2025narrative}. Integram também essa vertente os algoritmos baseados em propagação local de restrições de adjacência inspirados em mecânica estatística, com destaque para o \textit{WaveFunctionCollapse} (WFC) \cite{karth2017wavefunctioncollapse, maleki2024procedural}. Nessa abordagem, o foco desloca-se da definição de passos procedurais imperativos para a formulação matemática e lógica do espaço de design admissível \cite{tanagra, smith2011answer, linden2014dungeons}.

\subsubsection{Geração Procedural de Fases}
\label{sec:pcg_fases}

A aplicação de PCG à criação de fases e níveis constitui um dos tópicos mais investigados na área, apresentando exigências técnicas particulares que a diferenciam da geração de ativos puramente cosméticos \cite{summerville2018procedural, maleki2024procedural}. A literatura distingue conteúdos cosméticos (como texturas e efeitos visuais, cuja alteração não modifica o resultado das ações do jogador) de conteúdos funcionais, que afetam diretamente a mecânica e a dinâmica da partida \cite{summerville2018procedural}. Fases de jogos são artefatos intrinsecamente funcionais: uma fase visualmente atraente, mas que apresente falhas topológicas ou bloqueios estruturais que impeçam sua conclusão, é considerada inválida para o propósito do jogo \cite{summerville2018procedural, linden2014dungeons}.

Na geração procedural de níveis e masmorras, um método típico envolve três componentes conceituais: (1) um modelo representacional abstrato que descreve a estrutura global da fase; (2) um método para construir esse modelo representacional; e (3) um método para sintetizar a geometria efetiva da fase a partir do modelo abstrato \cite{linden2014dungeons, shaker2016constructive}. Diversas famílias algorítmicas foram desenvolvidas para essas etapas. Algoritmos de particionamento de espaço, como o Particionamento Binário do Espaço (\textit{Binary Space Partitioning} --- BSP), \textit{quadtrees} (em 2D) e \textit{octrees} (em 3D), subdividem recursivamente o domínio espacial em células disjuntas, garantindo estruturalmente que salas alocadas em partições distintas não se sobreponham \cite{shaker2016constructive}. Em contrapartida, abordagens baseadas em agentes escavadores (\textit{agent-based dungeon growing}) operam em escala local, deslocando um agente virtual que escava túneis e posiciona salas de maneira passo a passo; embora produzam configurações orgânicas, agentes sem mecanismos de inspeção antecipada (\textit{look-ahead}) estão sujeitos a gerar salas sobrepostas e corredores sem saída \cite{shaker2016constructive}.

De modo complementar, autômatos celulares operam sobre grades discretas aplicando regras de transição baseadas na vizinhança local (como as vizinhanças de Moore ou de von Neumann) para formar cavernas naturais, mas não asseguram, por si sós, a conectividade global entre câmaras isoladas, exigindo etapas adicionais de verificação e abertura de passagens \cite{linden2014dungeons, shaker2016constructive}. Gramáticas generativas --- incluindo gramáticas de grafos e gramáticas de formas (\textit{shape grammars}) --- expandem um axioma inicial por meio de regras de produção e reescrita, permitindo capturar princípios de progressão e hierarquia em níveis de ação e aventura \cite{dormans2010, linden2013designing, shaker2016constructive}. Apesar da maturidade dessas técnicas em domínios bidimensionais, a transição para a geração de fases tridimensionais complexas permanece um desafio aberto na literatura devido ao aumento expressivo das restrições geométricas, de navegação vertical e de acoplamento espacial \cite{linden2014dungeons, maleki2024procedural, xu2025constraint}.

% ---------------------------------------------------------------------------
\subsection{PCG Baseada em Aprendizado de Máquina e Uso de LLMs}
\label{sec:pcgml_llms}
% ---------------------------------------------------------------------------

\subsubsection{\textit{Procedural Content Generation via Machine Learning} --- PCGML}
\label{sec:pcgml}

A Geração Procedural de Conteúdo via Aprendizado de Máquina (\textit{Procedural Content Generation via Machine Learning} --- PCGML) é definida formalmente como a geração de conteúdos de jogos diretamente por modelos estatísticos ou de aprendizado de máquina treinados sobre conjuntos de conteúdos preexistentes de jogos \cite{summerville2018procedural}. A distinção fundamental entre o PCGML e as abordagens clássicas baseadas em busca ou em solucionadores reside no mecanismo gerador: enquanto a busca evolucionária pode eventualmente empregar um modelo neural apenas como função de avaliação para pontuar candidatos sintetizados por outros meios, no PCGML o modelo treinado internaliza a distribuição estatística do \textit{corpus} de entrada e amostra diretamente novas instâncias a partir de sua inferência \cite{summerville2018procedural}.

A taxonomia clássica de PCGML organiza os métodos segundo a representação subjacente dos dados (sequências unidimensionais, grades bidimensionais/tridimensionais e grafos) e o paradigma de modelagem empregado \cite{summerville2018procedural}. Entre as técnicas consolidadas figuram: (i) modelos probabilísticos de contagem de frequência, como modelos de $n$-gramas, Cadeias de Markov Multidimensionais (MdMCs) e Campos Aleatórios de Markov (MRFs); (ii) métodos de redução de dimensionalidade e fatoração de matrizes (como PCA); (iii) Redes Neurais Recorrentes, especialmente redes \textit{Long Short-Term Memory} (LSTMs), utilizadas na geração sequencial de fases baseadas em blocos; (iv) \textit{Autoencoders} e \textit{Variational Autoencoders} (VAEs), que mapeiam fases em espaços latentes contínuos para fins de interpolação e reparo de segmentos; e (v) Redes Adversárias Generativas (GANs) \cite{summerville2018procedural, maleki2024procedural}.

No âmbito do aprendizado por reforço aplicado à geração (PCGRL), o paradigma difere do treinamento supervisionado sobre \textit{corpora} estáticos \cite{khalifa2020pcgrl}: o agente gerador interage dinamicamente com um ambiente de desenvolvimento de fases, modificando células ou inserindo peças por meio de ações discretas e sendo otimizado para maximizar uma função de recompensa cumulativa que quantifica propriedades de jogabilidade e solvabilidade \cite{khalifa2020pcgrl, maleki2024procedural}.

É conceitualmente relevante distinguir a utilização de Modelos de Linguagem de Grande Porte (LLMs) em relação à definição estrita de PCGML formulada por Summerville et al. \cite{summerville2018procedural}. Quando um LLM fundacional é empregado por meio de engenharia de \textit{prompt} em regime de zero-shot ou few-shot, o modelo não foi treinado sobre um \textit{corpus} de fases de jogos específico, mas opera por meio de aprendizado em contexto (\textit{in-context learning}) suportado pelo conhecimento estrutural e semântico adquirido durante o pré-treinamento geral em linguagem natural e código \cite{brown2020language, Gallotta_2024}. Por outro lado, caso um modelo de linguagem seja expressamente submetido a ajuste fino (\textit{fine-tuning}) sobre conjuntos de dados de fases (como representações em texto, ASCII ou grafos de níveis), sua operação alinha-se estritamente à definição clássica de PCGML \cite{summerville2018procedural, todd2023level}.

Embora o aprendizado a partir de exemplos atenue a necessidade de codificação manual de regras estilísticas, a literatura aponta duas limitações estruturais no paradigma de PCG baseado em modelos conexionistas \cite{summerville2018procedural, maleki2024procedural, her2025zeroshot3dmapgeneration}: a escassez de dados (\textit{data scarcity}), decorrente do número restrito de fases digitalizadas e da ausência de representações universais entre diferentes jogos \cite{summerville2018procedural}; e a dificuldade de garantir propriedades funcionais globais, visto que modelos puramente estatísticos tendem a reproduzir padrões visuais plausíveis localmente, mas não oferecem garantias determinísticas de conectividade, acessibilidade ou solvabilidade global sem módulos externos de validação \cite{summerville2018procedural, maleki2024procedural}.

\subsubsection{Modelos de Linguagem de Grande Porte Aplicados à Geração de Conteúdo}
\label{sec:llms_pcg}

Os Modelos de Linguagem de Grande Porte (\textit{Large Language Models} --- LLMs) constituem uma classe de modelos fundacionais de aprendizado profundo fundamentados na arquitetura \textit{Transformer} \cite{vaswani2017attention, Gallotta_2024}. Historicamente, os modelos neurais sequenciais precursores em Processamento de Linguagem Natural (PLN), tais como Redes Neurais Recorrentes (RNNs) e LSTMs, enfrentavam severos entraves para capturar dependências de longo alcance devido ao fenômeno do esvanecimento ou explosão de gradientes (\textit{vanishing/exploding gradients}) durante o treinamento por retropropagação ao longo do tempo \cite{bengio1994learning}. A arquitetura \textit{Transformer} superou essa barreira teórica ao prescindir da recorrência sequencial e introduzir o mecanismo de autoatenção (\textit{self-attention}) \cite{vaswani2017attention}. O mecanismo de autoatenção calcula diretamente coeficientes de ponderação entre todos os pares de elementos (\textit{tokens}) da sequência de entrada, viabilizando o processamento altamente paralelizado e permitindo correlacionar informações estruturais e semânticas independentemente de sua distância posicional no texto \cite{vaswani2017attention, Gallotta_2024}.

Na vertente generativa preponderante, os modelos autorregressivos causais operam estimando a distribuição de probabilidade condicional do próximo \textit{token} $x_t$ com base no histórico de \textit{tokens} precedentes:
\begin{equation}
    P(x_1, x_2, \dots, x_T) = \prod_{t=1}^{T} P(x_t \mid x_1, x_2, \dots, x_{t-1})
\end{equation}
O ciclo de desenvolvimento desses modelos envolve tipicamente o pré-treinamento autossupervisionado em larga escala sobre extensos \textit{corpora} de textos e linguagens de programação, seguido por ajuste fino supervisionado para seguimento de instruções (\textit{instruction tuning}) e alinhamento por preferências (como RLHF) \cite{brown2020language, Gallotta_2024, her2025zeroshot3dmapgeneration}. Em modelos com bilhões de parâmetros, destaca-se a capacidade de aprendizado em contexto (\textit{In-Context Learning} --- ICL), que possibilita a resolução de novas tarefas a partir exclusivamente de especificações no \textit{prompt} em regime \textit{zero-shot} ou condicionada por poucos exemplos demonstrativos (\textit{few-shot}), sem modificação dos pesos sinápticos da rede \cite{brown2020language}. Estruturações avançadas de \textit{prompting}, como Cadeia de Pensamento (\textit{Chain-of-Thought} --- CoT) e a incorporação de Geração Aumentada por Recuperação (\textit{Retrieval-Augmented Generation} --- RAG), estendem essa capacidade permitindo decomposição analítica de etapas e consulta a fontes de dados estruturados \cite{Gallotta_2024, tang2025unrealllm, wang2026designerlyloop}.

Na literatura de PCG, o emprego de LLMs diversificou-se em múltiplas funções especializadas \cite{Gallotta_2024, maleki2024procedural}:
\begin{enumerate}
    \item \textbf{Gerador direto de representações e conteúdo:} síntese de artefatos textuais como diálogos e narrativas, bem como mapas bidimensionais codificados diretamente em texto ou caracteres ASCII (como grades para \textit{Sokoban} ou \textit{Super Mario Bros.}) \cite{todd2023level, nasir2024word2worldgeneratingstoriesworlds};
    \item \textbf{Planejador lógico e semântico de alto nível:} geração de estruturas conceituais abstratas, incluindo árvores de missões, grafos de progressão lógica e atribuição contextual de temas narrativos a salas e regiões \cite{Gallotta_2024, xu2025databasedrivenframework3dlevel, wang2026designerlyloop};
    \item \textbf{Sintetizador de parâmetros e código para ferramentas procedurais:} produção de rotinas em linguagens de programação (como scripts para motores ou modeladores 3D), parametrização de geradores construtivos e formulação de restrições para solucionadores matemáticos \cite{hu2024scenecraft, tang2025unrealllm, xu2025constraint, yang2026codeasroomgenerating3drooms};
    \item \textbf{Assistente de design em sistemas de iniciativa mista:} atuação colaborativa na interpretação de comandos do designer, oferecendo sugestões estruturais e refinando especificações conceituais \cite{Gallotta_2024, bandeali2026generative, wang2026designerlyloop}.
\end{enumerate}

\subsubsection{Geração de Representações Estruturadas}
\label{sec:representacoes_estruturadas}

Quando modelos de linguagem são integrados a sistemas computacionais externos --- tais como motores gráficos, otimizadores matemáticos ou verificadores formais ---, a emissão de respostas em linguagem natural livre torna-se inadequada devido à sua ambiguidade intrínseca e à suscetibilidade a falhas sintáticas na análise automatizada (\textit{parsing}) \cite{xu2025constraint, regmi2025system}. Para garantir a interoperabilidade entre a inferência probabilística do LLM e os procedimentos algorítmicos determinísticos, a literatura preconiza a geração de representações estruturadas intermediárias \cite{hu2024scenecraft, tang2025unrealllm, xu2025constraint, chen2025narrative, wang2026designerlyloop}.

Sob o ponto de vista conceitual e de engenharia de software, é indispensável diferenciar o \textit{modelo conceitual} do \textit{formato de serialização} \cite{xu2025constraint, xu2025databasedrivenframework3dlevel, chen2025narrative}:
\begin{itemize}
    \item O \textbf{modelo conceitual} (ou estrutura abstrata de informação) é a formalização lógica ou matemática que organiza o domínio do problema. Exemplos compreendem grafos de cena e de progressão \cite{xu2025databasedrivenframework3dlevel, wang2026designerlyloop}, conjuntos de triplas relacionais do tipo \texttt{[Objeto-Relação-Objeto]} \cite{chen2025narrative}, hierarquias de componentes de cena, linguagens de descrição de domínio específico (DSL) como a GLDL (\textit{Game Level Description Language}) \cite{xu2025constraint}, ou sínteses procedurais via código executável \cite{hu2024scenecraft, yang2026codeasroomgenerating3drooms};
    \item O \textbf{formato de serialização} é o padrão textual ou binário adotado para codificar, armazenar e transmitir esse modelo conceitual entre diferentes subsistemas \cite{sun2024roomplanner, xu2025databasedrivenframework3dlevel}. O formato JSON (\textit{JavaScript Object Notation}), bem como esquemas em XML ou YAML, constitui a representação hegemônica por estruturar dados hierárquicos em pares chave-valor e arranjos vetoriais com suporte a validação rigorosa de sintaxe por meio de esquemas formais (\textit{JSON Schema}) \cite{tang2025unrealllm, sun2024roomplanner, xu2025databasedrivenframework3dlevel}. Abordagens pontuais na literatura, a exemplo de Öcal et al. \cite{ocal2024sceneteller}, adaptaram experimentalmente a sintaxe de folhas de estilo (CSS) para declarar parâmetros de caixas delimitadoras; não obstante, o padrão JSON consolida-se como a convenção primária de serialização para intercâmbio de dados estruturados.
\end{itemize}

O formato de serialização não se confunde, portanto, com a lógica do nível: ele é o meio físico-textual de transmissão que transporta a abstração formal da fase para que ferramentas computacionais possam interpretá-la e validá-la de modo determinístico \cite{tang2025unrealllm, xu2025databasedrivenframework3dlevel}.

% ---------------------------------------------------------------------------
\subsection{Raciocínio e Ancoragem Espacial em LLMs}
\label{sec:spatial_grounding}
% ---------------------------------------------------------------------------

\subsubsection{Raciocínio Espacial}
\label{sec:raciocinio_espacial}

O raciocínio espacial compreende a capacidade de representar formalmente entidades em um domínio multidimensional, inferir suas propriedades geométricas e deduzir os desdobramentos de suas transformações e arranjos mútuos \cite{cohn2008qualitative, sun2024roomplanner, xu2025constraint}. No domínio do \textit{Level Design} e da síntese de ambientes virtuais 3D, o raciocínio espacial subdivide-se entre relações qualitativas (topológicas e mereológicas) \cite{cohn2008qualitative} e relações quantitativas métricas contínuas \cite{xu2025constraint, chen2025narrative}.

Na formulação da organização espacial de uma fase, identificam-se propriedades e relações geométricas essenciais \cite{cohn2008qualitative, sun2024roomplanner, xu2025constraint, chen2025narrative}:
\begin{itemize}
    \item \textbf{Posição:} localização absoluta ou relativa de uma entidade no espaço $\mathbb{R}^3$, expressa em relação à origem global ou a um ponto de referência local;
    \item \textbf{Distância e proximidade:} separação métrica euclidiana entre entidades, determinando adjacência funcional, raio de alcance ou zonas de segurança;
    \item \textbf{Direção versus Orientação:} sob o rigor do raciocínio espacial \cite{cohn2008qualitative}, a \textit{direção} constitui uma relação espacial relativa entre duas entidades (indicando, por exemplo, que a sala $B$ está situada ao norte, a leste ou acima da sala $A$), enquanto a \textit{orientação} refere-se ao alinhamento angular intrínseco de uma entidade (sua \textit{pose} ou vetor de mira frontal --- \textit{facing} --- no espaço, determinado por suas rotações em torno dos eixos locais);
    \item \textbf{Adjacência e conectividade:} contato geométrico ou comunicação de passagem entre regiões contíguas, tais como salas que compartilham portais comuns;
    \item \textbf{Contenção:} exigência topológica e volumétrica de que determinado conjunto de elementos (como itens ou mobílias) esteja integralmente situado dentro do perímetro delimitador de uma sala ou compartimento hospedeiro;
    \item \textbf{Interseção e ocupação volumétrica:} verificação de compartilhamento de pontos no espaço tridimensional entre dois ou mais corpos físicos.
\end{itemize}

A preservação concorrente da consistência de todas essas relações ao longo da alocação de dezenas de elementos arquitetônicos e funcionais constitui um problema combinatório e contínuo de alta complexidade \cite{xu2025constraint, feng2026repurposing3dgenerativemodel}.

\subsubsection{\textit{Spatial Grounding}}
\label{sec:conceito_spatial_grounding}

O problema de ancoragem semântica (\textit{symbol grounding problem}), formulado originalmente por Harnad \cite{harnad1990symbol}, investiga como símbolos abstratos manipulados puramente com base em regras sintáticas podem adquirir significado intrínseco e correspondência funcional com entidades concretas de um mundo físico ou simulado. No âmbito dos sistemas generativos aplicados a ambientes 3D, a \textit{ancoragem espacial} (\textit{spatial grounding}) designa a capacidade de mapear comandos linguísticos e especificações conceituais (como ``posicionar uma mesa no centro da câmara'' ou ``vincular o corredor de acesso à ala norte'') em coordenadas métricas coerentes, dimensões tridimensionais compatíveis e restrições espaciais verificáveis \cite{harnad1990symbol, hu2024scenecraft, sun2024roomplanner, chen2025narrative}.

Um componente crítico da ancoragem espacial consiste na interpretação correta de \textit{affordances}. Na definição clássica da psicologia ecológica de Gibson \cite{gibson1979ecological}, uma \textit{affordance} traduz as possibilidades reais de ação que os elementos do ambiente físico oferecem a um agente (por exemplo, uma porta aberta oferece passagem, uma superfície horizontal firme oferece suporte para caminhar ou sentar, e uma parede oferece bloqueio ou proteção) \cite{gibson1979ecological}. A ancoragem espacial requer que o sistema não apenas reconheça a coocorrência semântica entre termos linguísticos, mas garanta que as configurações geométricas resultantes respeitem as propriedades funcionais de uso: corredores devem manter altura e largura livres para circulação de agentes, portas devem possuir desobstrução frontal para acesso e objetos de apoio devem respeitar o contato estável de superfícies no espaço tridimensional \cite{sun2024roomplanner, xu2025constraint}.

\subsubsection{Limitações dos LLMs em Tarefas Espaciais Tridimensionais}
\label{sec:limitacoes_espaciais_llms}

No contexto dos modelos generativos de linguagem, define-se por alucinação a síntese de respostas que apresentam correção superficial ou gramatical, mas que violam a coerência factual, contêm premissas contraditórias ou desrespeitam restrições explícitas do domínio \cite{Gallotta_2024, xu2025constraint, bandeali2026generative}. Quando os LLMs são incumbidos de determinar diretamente atributos geométricos contínuos em ambientes 3D --- tais como coordenadas cartesianas $(x, y, z)$, caixas delimitadoras e ângulos de rotação para múltiplos elementos interconectados ---, esse comportamento manifesta-se sob a forma de \textit{alucinação espacial} ou inconsistência métrica \cite{Gallotta_2024, sun2024roomplanner, yang2026codeasroomgenerating3drooms, feng2026repurposing3dgenerativemodel}.

A fundamentação teórica dessa limitação decorre de características constitutivas dos modelos de linguagem \cite{Gallotta_2024, hu2024scenecraft, sun2024roomplanner, feng2026repurposing3dgenerativemodel}:
\begin{enumerate}
    \item \textbf{Ausência de modelagem física e geométrica na inferência:} LLMs são treinados com base na verossimilhança estatística de sequências de \textit{tokens} textuais, sem acesso nativo a um motor de cálculo vetorial ou representação interna de espaço euclidiano tridimensional durante o cálculo de autoatenção \cite{Gallotta_2024, sun2024roomplanner};
    \item \textbf{Severa compressão de informação geométrica:} ao representar geometrias poligonais complexas exclusivamente por rótulos simbólicos e tuplas numéricas discretizadas, omitem-se as interações de superfície contínuas e as relações de contato entre sólidos \cite{feng2026repurposing3dgenerativemodel};
    \item \textbf{Deriva autorregressiva e acúmulo de erros métricos:} em processos de geração sequencial ou em planejamentos em passo único (\textit{single-step}), pequenos desvios numéricos cometidos nas primeiras coordenadas preditas propagam-se recursivamente pelo contexto, degradando a consistência geométrica global conforme o número de elementos cresce no ambiente \cite{sun2024roomplanner, feng2026repurposing3dgenerativemodel, yang2026codeasroomgenerating3drooms}.
\end{enumerate}

Diversos estudos empíricos recentes ilustram concretamente essas falhas:
\begin{itemize}
    \item No sistema \textit{SceneCraft}, Hu et al. \cite{hu2024scenecraft} evidenciaram que, ao gerar scripts em Python para modeladores 3D como o Blender, os LLMs conseguem estruturar a hierarquia de objetos, mas cometem graves imprecisões nas coordenadas numéricas, gerando colisões e desalinhamentos severos que exigem otimizadores numéricos acoplados para corrigir a cena;
    \item No framework \textit{RoomPlanner}, Sun et al. \cite{sun2024roomplanner} constataram empiricamente sobre o conjunto de dados Matterport3D que a solicitação direta de caixas delimitadoras a LLMs resulta em taxas substanciais de interpenetração de móveis e posicionamento de elementos fora das paredes dos cômodos;
    \item No estudo \textit{Code-as-Room}, Yang et al. \cite{yang2026codeasroomgenerating3drooms} reportaram que a síntese de recintos 3D via código por LLMs padece de alucinações espaciais e violações funcionais de circulação quando desprovida de verificação externa de navegabilidade.
\end{itemize}

As falhas observadas agrupam-se em quatro categorias recorrentes \cite{sun2024roomplanner, xu2025constraint, feng2026repurposing3dgenerativemodel, yang2026codeasroomgenerating3drooms}: (i) \textbf{colisões e sobreposições volumétricas (\textit{overlapping} / \textit{clipping})}, nas quais malhas sólidas ocupam o mesmo volume simultaneamente; (ii) \textbf{violações de limites e suporte (\textit{out-of-bounds} e objetos flutuantes)}, em que elementos atravessam paredes externas ou ficam suspensos sem contato com o piso; (iii) \textbf{erros de rotação e alinhamento}, como portas e armários posicionados de frente para paredes cegas; e (iv) \textbf{obstruções de circulação}, nas quais objetos bloqueiam corredores vitais e fragmentam o espaço navegável \cite{sun2024roomplanner, xu2025constraint, xu2025databasedrivenframework3dlevel}.

Tais constatações consolidam a compreensão de que, embora os LLMs demonstrem notável competência para o raciocínio semântico e planejamento conceitual, a determinação direta de layouts tridimensionais estritamente por inferência de linguagem é inerentemente falha, demandando arquiteturas estruturadas que deleguem a instanciação geométrica e a validação espacial a componentes algorítmicos dedicados \cite{hu2024scenecraft, sun2024roomplanner, xu2025constraint, xu2025databasedrivenframework3dlevel, yang2026codeasroomgenerating3drooms}.

% ---------------------------------------------------------------------------
\subsection{Representação Lógica de Fases por Grafos}
\label{sec:grafos}
% ---------------------------------------------------------------------------

\subsubsection{Conceitos Fundamentais de Teoria dos Grafos}
\label{sec:conceitos_grafos}

A Teoria dos Grafos fornece a base matemática formal para modelar relações discretas e conectividades estruturais em sistemas complexos \cite{cormen2009introduction, summerville2018procedural}. Formalmente, um grafo é definido como um par ordenado $G = (V, E)$, no qual $V$ consiste em um conjunto finito e não vazio de vértices (ou nós) e $E$ representa um conjunto de arestas (ou arcos) que conectam pares de vértices pertencentes a $V$ \cite{cormen2009introduction, maleki2024procedural}.

Quanto à orientação de suas arestas, um grafo pode ser não direcionado --- em que as arestas indicam conexões simétricas bidirecionais entre pares de vértices $\{u, v\}$ ---, ou direcionado (dígrafo), quando cada aresta consiste em um par ordenado $(u, v) \in E$ orientado do vértice de origem $u$ para o vértice de destino $v$ \cite{cormen2009introduction, dormans2010}. Em dígrafos, distingue-se o grau de entrada ($d^-(v)$, número de arestas que incidem em $v$) do grau de saída ($d^+(v)$, número de arestas que partem de $v$) \cite{cormen2009introduction}. Um grafo ponderado associa valores numéricos reais às suas arestas (como distâncias métricas, tempos de travessia ou dificuldades), enquanto um grafo rotulado ou de atributos vincula tipos, símbolos e dicionários de propriedades tanto a vértices quanto a arestas \cite{dormans2010, linden2013designing, linden2014dungeons, chen2025narrative}.

Sobre essas estruturas, fundamentam-se propriedades topológicas cruciais \cite{cormen2009introduction, linden2014dungeons, summerville2018procedural}:
\begin{itemize}
    \item \textbf{Caminho (\textit{path}):} sequência de vértices $\langle v_0, v_1, \dots, v_k \rangle$ tal que $(v_{i-1}, v_i) \in E$ para todo $i \in \{1, \dots, k\}$;
    \item \textbf{Ciclo:} caminho fechado no qual $v_0 = v_k$ e que contém pelo menos uma aresta distinta. Um grafo que não possui ciclos é dito acíclico;
    \item \textbf{Grafo Acíclico Direcionado (DAG --- \textit{Directed Acyclic Graph}):} dígrafo desprovido de caminhos orientados fechados, estrutura fundamental para expressar relações de precedência estrita e causalidade \cite{cormen2009introduction, wang2026designerlyloop};
    \item \textbf{Conectividade e Conectividade Forte:} em grafos não direcionados, a conectividade assegura que exista ao menos um caminho entre qualquer par de vértices. Em dígrafos, define-se conectividade forte quando, para todo par de vértices $u, v \in V$, existem simultaneamente caminhos direcionados de $u$ para $v$ e de $v$ para $u$ \cite{cormen2009introduction};
    \item \textbf{Alcançabilidade (\textit{reachability}):} condição na qual existe um caminho orientado viável partindo de um nó inicial $v_{\text{start}}$ até um nó de destino $v_{\text{target}}$ \cite{cormen2009introduction, xu2025databasedrivenframework3dlevel};
    \item \textbf{Planaridade:} propriedade geométrica de um grafo que permite desenhá-lo em um plano bidimensional $\mathbb{R}^2$ de modo que suas arestas se interceptem unicamente nos vértices aos quais estão conectadas \cite{cormen2009introduction, linden2013designing}.
\end{itemize}

\subsubsection{Grafos Aplicados ao \textit{Level Design}}
\label{sec:grafos_level_design}

Na modelagem de fases e masmorras procedurais, a literatura apoia-se em grafos sob duas perspectivas complementares: grafos de espaço (topológicos) e grafos de missão (ou de ações) \cite{dormans2010, linden2013designing, linden2014dungeons, summerville2018procedural}.

Nos \textit{grafos topológicos de espaço}, cada vértice $v \in V$ representa uma região física discreta do nível (como uma sala, corredor, câmara ou pátio), e cada aresta $e \in E$ representa uma passagem física ou conexão de adjacência (portas, arcos de passagem, escadarias ou pontes) \cite{dormans2010, linden2014dungeons, summerville2018procedural}. Os vértices recebem atributos temáticos e funcionais (sala de início, sala de combate, câmara de tesouro, sala de chefe), enquanto as arestas recebem atributos de passagem (passagem aberta, porta trancada, passagem unidirecional com queda) \cite{linden2014dungeons, xu2025databasedrivenframework3dlevel}.

Nos \textit{grafos de missão} ou \textit{grafos de ações} (\textit{action graphs}), por outro lado, os vértices representam eventos, tarefas ou ações que o jogador executa (ex.: obter chave, ativar alavanca, derrotar sentinela, transpor portal), e as arestas direcionadas codificam restrições causais de precedência no tempo \cite{dormans2010, linden2013designing}. Conforme demonstrado seminalmente por Dormans \cite{dormans2010} e aprofundado por van der Linden et al. \cite{linden2013designing, linden2014dungeons}, um grafo de ações pode ser mapeado em um grafo de salas mediante o agrupamento de ações que compartilham a mesma localização física (\textit{co-located actions}), unificando a estrutura de progressão e a topologia arquitetônica.

\subsubsection{Grafos de Progressão e Dependências}
\label{sec:grafos_progressao}

A integração formal entre a conectividade espacial e as mecânicas de jogo estabelece os \textit{grafos de progressão}, nos quais arestas e vértices são condicionados a regras de acesso lógico \cite{dormans2010, linden2013designing, linden2014dungeons, xu2025databasedrivenframework3dlevel}. O paradigma central para estruturar a progressão em masmorras e fases de aventura é o quebra-cabeça de chave e fechadura (\textit{lock-and-key puzzle}) \cite{dormans2010, linden2014dungeons, xu2025databasedrivenframework3dlevel}. Uma fechadura denota qualquer barreira ou pré-requisito funcional (uma porta física trancada, uma barreira mágica, um obstáculo ambiental intransponível), enquanto a chave corresponde ao elemento ou estado necessário para desativá-la (uma chave física, uma habilidade de salto adquirida, o acionamento de um botão remoto ou a eliminação de inimigos) \cite{dormans2010, linden2013designing, shaker2016constructive}.

Do ponto de vista formal, a estrutura de dependências lógicas induz um subgrafo direcionado de pré-condições $G_{\text{dep}} = (V_{\text{dep}}, E_{\text{dep}})$, o qual deve satisfazer estritamente a condição de ser um Grafo Acíclico Direcionado (DAG) para garantir a inexistência de impasses circulares insolúveis (\textit{deadlocks}) \cite{cormen2009introduction, xu2025databasedrivenframework3dlevel}. Em tal DAG de dependências, existe ao menos uma ordenação topológica $\tau: V_{\text{dep}} \to \{1, 2, \dots, |V_{\text{dep}}|\}$ tal que, para qualquer aresta orientada $(u, v) \in E_{\text{dep}}$, tem-se $\tau(u) < \tau(v)$ \cite{cormen2009introduction}. Isso significa que todo recurso ou chave $u$ indispensável para a superação do bloqueio em $v$ deve ser alcançável e coletado antes que o estado de travessia de $v$ seja solicitado \cite{xu2025databasedrivenframework3dlevel}.

É essencial pontuar o contraste estrutural entre a lógica de dependências e a circulação física: enquanto o grafo causal de pré-condições ($G_{\text{dep}}$) deve ser estritamente acíclico para ser solucionável, o grafo espacial de navegação do nível pode conter ciclos controlados \cite{dormans2010, linden2014dungeons}. Esses ciclos na geometria física enriquecem a exploração ao fornecer rotas alternativas, atalhos destraváveis e a necessidade intencional de retorno a áreas previamente visitadas (\textit{backtracking}) após a obtenção de novos recursos \cite{dormans2010, linden2013designing, linden2014dungeons}. A solvabilidade do nível modela-se formalmente no espaço de estados do jogo ($S$), no qual cada estado $s \in S$ é composto pela localização espacial do agente associada ao seu inventário corrente de chaves, itens e modificadores de mundo \cite{dormans2010, linden2014dungeons, xu2025databasedrivenframework3dlevel}.

\subsubsection{Estrutura Lógica versus Representação Espacial}
\label{sec:logica_vs_espacial}

A Teoria dos Grafos elucida a separação ontológica entre a estrutura lógica e a representação geométrica: um grafo de progressão é uma entidade abstrata, topológica e combinatória, e não uma especificação espacial métrica \cite{dormans2010, summerville2018procedural}. Seus vértices e arestas estipulam a conectividade e a causalidade entre elementos, mas não contêm informações a respeito de posições tridimensionais $(x, y, z)$, comprimentos em metros, geometrias de paredes ou desníveis verticais \cite{dormans2010, linden2013designing, linden2014dungeons}.

Essa dissociação estabelece uma relação de mapeamento de um-para-muitos: um mesmo grafo de progressão admite infinitas configurações espaciais distintas no espaço contínuo tridimensional, contanto que as adjacências e ordenações lógicas sejam satisfeitas \cite{dormans2010, linden2013designing}. Todavia, nem todo grafo lógico abstrato é imediatamente realizável em uma grade ortogonal ou em um plano arquitetônico contínuo \cite{linden2013designing, linden2014dungeons}. Se um vértice do grafo contiver um número de arestas incidentes (grau) superior ao número físico de portas comportadas pela geometria do módulo arquitetônico, ou se a topologia for não planar no plano bidimensional, a instanciação física necessita de técnicas de planarização e expansão estrutural, tais como a introdução algorítmica de nós intermediários (corredores ramificados, bifurcações ou transições verticais entre múltiplos pavimentos), conforme formalizado por van der Linden et al. \cite{linden2013designing, linden2014dungeons}.

% ---------------------------------------------------------------------------
\subsection{Representação de \textit{Layouts} Espaciais 3D}
\label{sec:layout_3d}
% ---------------------------------------------------------------------------

\subsubsection{Representação Modular de Ambientes}
\label{sec:representacao_modular}

Para converter a especificação abstrata de um grafo em um ambiente interativo tridimensional, a literatura de PCG e os motores de jogos apoiam-se no paradigma da modelagem espacial modular \cite{linden2014dungeons, shaker2016constructive, xu2025databasedrivenframework3dlevel}. Nessa abordagem, o espaço não é concebido como uma malha estática singular, mas através da composição combinatória de unidades tridimensionais reutilizáveis e pré-autoradas \cite{linden2014dungeons, shaker2016constructive, xu2025databasedrivenframework3dlevel}.

As granularidades modulares típicas abrangem \cite{linden2014dungeons, shaker2016constructive, vieira2024metodologia}:
\begin{itemize}
    \item \textbf{Grades regulares e \textit{voxels}:} partição do espaço em células discretas uniformes em 2D ou 3D, atribuindo-se a cada célula um elemento primitivo de geometria (bloco sólido, piso, parede, abertura ou ar livre) \cite{pcgbook, shaker2016constructive, summerville2018procedural};
    \item \textbf{Segmentos compostos (\textit{chunks}):} subestruturas volumétricas que agrupam conjuntos pré-definidos de blocos e acessórios temáticos \cite{linden2014dungeons, shaker2016constructive};
    \item \textbf{Módulos de salas e \textit{prefabs} 3D:} blocos arquitetônicos inteiros (como salas com dimensões estipuladas, corredores em L, junções em T e escadarias) e ativos pré-fabricados (\textit{prefabs}) de objetos internos \cite{linden2014dungeons, xu2025databasedrivenframework3dlevel, regmi2025system}. Um \textit{prefab} encapsula a malha gráfica 3D (\textit{mesh}), materiais, texturas, colunas de suporte, volumes de colisão física e pontos de ancoragem semânticos em um único pacote modular \cite{xu2025databasedrivenframework3dlevel, regmi2025system}.
\end{itemize}

O acoplamento topológico e geométrico entre módulos adjacentes é operacionalizado por meio de conectores ou portas (\textit{connectors} / \textit{doorways} / \textit{sockets}) posicionados estrategicamente nas extremidades dos blocos modulares \cite{linden2013designing, linden2014dungeons, shaker2016constructive, sun2024roomplanner}. Cada conector define uma posição métrica relativa em relação ao pivô do módulo, a dimensão da abertura e um vetor de orientação normal que especifica a direção para a qual a passagem se projeta no espaço \cite{linden2013designing, shaker2016constructive, sun2024roomplanner, xu2025databasedrivenframework3dlevel}.

\subsubsection{Relações e Transformações Espaciais}
\label{sec:transformacoes_espaciais}

A instanciação de peças modulares no espaço contínuo $\mathbb{R}^3$ fundamenta-se nas transformações afins da álgebra linear, articulando o referencial de coordenadas local de cada objeto e o referencial global do cenário (\textit{world space}) \cite{ericson2004realtime, sun2024roomplanner, xu2025constraint}.

A especificação espacial completa de um componente $i$ exige distinguir sua \textit{pose}, suas dimensões físicas e seus fatores de escala \cite{ericson2004realtime, sun2024roomplanner, xu2025constraint}:
\begin{itemize}
    \item A \textbf{pose} define a localização e a atitude espacial do componente no mundo, decomposta na translação $\mathbf{p}_i = (x_i, y_i, z_i) \in \mathbb{R}^3$ (posição de seu ponto de pivô) e na rotação $\mathbf{R}_i \in SO(3)$ (orientação angular, expressa por matrizes ortogonais, quatérnios ou ângulos de Euler, tipicamente parametrizada em layouts pelo ângulo de guinada $\theta_i$ em torno do eixo vertical $Y$ para rotações ortogonais de $90^\circ$ ou $180^\circ$) \cite{sun2024roomplanner, xu2025databasedrivenframework3dlevel};
    \item As \textbf{dimensões nominais} $(w_i, l_i, h_i) \in \mathbb{R}^3$ representam a largura, o comprimento e a altura intrínsecos do volume do objeto no espaço local;
    \item O \textbf{fator de escala} $\mathbf{s}_i = (s_{x,i}, s_{y,i}, s_{z,i}) \in \mathbb{R}^3$ expressa a ampliação ou redução proporcional aplicada à geometria nos eixos coordenados.
\end{itemize}

Para a detecção de colisões e verificação de proximidade geométrica em tempo computacional eficiente, a computação gráfica recorre a volumes delimitadores (\textit{bounding volumes}) \cite{ericson2004realtime, sun2024roomplanner, xu2025constraint}. Distinguem-se fundamentalmente dois tipos:
\begin{enumerate}
    \item \textbf{Caixa Delimitadora Alinhada aos Eixos (\textit{Axis-Aligned Bounding Box} --- AABB):} definida pelos intervalos extremos $[x_{\min}, x_{\max}] \times [y_{\min}, y_{\max}] \times [z_{\min}, z_{\max}]$ paralelamente aos eixos globais do sistema de coordenadas \cite{ericson2004realtime}. Testes de interseção entre AABBs possuem custo computacional extremamente baixo por meio de comparações escalares de intervalos em cada eixo; contudo, quando um objeto sofre rotação angular arbitrária, sua AABB precisa ser recalculada e expandida para envolver os novos limites, introduzindo espaço vazio excessivo que pode gerar falsos positivos de colisão \cite{ericson2004realtime};
    \item \textbf{Caixa Delimitadora Orientada (\textit{Oriented Bounding Box} --- OBB):} paralelepípedo delimitador centrado no objeto cujos eixos permanecem solidários à rotação e orientação da própria geometria local \cite{ericson2004realtime}. Embora exija testes de separação mais sofisticados (como o Teorema dos Eixos Separadores --- SAT), a OBB ajusta-se com precisão estrita ao volume sólido sob qualquer rotação espacial \cite{ericson2004realtime}.
\end{enumerate}

\subsubsection{Restrições Geométricas}
\label{sec:restricoes_geometricas}

A montagem de um \textit{layout} tridimensional coerente requer a satisfação concorrente de um conjunto estrito de restrições geométricas \cite{linden2014dungeons, ericson2004realtime, sun2024roomplanner, xu2025constraint, xu2025databasedrivenframework3dlevel}:
\begin{itemize}
    \item \textbf{Restrição de não sobreposição volumétrica (\textit{non-overlap}):} sólidos arquitetônicos e objetos distintos não podem ocupar concomitantemente a mesma porção do espaço euclidiano tridimensional. Formalmente, a interseção volumétrica entre seus interiores sólidos deve ser estritamente nula ($V_{\text{overlap}} = 0$ ou $\text{IoU}_{3D} = 0$), evitando interpenetrações anômalas (\textit{clipping}) \cite{ericson2004realtime, sun2024roomplanner, xu2025constraint};
    \item \textbf{Restrição de contenção e fronteiras (\textit{boundary constraints}):} objetos internos associados a uma sala devem ter seus volumes contidos estritamente no interior do polígono ou prisma formado por suas paredes, piso e teto, sem atravessar a alvenaria perimetral \cite{sun2024roomplanner, xu2025constraint};
    \item \textbf{Compatibilidade e alinhamento de portais:} para que duas salas ou uma sala e um corredor compartilhem uma passagem transitável, seus respectivos conectores devem coincidir métrica e volumetricamente na fronteira comum, apresentando vetores normais opostos e colineares, sem que ocorra bloqueio por paredes cegas de um dos módulos \cite{linden2013designing, shaker2016constructive, sun2024roomplanner, xu2025databasedrivenframework3dlevel};
    \item \textbf{Desobstrução e espaço de circulação (\textit{clearance}):} a distribuição de objetos em um cômodo deve resguardar zonas livres desimpedidas ao redor de portas, escadas e rotas centrais de passagem, prevenindo o congestionamento ou adensamento físico indevido (\textit{overcrowding} / \textit{clutter}) que impeça o trânsito físico dos agentes virtuais \cite{sun2024roomplanner, xu2025constraint, xu2025databasedrivenframework3dlevel}.
\end{itemize}

% ---------------------------------------------------------------------------
\subsection{Arquiteturas Híbridas e Geração Baseada em Restrições}
\label{sec:arquiteturas_hibridas}
% ---------------------------------------------------------------------------

\subsubsection{Sistemas Híbridos e Neuro-Simbólicos}
\label{sec:sistemas_hibridos}

A análise aprofundada dos paradigmas generativos revela uma complementaridade ontológica: métodos puramente simbólicos, construtivos ou orientados a restrições fornecem rigor matemático e preservação determinística de invariantes funcionais, porém demandam complexa formalização manual de regras e tendem a apresentar menor flexibilidade estilística \cite{shaker2016constructive, summerville2018procedural}; inversamente, modelos baseados em redes profundas e LLMs exibem expressividade de linguagem, compreensão semântica e facilidade na geração de conteúdo criativo, porém sofrem de alucinações espaciais e ausência de garantias formais de consistência \cite{Gallotta_2024, maleki2024procedural}.

Essa constatação catalisou a emergência de \textit{arquiteturas híbridas neuro-simbólicas} em inteligência artificial e geração de fases \cite{garcez2023neurosymbolic, Gallotta_2024, xu2025constraint}. Conforme teorizado por Garcez e Lamb \cite{garcez2023neurosymbolic}, o paradigma neuro-simbólico postula a integração sinérgica entre o aprendizado conexionista baseado em dados (capacidade de inferência intuitiva, correlação estatística e semântica de alto nível) e o raciocínio simbólico (estruturas lógicas, grafos, solucionadores e verificação formal baseada em regras). No contexto de PCG para ambientes 3D, essa divisão funcional traduz-se na delegação das etapas de planejamento conceitual, estruturação lógica e contextualização ao módulo de linguagem (neural), ao passo que as rotinas de resolução métrica contínua, instanciação modular e validação geométrica são executadas por algoritmos determinísticos e solucionadores formais (simbólicos) \cite{hu2024scenecraft, tang2025unrealllm, chen2025narrative, xu2025constraint, xu2025databasedrivenframework3dlevel, her2025zeroshot3dmapgeneration}.

De modo análogo, essa divisão conceitual dialoga com os princípios da geração de iniciativa mista (\textit{Mixed-Initiative PCG}), fundamentada seminalmente em ferramentas como \textit{Tanagra} \cite{tanagra} e \textit{Sentient Sketchbook} \cite{liapis2013sentient}, nas quais o sistema computacional assume o papel de parceiro proativo que calcula métricas espaciais em tempo de autoria, resolve restrições e sugere correções estruturais mantendo o controle deliberativo do designer \cite{tanagra, liapis2013sentient, lai2020towards, walton2020evaluating, wang2026designerlyloop}.

\subsubsection{Geração Baseada em Regras e Restrições}
\label{sec:geracao_restricoes}

O suporte matemático para a imposição de regras em sistemas generativos reside na teoria de Problemas de Satisfação de Restrições (CSP --- \textit{Constraint Satisfaction Problem}) \cite{russell2021artificial}. Formalmente, um CSP é definido pela tupla $(X, D, C)$, na qual $X = \{X_1, X_2, \dots, X_n\}$ denota um conjunto de variáveis de decisão (como coordenadas espaciais, orientações ou tipos de sala), $D = \{D_1, D_2, \dots, D_n\}$ representa os domínios de valores admissíveis para cada variável, e $C = \{C_1, C_2, \dots, C_m\}$ constitui um conjunto de restrições que delimitam as combinações simultâneas de valores permitidas \cite{russell2021artificial}. A resolução de um CSP emprega algoritmos de retrocesso sistemático (\textit{backtracking}) potencializados por heurísticas de ordenação e mecanismos de propagação de restrições, a exemplo da consistência de arcos (\textit{arc consistency}), que elimina antecipadamente valores de domínios que comprovadamente violariam condições futuras \cite{russell2021artificial}.

No escopo do arranjo espacial e planejamento de fases, as restrições são categorizadas em duas ordens de rigor \cite{tanagra, xu2025constraint, xu2025databasedrivenframework3dlevel}:
\begin{itemize}
    \item \textbf{Restrições rígidas (\textit{hard constraints}):} condições mandatórias e invioláveis para a viabilidade do nível, cuja inobservância inviabiliza funcional ou fisicamente o artefato (tais como ausência de sobreposição sólida, contenção no perímetro e ordenação causal acíclica de chaves e portas) \cite{tanagra, sun2024roomplanner, xu2025constraint};
    \item \textbf{Restrições flexíveis (\textit{soft constraints}):} diretrizes de qualidade, preferências estéticas ou métricas de balanceamento que expressam objetivos desejáveis (tais como distâncias médias de deslocamento, agrupamento estilístico de móveis ou dispersão de desafios), cujo atendimento é otimizado sem configurar critério binário de rejeição \cite{xu2025constraint, xu2025databasedrivenframework3dlevel}.
\end{itemize}

Quando o problema de disposição espacial envolve a ponderação conjunta de restrições rígidas e flexíveis em domínios contínuos, ele é frequentemente modelado por analogia ao clássico Problema de \textit{Layout} de Instalações (\textit{Facility Layout Problem} --- FLP) da pesquisa operacional \cite{xu2025constraint, xu2025databasedrivenframework3dlevel}. Sob essa formulação, define-se uma função objetivo com penalidades ponderadas associadas a cada infração de distância ou posicionamento, buscando configurações que minimizem o custo global de violações \cite{xu2025constraint}.

\subsubsection{Validação e Reparo de Conteúdo Gerado}
\label{sec:validacao_reparo}

Dado que geradores conexionistas, estocásticos ou heurísticos não garantem a priori a total conformidade com restrições rígidas, a literatura formaliza quatro paradigmas teóricos para a garantia da correção do conteúdo \cite{pcgbook, summerville2018procedural, maleki2024procedural, xu2025databasedrivenframework3dlevel}:
\begin{enumerate}
    \item \textbf{Geração e teste com filtragem (Rejeição integral):} o gerador sintetiza o candidato e o submete a uma bateria de verificações algorítmicas; constatada qualquer infração rígida, o candidato é sumariamente descartado e o procedimento reiniciado \cite{pcgbook, maleki2024procedural};
    \item \textbf{Construção com retrocesso (\textit{Backtracking}):} o sistema aloca componentes de forma incremental; identificada a impossibilidade de satisfação futura das restrições no nó atual da árvore de busca, o algoritmo desfaz a última atribuição e explora caminhos alternativos \cite{tanagra, linden2014dungeons, russell2021artificial};
    \item \textbf{Reparo pós-geração (\textit{Post-generation repair}):} operadores algorítmicos locais atuam sobre o artefato gerado para solucionar violações pontuais sem descartá-lo, tais como transladar pequenos móveis em colisão, regenerar conectores desalinhados ou desobstruir passagens bloqueadas \cite{summerville2018procedural, sun2024roomplanner, xu2025databasedrivenframework3dlevel};
    \item \textbf{Refinamento iterativo em malha fechada (\textit{Feedback loop}):} diagnósticos detalhados extraídos pelos validadores formais são formatados como mensagens estruturadas de erro e reinjetados no módulo generativo de alto nível, permitindo correções informadas em iterações subsequentes \cite{hu2024scenecraft, nasir2024word2worldgeneratingstoriesworlds, xu2025constraint, her2025zeroshot3dmapgeneration}.
\end{enumerate}

% ---------------------------------------------------------------------------
\subsection{Validade, Navegabilidade e Solvabilidade de Fases}
\label{sec:validade_navegabilidade_solvabilidade}
% ---------------------------------------------------------------------------

A avaliação formal da correção de uma fase de jogo exige decompor rigorosamente o conceito genérico de ``jogabilidade funcional'' em quatro dimensões teóricas distintas: validade lógica, validade geométrica, navegabilidade física e solvabilidade global. Essas dimensões operam em planos de abstração distintos e não configuram sinônimos \cite{dormans2010, linden2013designing, linden2014dungeons, sun2024roomplanner, xu2025constraint}.

\subsubsection{Validade Lógica e Estrutural}
\label{sec:validade_logica}

A validade lógica e estrutural diz respeito à consistência relacional e causal do nível em sua representação topológica abstrata, desconsiderando dimensões métricas ou físicas do espaço tridimensional \cite{dormans2010, linden2013designing, linden2014dungeons}. Uma fase é considerada logicamente válida quando satisfaz as invariantes fundamentais do sistema de regras:
\begin{itemize}
    \item Presença de nós de início e fim devidamente parametrizados;
    \item Conectividade do grafo subjacente, atestando que todas as salas e objetivos necessários integram a mesma componente topológica conexa \cite{linden2014dungeons, xu2025databasedrivenframework3dlevel};
    \item Consistência e aciclicidade das dependências causais ($G_{\text{dep}}$ é um DAG), assegurando que não existem bloqueios circulares insolúveis, tais como trancar o acesso à chave de uma porta atrás da própria porta trancada \cite{dormans2010, linden2014dungeons, xu2025databasedrivenframework3dlevel}.
\end{itemize}

\subsubsection{Validade Geométrica}
\label{sec:validade_geometrica}

A validade geométrica refere-se à integridade física do ambiente quando seus módulos são instanciados no espaço contínuo tridimensional $\mathbb{R}^3$ \cite{ericson2004realtime, sun2024roomplanner, xu2025constraint, feng2026repurposing3dgenerativemodel}. Um layout 3D possui validade geométrica se, e somente se:
\begin{itemize}
    \item Nenhuma malha arquitetônica sólida (paredes, pisos ou mobílias densas) sobrepõe-se a outra com volume de interseção positivo ($V_{\text{overlap}} = 0$) \cite{ericson2004realtime, xu2025constraint};
    \item Todos os elementos internos respeitam os limites envoltórios do compartimento que os abriga \cite{sun2024roomplanner, xu2025constraint};
    \item Todas as portas e conectores adjacentes previstos na topologia estão devidamente alinhados e desimpedidos por paredes estruturais \cite{linden2013designing, sun2024roomplanner, xu2025databasedrivenframework3dlevel}.
\end{itemize}

A verificação dessa dimensão em motores gráficos e sistemas computacionais apoia-se em estruturas espaciais de consulta (como AABBs, OBBs, octrees e algoritmos de varredura como \textit{raycasting}) \cite{shaker2016constructive, ericson2004realtime, sun2024roomplanner}. Sob a ótica teórica, essas estruturas operam como ferramentas determinísticas de detecção e consulta espacial, fornecendo subsídios para que algoritmos de verificação confirmem ou rejeitem a conformidade geométrica do layout \cite{shaker2016constructive, xu2025constraint}.

\subsubsection{Navegabilidade e Malhas de Navegação}
\label{sec:navegabilidade}

A navegabilidade (\textit{navigability} ou \textit{traversability}) atesta se o ambiente tridimensional oferece superfícies contínuas e transponíveis sobre as quais um agente cinemático (jogador ou NPC) é fisicamente capaz de se locomover \cite{snook2000simplified, sun2024roomplanner, xu2025constraint, xu2025databasedrivenframework3dlevel}. A navegabilidade é uma propriedade dinâmica e dependente dos atributos físicos do agente, tais como altura livre requerida, raio cilíndrico de colisão, inclinação máxima de superfícies e altura máxima de degraus superáveis \cite{snook2000simplified, xu2025constraint}. Um nível pode exibir perfeita validade geométrica (ausência de colisão entre malhas), mas ser desprovido de navegabilidade caso um corredor seja mais estreito que o diâmetro do agente ou caso ocorra um desnível vertical desprovido de escada ou rampa de acesso \cite{sun2024roomplanner, xu2025databasedrivenframework3dlevel}.

A representação computacional padrão para o espaço contínuo transitável em ambientes tridimensionais consiste na Malha de Navegação (\textit{Navigation Mesh} --- NavMesh), introduzida na literatura de jogos por Snook \cite{snook2000simplified}. A NavMesh discretiza as superfícies caminháveis (\textit{walkable areas}) da geometria 3D decompondo-as em um conjunto de polígonos convexos bidimensionais interconectados no espaço tridimensional \cite{snook2000simplified, regmi2025system}. A convexidade de cada polígono assegura que qualquer trajetória linear traçada entre dois pontos pertencentes ao mesmo polígono permanece estritamente contida dentro da área caminhável, sem colidir com obstáculos \cite{snook2000simplified}. O cálculo de rotas em uma NavMesh opera sobre o \textit{grafo dual de adjacência} formado pelos polígonos convexos (em que os polígonos constituem vértices e as arestas compartilhadas constituem arestas do grafo dual), aplicando-se posteriormente algoritmos de suavização geométrica (como o \textit{funnel algorithm} ou \textit{string pulling}) para gerar caminhos esteticamente suaves e métricamente otimizados \cite{snook2000simplified, xu2025databasedrivenframework3dlevel}.

\subsubsection{Busca de Caminhos (\textit{Pathfinding})}
\label{sec:pathfinding}

A busca de caminhos (\textit{pathfinding}) consiste no procedimento algorítmico responsável por determinar uma rota contínua de travessia entre uma posição inicial $\mathbf{p}_{\text{start}}$ e uma posição alvo $\mathbf{p}_{\text{target}}$ sobre o espaço transitável discretizado (como a NavMesh) \cite{hart1968formal, snook2000simplified, sun2024roomplanner, xu2025databasedrivenframework3dlevel}.

O algoritmo $A^*$ (\textit{A-star}), formulado por Hart, Nilsson e Raphael \cite{hart1968formal}, é a solução canônica para busca informada de caminhos em grafos de navegação. O $A^*$ avalia vértices candidatos por meio de uma função de custo heurístico:
\begin{equation}
    f(n) = g(n) + h(n)
\end{equation}
onde $g(n)$ corresponde ao custo exato do caminho já percorrido da origem até o nó $n$, e $h(n)$ denota a estimativa heurística de custo de $n$ até o objetivo \cite{hart1968formal}. Uma propriedade fundamental demonstrada por Hart et al. \cite{hart1968formal} reside no critério de admissibilidade: se a função heurística $h(n)$ for admissível --- ou seja, se nunca superestimar a distância real mínima até o alvo ($h(n) \le h^*(n)$, como ocorre com a distância euclidiana em linha reta no espaço contínuo) ---, o algoritmo $A^*$ tem garantia matemática de encontrar um caminho de custo ótimo (mínimo), caso exista ao menos uma rota transponível \cite{hart1968formal}.

No contexto da verificação algorítmica de fases, a execução de testes de $A^*$ entre pontos críticos comprova que o espaço modelado mantém conectividade física e operacional real entre compartimentos distintos \cite{sun2024roomplanner, xu2025databasedrivenframework3dlevel}.

\subsubsection{Solvabilidade de uma Fase}
\label{sec:solvabilidade}

A solvabilidade (\textit{solvability}, \textit{winnability} ou completabilidade funcional) constitui o critério avaliativo superior de uma fase: ela estipula a existência de pelo menos uma trajetória válida e realizável pelo jogador no espaço de estados e no espaço geométrico tridimensional, partindo da configuração inicial e atingindo a condição de vitória sem violar nenhuma restrição física ou mecânica do jogo \cite{tanagra, dormans2010, linden2014dungeons, summerville2018procedural, nasir2024word2worldgeneratingstoriesworlds, xu2025constraint}.

A distinção teórica entre validade geométrica, navegabilidade e solvabilidade é indispensável porque nenhuma dessas condições é suficiente para assegurar as demais \cite{dormans2010, linden2014dungeons, summerville2018procedural, xu2025constraint}:
\begin{enumerate}
    \item \textbf{Validade geométrica sem navegabilidade:} dois recintos podem ser posicionados sem nenhuma colisão volumétrica entre suas paredes, mas separados por uma vala intransponível ou por uma rampa com inclinação superior ao limite de locomoção do agente, inviabilizando a transposição física \cite{sun2024roomplanner, xu2025databasedrivenframework3dlevel};
    \item \textbf{Navegabilidade física sem solvabilidade dinâmica:} o algoritmo $A^*$ pode computar com sucesso um caminho contínuo na NavMesh estática entre a entrada e a saída do nível; entretanto, se esse trajeto interceptar uma porta trancada cuja respectiva chave estiver situada em uma câmara cujo acesso físico dependa da prévia transposição dessa mesma porta, o nível apresenta navegabilidade métrica, porém é absolutamente insolúvel devido ao impasse lógico de progressão \cite{dormans2010, linden2013designing, xu2025constraint};
    \item \textbf{Validade lógica abstrata sem solvabilidade espacial:} um grafo de progressão pode ter todas as suas relações de precedência perfeitamente ordenadas no plano abstrato ($\tau(chave) < \tau(porta)$), mas o nível resultante ser insolúvel porque a instanciação geométrica 3D posicionou um obstáculo físico denso que bloqueou o corredor de acesso à sala da chave, rompendo a transitabilidade real da fase \cite{linden2014dungeons, xu2025databasedrivenframework3dlevel}.
\end{enumerate}

Portanto, a solvabilidade plena em ambientes tridimensionais com progressão não linear requer que a alcançabilidade seja verificada simultaneamente sob duas dimensões acopladas: a validade causal no espaço de estados lógicos e a conectividade física contínua comprovada no espaço métrico navegável da fase \cite{dormans2010, linden2014dungeons, xu2025constraint, xu2025databasedrivenframework3dlevel}.

\subsubsection{Síntese Conceitual e Perspectiva Metodológica}
\label{sec:sintese_transicao}

A articulação dos conceitos expostos ao longo deste capítulo evidencia que a síntese de fases tridimensionais complexas impõe um duplo desafio formal: requer, por um lado, alta capacidade semântica e contextual para conceber narrativas, estruturas de desafio e progressões lógicas que regulem o ritmo do jogo; e exige, por outro lado, rigor determinístico absoluto para satisfazer restrições geométricas volumétricas, assegurar alinhamentos modulares e garantir a transitabilidade da malha de navegação.

O reconhecimento das limitações intrínsecas dos modelos conexionistas na manipulação direta de coordenadas métricas, contraposto ao poder estruturante dos formalismos de grafos e da satisfação de restrições, consolida a fundamentação teórica que orienta as abordagens modernas de cocriação e automação em jogos. Esse arcabouço conceitual --- alicerçado na estrita distinção entre organização lógica e realização espacial, e na decomposição analítica entre validade geométrica, navegabilidade e solvabilidade --- estabelece as premissas indispensáveis para a formulação dos modelos, representações intermediárias e rotinas determinísticas de instanciação que serão detalhados no capítulo de Metodologia.

% ---------------------------------------------------------------------------
\section{Trabalhos Relacionados}
\label{sec:trabalhos_relacionados}
% ---------------------------------------------------------------------------

A aplicação de Modelos de Linguagem de Grande Porte (LLMs) à Geração Procedural de Conteúdo (PCG) constitui um campo de investigação recente, emergente e em rápida consolidação na literatura científica de jogos digitais e computação gráfica entre os anos de 2023 e 2026 \cite{Gallotta_2024, maleki2024procedural}. O atrativo central desse paradigma reside na capacidade demonstrada pelos transformadores de interpretar instruções abertas em linguagem natural, compreender relações semânticas complexas e sintetizar estruturas conceituais coerentes a partir de contextos mínimos de entrada \cite{todd2023level, Gallotta_2024}. Contudo, a transposição direta dessas competências conexionistas para a produção de ambientes tridimensionais (3D) navegáveis e funcionais esbarra em limitações estruturais severas, em especial na incapacidade dos modelos autorregressivos de realizar raciocínio métrico, espacial e geométrico contínuo com precisão exata \cite{xu2025constraint, feng2026repurposing3dgenerativemodel}.

Para situar a pesquisa no estado da arte contemporâneo e fundamentar as decisões de projeto adotadas neste trabalho, esta seção organiza as propostas correlatas em quatro famílias arquiteturais fundamentais, acrescidas de uma análise de diagnósticos empíricos e da percepção da indústria de desenvolvimento:
\begin{enumerate}[label=(\roman*)]
    \item \textit{Geração Simbólica e Síntese de Ambientes 2D}: abordagens que exploram a capacidade \textit{in-context} de LLMs sobre representações matriciais discretas (grades ASCII, \textit{tilemaps} e triplas semânticas);
    \item \textit{Síntese de Cenas 3D via Código Intermediário e Renderização Offline}: sistemas que empregam LLMs como geradores de scripts imperativos (como código Python para Blender) combinados a modelos multimodais de visão computacional em ciclos de retroalimentação;
    \item \textit{Planejamento de Layouts 3D Guiado por Otimizadores e Restrições}: propostas que utilizam LLMs para mapear intenções textuais em linguagens de restrições formais ou modelos desacoplados, delegando a geometria a solucionadores numéricos ou heurísticos de reparo;
    \item \textit{Integração com Motores Gráficos e Arquiteturas Multiagente}: pesquisas que acoplam LLMs diretamente a motores de jogos comerciais (\textit{engines}) ou organizam redes de agentes cooperativos em funções de geração e crítica textual.
\end{enumerate}

Cada trabalho correlato é submetido a uma análise dialética tripartite composta por: (1) \textbf{Mecanismo}, descrevendo o arcabouço técnico e o fluxo operacional empregado; (2) \textbf{Mérito}, destacando as contribuições científicas e avanços perante métodos anteriores; e (3) \textbf{Ponto de Ruptura}, dissecando analiticamente o gargalo metodológico ou a falha técnica impeditiva perante o problema de síntese de fases tridimensionais completas, jogáveis e com solvabilidade garantida.

% ---------------------------------------------------------------------------
\subsection{Geração Simbólica e Síntese de Ambientes 2D via Modelos de Linguagem}
\label{sec:relacionados_2d}
% ---------------------------------------------------------------------------

As primeiras investigações sistemáticas sobre o uso de LLMs para a síntese estrutural de fases concentraram-se no domínio bidimensional discreto, no qual a geometria é representada de forma matricial por caracteres ou índices alfanuméricos.

No trabalho pioneiro de Todd et al. \cite{todd2023level}, intitulado \textit{Level Generation Through Large Language Models}, os autores avaliaram a capacidade de modelos de linguagem de grande porte pré-treinados (como GPT-3) gerarem níveis completos para o jogo de quebra-cabeça clássico \textit{Sokoban}. 
\begin{itemize}
    \item \textbf{Mecanismo:} Os autores converteram as fases em representações textuais em formato de grade ASCII, em que cada caractere codifica um elemento específico do jogo (paredes, chão transitável, caixas, alvos e posição do jogador). Por meio do aprendizado em contexto (\textit{in-context learning}), pares de exemplos de entrada e saída foram fornecidos no \textit{prompt} para instruir o modelo a completar e sintetizar novas matrizes caractere a caractere de maneira autorregressiva direta.
    \item \textbf{Mérito:} O estudo comprovou experimentalmente, de maneira pioneira, que os LLMs detêm a capacidade latente de assimilar e reproduzir a gramática superficial e as regularidades estilísticas locais de níveis de jogos sem a necessidade de ajuste fino supervisionado (\textit{fine-tuning}) em bases proprietárias.
    \item \textbf{Ponto de Ruptura:} A geração puramente autorregressiva opera sob a ausência de invariantes geométricos e espaciais globais. À medida que o tamanho da grade aumenta, a deriva estocástica token a token acumula inconsistências métricas, resultando em barreiras incompletas, alvos inacessíveis e uma taxa expressiva de fases insolúveis. Além disso, a abordagem é estritamente confinada a espaços 2D discretos regulares, sendo estruturalmente incapaz de extrapolar para coordenadas métricas contínuas em $\mathbb{R}^3$, volumes de colisão ou malhas poligonais de navegação.
\end{itemize}

Ampliando o escopo para a geração procedimental integrada entre narrativa e espaço, Nasir, James e Togelius \cite{nasir2024word2worldgeneratingstoriesworlds} apresentaram o sistema \textit{Word2World: Generating Stories and Worlds through Large Language Models}.
\begin{itemize}
    \item \textbf{Mecanismo:} O sistema emprega um pipeline modular no qual LLMs sintetizam recursivamente narrativas descritivas e convertem esses enredos em mundos bidimensionais subdivididos em blocos (\textit{tile-based}), interligados por pontos de interesse narrativo. Para atestar a viabilidade básica das áreas geradas, o sistema acopla um algoritmo de busca heurística $A^*$ tradicional que verifica a existência de caminhos livres sobre a matriz discreta de tiles ortogonais.
    \item \textbf{Mérito:} A pesquisa estabelece um avanço relevante na coesão entre enredo e mapa de jogo, demonstrando a utilidade de submeter a saída textual do LLM a um validador clássico de conectividade física.
    \item \textbf{Ponto de Ruptura:} O escopo espacial permanece estritamente restrito a \textit{tilemaps} 2D estáticos e ortogonais. Não há qualquer tratamento de transformações afins tridimensionais (translação, rotação de guinada e escala volumétrica), relevo vertical ou detecção de colisão contínua. Ademais, a validação por $A^*$ restringe-se à conectividade geométrica estática de caminhos livres, ignorando a dinâmica temporal de dependências mecânicas (como sistemas de chaves e portas trancadas, onde o estado do agente altera dinamicamente a transitabilidade de barreiras no espaço).
\end{itemize}

Em outra vertente 2D, Chen e Jhala \cite{chen2025narrative} propuseram o arcabouço \textit{Narrative-to-Scene Generation: An LLM-Driven Pipeline for 2D Game Environments}.
\begin{itemize}
    \item \textbf{Mecanismo:} O pipeline traduz histórias e descrições textuais livres em um conjunto intermediário de triplas semânticas relacionais (da forma $\langle$\textit{entidade}, \textit{relação}, \textit{entidade}$\rangle$, como ``árvore à esquerda da cabana''). Em seguida, essas relações espaciais qualitativas são traduzidas em restrições de um Problema de Satisfação de Restrições (CSP) bidimensional, resolvido algoritmicamente para dispor os elementos visuais sobre a cena.
    \item \textbf{Mérito:} O trabalho mitiga a fragilidade numérica do LLM ao abstrair coordenadas absolutas em relações qualitativas estruturadas, delegando a acomodação geométrica a um solucionador simbólico formal.
    \item \textbf{Ponto de Ruptura:} O método foca unicamente em cenas 2D decorativas e estáticas, assemelhando-se a dioramas visuais planos. Ele desconsidera por completo as exigências do \textit{Level Design} interativo de fases jogáveis: não modela a locomoção física de agentes cinemáticos, não calcula Malhas de Navegação (NavMesh), não suporta ambientes modulares conectados por corredores nem trata a progressão funcional e a solvabilidade lógica de quebra-cabeças mecânicos.
\end{itemize}

% ---------------------------------------------------------------------------
\subsection{Síntese de Cenas 3D via Código Intermediário e Renderização \textit{Offline}}
\label{sec:relacionados_codigo_render}
% ---------------------------------------------------------------------------

Diante da constatação empírica de que os LLMs falham ao emitir diretamente listas de coordenadas métricas $(x, y, z)$, uma linha de pesquisa buscou explorar a síntese de código como linguagem intermediária, combinando geradores neurais a suítes de modelagem e renderização tridimensional \textit{offline}.

O principal expoente desse paradigma é o sistema \textit{SceneCraft}, desenvolvido por Hu et al. \cite{hu2024scenecraft} e publicado na ICML 2024.
\begin{itemize}
    \item \textbf{Mecanismo:} O \textit{SceneCraft} estrutura um agente autônomo baseado em LLM encarregado de traduzir descrições em linguagem natural em código executável em linguagem Python, o qual invoca comandos da API interna do software de computação gráfica Blender. O código programa a instanciação de primitivas e malhas 3D, ajustando transformações no espaço contínuo. Para mitigar colisões visuais e erros de proporção, o sistema fecha um laço de retroalimentação multimodal (\textit{feedback loop}): uma câmera virtual renderiza a cena em imagens bidimensionais a partir de múltiplos ângulos, as quais são submetidas a um Modelo de Visão e Linguagem (VLM); o VLM inspeciona os artefatos visuais e devolve críticas em texto para que o LLM reescreva e refine o código Python iterativamente.
    \item \textbf{Mérito:} O emprego de código como camada intermediária formaliza as operações de posicionamento, contornando a emissão de números puros sem contexto sintático. A incorporação do feedback visual via VLM introduz uma verificação empírica automatizada de sobreposições visíveis na renderização.
    \item \textbf{Ponto de Ruptura:} O sistema opera sob um fluxo estritamente \textit{offline} e computacionalmente oneroso, consumindo minutos de processamento por cena devido à renderização fotorrealista e às chamadas de modelos multimodais de alta escala, o que inviabiliza sua aplicação em tempo de execução ou como ferramenta interativa em motores de jogos. Fundamentalmente, o artefato resultante do \textit{SceneCraft} é uma cena gráfica estática e contemplativa: não há Malha de Navegação (NavMesh), não há conectores ou portas funcionais, inexiste física de colisão de tempo real com agentes jogáveis e não há suporte à progressão de objetivos ou validação de solvabilidade de fases.
\end{itemize}

Mais recentemente, Yang et al. \cite{yang2026codeasroomgenerating3drooms} apresentaram o estudo exploratório \textit{Code-as-Room: Generating 3D Rooms from Top-Down View Images via Agentic Code Synthesis}.
\begin{itemize}
    \item \textbf{Mecanismo:} O arcabouço recebe como entrada imagens de plantas baixas residenciais (\textit{top-down view}) e coordena agentes de síntese de código especializados em decompor o perímetro do recinto e alocar mobílias e adereços no espaço tridimensional interno por meio de scripts procedurais.
    \item \textbf{Mérito:} Apresenta alta eficácia na acomodação espacial de móveis respeitando os limites perimetrais de paredes em ambientes residenciais fechados, aproveitando a capacidade de abstração modular do código para evitar sobreposições grosseiras.
    \item \textbf{Ponto de Ruptura:} O escopo do \textit{Code-as-Room} é categoricamente confinado a cômodos residenciais individuais e isolados (\textit{single room layout}). O sistema não possui mecanismos para sintetizar masmorras ou fases completas compostas por múltiplas câmaras articuladas por corredores, bifurcações e transições verticais. Além disso, ignora completamente as variáveis de jogabilidade, fluxo de circulação de jogadores, posicionamento de desafios e encadeamento causal de mecânicas de jogo.
\end{itemize}

% ---------------------------------------------------------------------------
\subsection{Planejamento de \textit{Layouts} 3D Guiado por Otimizadores e Restrições}
\label{sec:relacionados_otimizadores}
% ---------------------------------------------------------------------------

Reconhecendo as falhas geométricas crônicas dos métodos puramente neurais, uma terceira família de pesquisas procurou desacoplar a geração em dois estratos: um planejamento de alto nível orientado por LLM e uma etapa de alocação espacial governada por otimizadores matemáticos ou acervos de dados estruturados.

No trabalho de Sun et al. \cite{sun2024roomplanner}, intitulado \textit{RoomPlanner: Explicit Layout Planner for Easier LLM-Driven 3D Room Generation}, os autores abordaram o arranjo interno de recintos tridimensionais.
\begin{itemize}
    \item \textbf{Mecanismo:} O sistema propõe a divisão explícita em dois estágios sequenciais: primeiramente, o LLM atua como um planejador macroestrutural, gerando a disposição relacional dos objetos sem especificar transformações milimétricas; em seguida, um módulo determinístico de recuperação e posicionamento de ativos (\textit{asset retrieval and placement}) realiza o encaixe físico das malhas dentro do ambiente delimitado.
    \item \textbf{Mérito:} Formaliza a necessidade de segregar a concepção estrutural abstrata da resolução geométrica fina, reconhecendo explicitamente que a imposição de cálculos métricos simultâneos sobrecarrega o LLM e eleva as taxas de alucinação espacial.
    \item \textbf{Ponto de Ruptura:} O modelo é estritamente restrito ao arranjo doméstico de mobília (\textit{indoor furniture placement}). Não contempla as exigências de uma fase de jogo: não existem corredores de transição, conectores modulares, zonas de combate, quebra-cabeças com dependências temporais nem requisitos cinemáticos de navegação transitável para agentes autônomos.
\end{itemize}

Em uma abordagem voltada especificamente para fases de jogos, Xu e Verbrugge \cite{xu2025constraint} apresentaram na FDG 2025 o estudo \textit{Constraint Is All You Need: Optimization-Based 3D Level Generation with LLMs}.
\begin{itemize}
    \item \textbf{Mecanismo:} O usuário fornece uma descrição em linguagem natural do nível desejado. O LLM é instruído a traduzir essa intenção em um conjunto formal de restrições espaciais e topológicas codificadas em uma linguagem de domínio proprietária, denominada GLDL (\textit{Game Level Description Language}). Essas restrições são então mapeadas em uma formulação contínua do Problema de Layout de Instalações (FLP --- \textit{Facility Layout Problem}) e resolvidas por otimizadores numéricos de programação não linear com restrições (tais como IPOPT e SLSQP), que calculam posições contínuas e dimensões de blocos retangulares assegurando a ausência de sobreposição e a adjacência especificada.
    \item \textbf{Mérito:} Estabelece um dos patamares mais rigorosos da literatura recente ao substituir o palpite probabilístico de coordenadas por formulações matemáticas clássicas de otimização, garantindo analiticamente a não sobreposição volumétrica ($V_{\text{overlap}} = 0$) e o cumprimento de adjacências mínimas entre volumes abstratos.
    \item \textbf{Ponto de Ruptura:} A técnica opera como um solucionador numérico puramente matemático e estritamente desacoplado de motores de jogos comerciais (\textit{game engines}). Como consequência desse isolamento: (i) as restrições são calculadas sobre caixas delimitadoras ideais, sem levar em consideração as malhas poligonais detalhadas dos \textit{prefabs} tridimensionais reais; (ii) o sistema não executa o cozimento (\textit{bake}) nem a validação de Malhas de Navegação (NavMesh), ignorando parâmetros cinemáticos de agentes reais (como altura de degraus e raio de passagem); e (iii) a formulação do FLP é puramente geométrica e estática, desconsiderando por completo o espaço de estados da progressão lógica do jogo (isto é, não modela nem verifica se o posicionamento de chaves e portas trancadas gera impasses ou bloqueios intransponíveis na dinâmica da partida).
\end{itemize}

Complementarmente, Xu e Verbrugge \cite{xu2025databasedrivenframework3dlevel} publicaram na AIIDE 2025 o artigo \textit{A Database-Driven Framework for 3D Level Generation with LLMs}.
\begin{itemize}
    \item \textbf{Mecanismo:} Os autores estruturaram um repositório relacional de salas e módulos pré-modelados em 3D, contendo anotações semânticas e pontos de conexão validados. O LLM atua na seleção e composição dessas salas a partir de consultas textuais. Para contornar desconexões geradas pelo modelo durante a montagem do nível, o framework incorpora um algoritmo de pós-processamento baseado em heurísticas de reparo topológico, encarregado de inserir corredores adicionais ou rotacionar salas desalinhadas.
    \item \textbf{Mérito:} O reúso de blocos arquiteturais pré-construídos eleva a fidelidade visual local e garante que as unidades atômicas da masmorra sejam internamente válidas.
    \item \textbf{Ponto de Ruptura:} O sistema padece de alta rigidez, dependendo intrinsecamente de um vocabulário fechado de ativos pré-cadastrados no banco de dados. A necessidade de etapas corretivas \textit{ad hoc} de reparo evidencia a incapacidade do gerador neural em preservar a integridade estrutural da fase de forma autônoma. Ademais, o framework não é integrado nativamente ao fluxo de trabalho de um motor gráfico comercial (como ferramenta ou \textit{plugin} de editor), limitando severamente a interatividade e a ergonomia de autoria assistida para os desenvolvedores.
\end{itemize}

% ---------------------------------------------------------------------------
\subsection{Integração com Motores de Jogos e Arquiteturas Multiagente}
\label{sec:relacionados_motores_multiagente}
% ---------------------------------------------------------------------------

Uma quarta frente de trabalhos tem se concentrado na ponte direta entre LLMs e motores de jogos comerciais (\textit{game engines}), bem como na colaboração entre múltiplos agentes neurais para mitigar erros de geração.

No trabalho de Her et al. \cite{her2025zeroshot3dmapgeneration}, intitulado \textit{Zero-shot 3D Map Generation with LLM Agents: A Dual-Agent Architecture for Procedural Content Generation}, os autores propuseram uma arquitetura multiagente cooperativa.
\begin{itemize}
    \item \textbf{Mecanismo:} O sistema divide o processo de geração entre dois agentes fundamentados em LLMs operando em uma configuração Ator-Crítico (\textit{Actor-Critic}): um agente gerador formula a especificação espacial do mapa a partir do prompt do usuário, e um segundo agente crítico analisa textualmente o rascunho produzido, apontando inconsistências conceituais e devolvendo orientações de refinamento para o gerador em um laço de diálogo \textit{zero-shot}.
    \item \textbf{Mérito:} O trabalho explora a auditoria cruzada sem exigir retreinamento de pesos ou ajustes finos complexos, demonstrando que o diálogo entre agentes pode detectar certas incoerências semânticas e contextuais.
    \item \textbf{Ponto de Ruptura:} A arquitetura incorre no que se denomina na literatura como a \textbf{falácia da auto-validação probabilística}. Ambos os agentes compartilham a mesma natureza conexionista autorregressiva, carecendo intrinsecamente de um modelo físico e geométrico do espaço contínuo. O agente crítico é incapaz de calcular com exatidão se duas caixas delimitadoras tridimensionais colidem no espaço contínuo $\mathbb{R}^3$, não tem acesso a rotinas de colisão física e não pode atestar a conectividade transitável de uma Malha de Navegação poligonal. O ciclo crítico-gerador opera, portanto, em um plano puramente retórico, permanecendo sujeito a alucinações métricas imperceptíveis a ambos os modelos.
\end{itemize}

No âmbito de integração nativa em engines comerciais de alta fidelidade, destaca-se o trabalho de Tang et al. \cite{tang2025unrealllm}, intitulado \textit{UnrealLLM: Towards Highly Controllable and Interactable 3D Scene Generation by LLM-powered Procedural Content Generation}, publicado na ACL 2025.
\begin{itemize}
    \item \textbf{Mecanismo:} O \textit{UnrealLLM} conecta um modelo de linguagem à \textit{Unreal Engine 5} (UE5) por meio de um protocolo baseado em mensagens estruturadas em JSON e ferramentas procedurais baseadas em curvas paramétricas (\textit{splines}). O LLM decompõe instruções em linguagem natural em comandos estruturados para instanciar terrenos, relevos, corpos hídricos, densidade de vegetação e trajetórias de vias, permitindo a autoria de cenas interativas controladas pelo motor gráfico.
    \item \textbf{Mérito:} Representa um dos maiores avanços práticos de acoplamento entre modelos de linguagem e um motor de jogos comercial consolidado na indústria, comprovando que a mediação via esquemas estruturados em JSON viabiliza o controle programático de ambientes gráficos de ponta.
    \item \textbf{Ponto de Ruptura:} O escopo do \textit{UnrealLLM} é quase que exclusivamente direcionado a ambientes abertos (\textit{outdoor}) e à dispersão cosmética e realista de ativos ao longo de superfícies contínuas e splines. O sistema não aborda a geração modular de fases fechadas ou masmorras compostas por salas, corredores e conectores discretos. Fundamentalmente, não há tratamento da progressão mecânica de \textit{Level Design}: o sistema não suporta a formalização de grafos causais acíclicos (DAGs) com pré-requisitos mecânicos (chaves, portas trancadas e interruptores), carecendo de rotinas determinísticas para comprovar a solvabilidade dinâmica da fase no espaço de estados e a integridade de rotas navegáveis em ambientes interiores fechados.
\end{itemize}

Em paralelo, pesquisas complementares como as de Regmi \cite{regmi2025system} e Re \cite{re2026integration} investigaram o emprego de LLMs acoplados a motores de jogos (especialmente Unity). Regmi \cite{regmi2025system} desenvolveu um sistema para geração procedural de atividades de agentes virtuais baseado na interpretação semântica de cenas tridimensionais, ao passo que Re \cite{re2026integration} explorou técnicas de engenharia de \textit{prompt} voltadas à síntese automática de código para agentes inteligentes no desenvolvimento de jogos. Embora ambos os trabalhos atestem a crescente permeabilidade dos LLMs nas rotinas de motores de jogos, seus focos concentram-se estritamente na orquestração de comportamentos de NPCs e na automação de scripts utilitários, sem abordar a síntese procedural da topologia e da geometria de fases tridimensionais, nem os problemas formais de colisão física e solvabilidade mecânica.

% ---------------------------------------------------------------------------
\subsection{Diagnósticos Empíricos de Falhas Espaciais e Percepção do \textit{Level Designer}}
\label{sec:relacionados_diagnosticos}
% ---------------------------------------------------------------------------

Para além das arquiteturas generativas, a consolidação recente da literatura trouxe investigações empíricas e relatos da indústria que evidenciam com rigor metodológico as fragilidades dos geradores conexionistas diretos e a recepção dos profissionais de desenvolvimento.

O estudo empírico de maior impacto nesse domínio foi conduzido por Feng et al. \cite{feng2026repurposing3dgenerativemodel}, sob o título \textit{Repurposing 3D Generative Model for Autoregressive Layout Generation}.
\begin{itemize}
    \item \textbf{Diagnóstico Técnico:} Os autores realizaram testes quantitativos exaustivos sobre a capacidade de modelos generativos neurais e transformadores em emitir diretamente arranjos espaciais tridimensionais. A conclusão central da pesquisa foi formalizada na máxima: \textbf{``Layout as Language: Semantic YES, Physical NO''} (Layout como Linguagem: Sim para Semântica, Não para Física).
    \item \textbf{Evidências Empíricas:} Os experimentos comprovaram que, conquanto os modelos alcancem escores elevados de plausibilidade semântica (como posicionar cadeiras próximas a mesas e associar objetos contextuais corretamente em texto), a geração numérica direta de posições e rotações 3D incorre em falhas geométricas frequentes e inaceitáveis: taxas severas de sobreposição volumétrica não autorizada (\textit{clipping} massivo), elementos suspensos flutuando sem apoio estrutural (\textit{floating objects}), violações perimétricas de colisão contra paredes e desvios de escala que quebram as restrições físicas mínimas do ambiente euclidiano.
    \item \textbf{Implicação para a Área:} Esse diagnóstico confirma teoricamente que a tentativa de delegar o cálculo exato de coordenadas geométricas contínuas a modelos estatísticos autorregressivos constitui uma premissa arquitetural falha, demandando imperativamente que as variáveis métricas sejam governadas por processos determinísticos externos com respaldo nas leis físicas e espaciais.
\end{itemize}

Sob a perspectiva humana e profissional do desenvolvimento de jogos, Bandeali \cite{bandeali2026generative} publicou o relatório investigativo \textit{Generative AI \& The Level Designer: An Exploratory Report of Tools, Workflows, and Challenges}.
\begin{itemize}
    \item \textbf{Diagnóstico da Indústria:} O autor conduziu um estudo de campo detalhado com \textit{level designers} profissionais em estúdios de desenvolvimento de jogos sobre a integração de ferramentas de inteligência artificial generativa em suas rotinas de trabalho.
    \item \textbf{Desafios e Rejeição de Ferramentas Autônomas:} O estudo revelou um ceticismo generalizado da indústria em relação a geradores autônomos do tipo ``caixa-preta'' que entregam cenas fechadas sem controle intermediário. Os profissionais relataram que o tempo despendido para auditar, desmanchar e corrigir as alucinações geométricas, caminhos bloqueados e quebras de ritmo gerados por ferramentas puramente neurais frequentemente supera o tempo necessário para criar a fase manualmente do zero.
    \item \textbf{Requisitos Mandatórios da Indústria:} O relatório conclui que as ferramentas generativas assistidas só alcançam viabilidade e adoção real na indústria quando atendem a três pilares: (i) operação no paradigma de \textbf{iniciativa mista} (\textit{mixed-initiative}), mantendo o designer no controle criativo através de parâmetros claros de edição; (ii) integração nativa e ergonômica ao ambiente do motor de jogo (\textit{game engine}) utilizado pelo estúdio; e (iii) existência de salvaguardas algorítmicas determinísticas que garantam categoricamente que o artefato gerado seja funcional, desprovido de sobreposições volumétricas destrutivas e com navegabilidade atestada.
\end{itemize}

Essa demanda da indústria encontra eco direto em estudos seminais e contemporâneos de iniciativa mista e autoria assistida, tais como os trabalhos de Lai, Latham e Leymarie \cite{lai2020towards} sobre os três pilares de ferramentas amigáveis para a indústria, de Walton, Rahat e Stovold \cite{walton2020evaluating} sobre avaliação cega de ferramentas em estúdios de jogos, e de Wang et al. \cite{wang2026designerlyloop} sobre o alinhamento da intenção do designer com LLMs por meio de circuitos iterativos de raciocínio.

% ---------------------------------------------------------------------------
\subsection{Quadro Comparativo do Estado da Arte}
\label{sec:relacionados_quadro_comparativo}
% ---------------------------------------------------------------------------

Para consolidar as distinções metodológicas e explicitar os limites de cada linha de investigação analisada, a Tabela~\ref{tab:quadro_comparativo_relacionados} (Quadro comparativo do estado da arte) estabelece uma comparação sistemática e multidimensional entre os 11 principais sistemas da literatura e a proposta delineada no presente TCC. A matriz avaliativa cruza os trabalhos frente a sete dimensões técnicas rigorosas:
\begin{enumerate}
    \item \textbf{Paradigma Arquitetural:} a taxonomia do pipeline (puramente neural, multiagente, assistido por otimizador numérico ou híbrido neuro-simbólico);
    \item \textbf{Espaço de Representação:} dimensionalidade do artefato produzido (2D discreto vs. 3D contínuo);
    \item \textbf{Representação Intermediária:} a estrutura de dados empregada entre a entrada e a saída (grade ASCII, triplas relacionais, código imperativo, formulação de restrições ou grafo serializado);
    \item \textbf{Garantia Anti-Clipping:} o mecanismo formal que assegura a não sobreposição física de malhas sólidas no espaço contínuo ($V_{\text{overlap}} = 0$);
    \item \textbf{Navegabilidade Dinâmica:} capacidade de gerar e validar malhas de navegação contínuas (NavMesh) com suporte a agentes cinemáticos sob física da engine;
    \item \textbf{Solvabilidade Mecânica:} capacidade de modelar e comprovar a alcançabilidade dinâmica de quebra-cabeças lógicos (como sistemas de chaves e fechaduras no espaço de estados);
    \item \textbf{Integração com Game Engine:} se o sistema opera incorporado como ferramenta nativa ao editor de um motor de jogos comercial.
\end{enumerate}

\begin{table}[H]
\centering
\caption{Quadro comparativo multidimensional do estado da arte em geração de fases e cenas com LLMs.}
\label{tab:quadro_comparativo_relacionados}
\resizebox{\linewidth}{!}{%
\begin{tabular}{@{}llllllll@{}}
\toprule
\textbf{Trabalho / Sistema} & \textbf{Paradigma} & \textbf{Espaço} & \textbf{Intermediário} & \textbf{Anti-Clipping} & \textbf{Navegabilidade} & \textbf{Solvabilidade} & \textbf{Game Engine} \\ \midrule
Todd et al. \cite{todd2023level} & LLM Puro (ICL) & 2D & Grade ASCII & Não & Não & Baixa (Probabilística) & Não (Simulador) \\
Word2World \cite{nasir2024word2worldgeneratingstoriesworlds} & LLM + Busca & 2D & Tilemap Discreto & Não (Grid) & $A^*$ em Grid 2D & Não & Não (Python) \\
Narrative-to-Scene \cite{chen2025narrative} & LLM + CSP & 2D & Triplas Semânticas & Heurística (2D) & Não & Não se aplica & Não \\ \midrule
SceneCraft \cite{hu2024scenecraft} & LLM + VLM & 3D & Código Python & Feedback Visual & Não & Não se aplica & Não (Blender Offline) \\
Code-as-Room \cite{yang2026codeasroomgenerating3drooms} & Multiagente Código & 3D & Código Python & Heurística Código & Não & Não se aplica & Não (Standalone) \\ \midrule
RoomPlanner \cite{sun2024roomplanner} & LLM + Slot Placement & 3D & Bounding Boxes & Heurística AABB & Não & Não se aplica & Não \\
Xu e Verbrugge (GLDL) \cite{xu2025constraint} & LLM + Otimizador FLP & 3D & GLDL (Restrições) & Numérica Exata & Não (Sem NavMesh) & Não (Estático) & Não (Solver Python) \\
Xu e Verbrugge (DB) \cite{xu2025databasedrivenframework3dlevel} & LLM + Base de Dados & 3D & Grafo Semântico & Heurística Reparo & Parcial (Heurística) & Não & Não \\ \midrule
Her et al. \cite{her2025zeroshot3dmapgeneration} & Multiagente LLM & 3D & Texto Estruturado & Não (Probabilística) & Não & Não & Não \\
UnrealLLM \cite{tang2025unrealllm} & LLM + Procedural UE5 & 3D & JSON / Splines & Física da Engine & Sim (Terreno Aberto) & Não (Sem Chaves) & Sim (Unreal Engine 5) \\
Regmi / Re \cite{regmi2025system, re2026integration} & LLM + Scripting & 3D & Código / Ações & Física da Engine & Parcial (NPCs) & Não se aplica & Sim (Unity) \\ \midrule
\textbf{Proposta deste TCC} & \textbf{Neuro-Simbólico} & \textbf{3D} & \textbf{Grafo JSON (DAG)} & \textbf{Determinística} & \textbf{NavMesh Bake ($A^*$)} & \textbf{Formal (Espaço de Estados)} & \textbf{Sim (Godot Engine)} \\ \bottomrule
\end{tabular}%
}
\vspace{0.1cm}
\begin{flushleft}
\footnotesize{\textbf{Fonte:} Elaborado pelo autor (2026).}
\end{flushleft}
\end{table}

A inspeção detalhada da Tabela~\ref{tab:quadro_comparativo_relacionados} permite deduzir regularidades estruturais claras sobre a literatura contemporânea:
\begin{itemize}
    \item As abordagens bidimensionais \cite{todd2023level, nasir2024word2worldgeneratingstoriesworlds, chen2025narrative} operam satisfatoriamente em representações discretas e permitem a execução de buscas heurísticas simples em grade, mas são fundamentalmente incapazes de transpor suas regras para a complexidade contínua do espaço métrico tridimensional $\mathbb{R}^3$, ignorando volumes de colisão, conectores angulares e relevo vertical;
    \item Os sistemas baseados em síntese de código para suítes de modelagem \textit{offline} \cite{hu2024scenecraft, yang2026codeasroomgenerating3drooms} atingem qualidade estética relevante, mas operam em tempos de processamento proibitivos e geram modelos estáticos e contemplativos, desprovidos de qualquer atributo de jogabilidade, malhas de navegação ou lógica de masmorras;
    \item As abordagens fundamentadas em solucionadores numéricos \cite{xu2025constraint, xu2025databasedrivenframework3dlevel} oferecem sólidas garantias de não colisão volumétrica em modelos ideais, porém permanecem confinadas a processos analíticos desacoplados de motores gráficos comerciais, falhando em gerar Malhas de Navegação e em validar a solvabilidade mecânica dinâmica em tempo de jogo;
    \item As soluções integradas a motores de jogos \cite{tang2025unrealllm, regmi2025system} aproveitam o potencial de engines consagradas, porém concentram-se na dispersão superficial de folhagens e terrenos em ambientes abertos ou no controle comportamental de NPCs, deixando desassistida a síntese modular de masmorras com progressão condicional chave-porta.
\end{itemize}

% ---------------------------------------------------------------------------
\subsection{Síntese da Lacuna na Literatura e Posicionamento da Proposta}
\label{sec:relacionados_lacuna_posicionamento}
% ---------------------------------------------------------------------------

A análise comparativa e o mapeamento dialético dos trabalhos correlatos evidenciam a existência de uma \textbf{tríplice lacuna científica} na literatura contemporânea de Geração Procedural de Fases assistida por Inteligência Artificial:

\begin{enumerate}
    \item \textbf{A falácia da auto-validação probabilística e a ausência de ancoragem física determinística:} Tentativas de utilizar agentes LLMs para auditar outros LLMs em tarefas espaciais \cite{her2025zeroshot3dmapgeneration} incorrem em uma circularidade conexionista estocástica. Modelos autorregressivos não possuem modelo de mundo físico incorporado nem aparato analítico para solucionar geometria contínua euclidiana. Sem um validador simbólico e determinístico operando sobre estruturas de dados espaciais precisas (como árvores de delimitadores e malhas poligonais), é teoricamente inviável assegurar a ausência de sobreposição física e a coerência métrica de ambientes tridimensionais, conforme comprovado empiricamente por Feng et al. \cite{feng2026repurposing3dgenerativemodel}.
    
    \item \textbf{A dissociação entre estética cosmética e arquitetura funcional de \textit{Level Design}:} O estado da arte encontra-se polarizado entre geradores puramente visuais, focados em renderizações fotorrealistas ou arranjos de mobília estática \cite{hu2024scenecraft, sun2024roomplanner}, e modelos abstratos de grafos que ignoram a transitabilidade real do espaço. Nenhuma proposta existente articula conjuntamente a Malha de Navegação (NavMesh com $A^*$) às dependências dinâmicas temporais de missões (como o problema de chaves e portas trancadas governado por um Grafo Acíclico Dirigido no espaço de estados). Ambientes esteticamente impressionantes tornam-se inúteis para a experiência interativa caso não ofereçam garantia formal de completabilidade funcional e navegabilidade.
    
    \item \textbf{O distanciamento entre solucionadores acadêmicos abstratos e motores de jogos comerciais:} Soluções metodologicamente sofisticadas baseadas em restrições \cite{xu2025constraint} operam em isolamento computacional, desconectadas dos pipelines e ferramentas onde os jogos digitais são efetivamente desenvolvidos e produzidos. A literatura carece de arquiteturas neuro-simbólicas integradas nativamente como ferramentas de autoria (\textit{plugins} ou \textit{addons}) em motores de jogos contemporâneos (como o Godot Engine), capazes de oferecer assistência criativa e validação imediata em tempo de desenvolvimento, atendendo categoricamente às demandas de ergonomia e controle expressas pelos profissionais da indústria \cite{bandeali2026generative}.
\end{enumerate}

Para superar de maneira unificada essa tríplice lacuna, o presente Trabalho de Conclusão de Curso propõe uma \textbf{arquitetura híbrida e neuro-simbólica} para a geração assistida de fases tridimensionais modulares:
\begin{itemize}
    \item O estrato cognitivo e semântico de alto nível é delegado ao Modelo de Linguagem de Grande Porte, que atua estritamente como um \textbf{planejador lógico}, traduzindo o prompt em linguagem natural em um grafo causal de progressão estruturado e serializado em formato JSON (contendo salas funcionais, corredores, conexões e pré-requisitos mecânicos de chaves e portas);
    \item O estrato geométrico, físico e mecânico é delegado integralmente a rotinas algorítmicas determinísticas implementadas de forma nativa no motor de jogos \textbf{Godot Engine 4}, que realizam a instanciação espacial contínua de malhas 3D modulares com prevenção rigorosa de colisões (AABB), o cálculo e cozimento programático da Malha de Navegação poligonal (NavMesh), o teste de transitabilidade com $A^*$ e a validação formal da solvabilidade dinâmica no espaço de estados dinâmicos.
\end{itemize}

A formalização dos modelos matemáticos, dos esquemas de dados, dos algoritmos de resolução geométrica e do pipeline de validação implementado na \textit{engine} constitui o escopo detalhado no capítulo de Metodologia que se segue.


% ===========================================================================
%  REFERENCIAS
% ===========================================================================
\SingleSpacing
\printbibliography[title={Bibliografia}]
\end{document}